% Regression and Correlation — Pure Mathematics with Statistics Upper Sixth — Cours signé OnBuch+
% Official MINESEC syllabus. Compiler avec : tectonic PMS09-regression-and-correlation.tex
\newcommand{\DOCMATIERE}{Pure Mathematics with Statistics}
\newcommand{\DOCNIVEAU}{Upper Sixth}
\newcommand{\DOCMODULE}{Upper Sixth — Pure mathematics with statistics}
\newcommand{\DOCLECON}{9}
\newcommand{\DOCTITRE}{Regression and Correlation}
% OnBuch+ — préambule « cours signés » ENRICHI (compilé avec tectonic / XeLaTeX)
% Système visuel « Pop » d'OnBuch+ : fond crème (jamais blanc), bordures encre
% épaisses, ombres « dures » décalées, titres Archivo Black en capitales.
% Version MATHÉMATIQUES : graphes (pgfplots/tikz), boîtes pédagogiques
% (définition, propriété, méthode, attention, le savais-tu, exercice résolu,
% exercice, TP, bilan), page de garde et signature OnBuch+ ancrée en bas.
%
% Macros attendues dans main.tex AVANT \input{preamble} :
%   \DOCMATIERE  \DOCNIVEAU  \DOCMODULE  \DOCLECON (numéro)  \DOCTITRE
\documentclass[11pt]{article}

\usepackage{fontspec}
\usepackage[english]{babel}
\usepackage[a4paper,margin=2cm,top=3.4cm,bottom=2.8cm,headheight=1.4cm,footskip=1.3cm]{geometry}
\usepackage{amsmath,amssymb,amsfonts}
\usepackage{mathtools} % \xmapsto, \coloneqq... souvent utilisés par les modèles
\usepackage{cancel} % simplifications barrées (bilans de forces)
% \abs{z} (module, valeur absolue) et \norm{u} (norme), utilisés par les modèles
\providecommand{\abs}{}\let\abs\relax\DeclarePairedDelimiter{\abs}{\lvert}{\rvert}
\providecommand{\norm}{}\let\norm\relax\DeclarePairedDelimiter{\norm}{\lVert}{\rVert}
\usepackage[table]{xcolor} % [table] : fournit aussi \rowcolors (alternance de lignes)
\usepackage{pagecolor}
\usepackage{tikz}
\usetikzlibrary{arrows.meta,positioning,calc,shapes.geometric,shapes.misc,shapes.arrows,decorations.pathmorphing,decorations.pathreplacing,decorations.markings,patterns,shadows,mindmap,trees,fit,backgrounds,matrix,intersections,angles,quotes}
\usepackage{pgfplots}
\usepackage[version=4]{mhchem} % équations nucléaires \ce{^{235}_{92}U}
\usepackage[european,siunitx]{circuitikz} % schémas électriques
\pgfplotsset{compat=1.18}
\usepgfplotslibrary{fillbetween,groupplots} % aires entre courbes (calcul intégral)
\usepackage{enumitem}
\usepackage{fancyhdr}
% [most] charge toutes les bibliothèques tcolorbox (listings, minted, raster,
% documentation...) et gonfle inutilement la mémoire TeX sur un document long
% (~300 boîtes) : on ne charge que ce qui est réellement utilisé.
\usepackage[skins,breakable]{tcolorbox}
\usepackage{graphicx}
\usepackage{colortbl}
\usepackage{booktabs}
\usepackage{tabularx}
\usepackage{diagbox} % case d'en-tête diagonale des tableaux à double entrée
% Colonnes extensibles centrée / gauche / droite (C, L, R), souvent utilisées par les modèles
\newcolumntype{C}{>{\centering\arraybackslash}X}
\newcolumntype{L}{>{\raggedright\arraybackslash}X}
\newcolumntype{R}{>{\raggedleft\arraybackslash}X}
\usepackage{array}
\usepackage{multicol}
\usepackage{siunitx}
% parse-numbers=false : accepte aussi \SI{2,5 \times 10^{-3}}{...}
\sisetup{locale=UK,output-decimal-marker={.},group-minimum-digits=4,parse-numbers=false}
\DeclareSIUnit\ohm{\ensuremath{\symup{\Omega}}} % U+2126 absent de la police
\DeclareSIUnit\division{div} % réglages d'oscilloscope (ms/div, V/div)
\DeclareSIUnit\tr{tr} % tours (vitesse de rotation, ex. \SI{1500}{\tr\per\minute})
\usepackage{qrcode}
% \degree : commande fréquemment utilisée par les modèles pour un angle ou une
% température en dehors de \SI{}{} (non fournie par les paquets chargés).
\providecommand{\degree}{^{\circ}}
% \qed (fin de démonstration), utilisé par les modèles sans amsthm
\providecommand{\qed}{\hfill$\square$}
% \barre : double barre des valeurs interdites dans un tableau de variations (tkz-tab)
\providecommand{\barre}{\ensuremath{\|}}
% environnement proof (amsthm non chargé) utilisé par les modèles
\ifdefined\proof\else\newenvironment{proof}[1][Proof]{\par\noindent\textit{#1.}\ }{\qed\par}\fi
\usepackage{lastpage}
\usepackage{hyperref}
\hypersetup{hidelinks}

% ============================================================
% Palette « Pop » OnBuch+ — mêmes hex que la classe P (plus_ui.dart)
% ============================================================
\definecolor{poporange}{HTML}{F59321}
\definecolor{popdark}{HTML}{E07A0C}
\definecolor{popcream}{HTML}{FFF6E8}
\definecolor{popink}{HTML}{1C1714}
\definecolor{poppurple}{HTML}{7A5AE0}
\definecolor{popgreen}{HTML}{2E9E6A}
\definecolor{poppink}{HTML}{E0457B}
\definecolor{popblue}{HTML}{2F7DE1}
\definecolor{popgold}{HTML}{C98A12}
\definecolor{popmuted}{HTML}{8A7F76}
\definecolor{popline}{HTML}{EADFD2}
\definecolor{popgray}{HTML}{8A7F76}
\colorlet{popgrey}{popgray}
% Alias tolérants (le modèle nomme parfois les couleurs différemment)
\colorlet{popwhite}{white}
\colorlet{popblack}{popink}
\colorlet{popred}{poppink}
\colorlet{popyellow}{popgold}
% Teintes claires pour fonds de tableaux / zones
\colorlet{popblueL}{popblue!10!white}
\colorlet{poporangeL}{poporange!14!white}
\colorlet{popgreenL}{popgreen!12!white}
\colorlet{poppurpleL}{poppurple!12!white}
\colorlet{poppinkL}{poppink!12!white}
\colorlet{popgoldL}{popgold!14!white}

\pagecolor{popcream}

% ============================================================
% Typographie
% ============================================================
\defaultfontfeatures{Ligatures=TeX}
\setmainfont{PlusJakartaSans-Medium.ttf}[Path=../../../fonts/,BoldFont=PlusJakartaSans-Bold.ttf]
% Maths en sans-serif (Fira Math) avec chiffres et lettres droites de Plus
% Jakarta, pour que les formules se fondent dans le texte.
\usepackage[math-style=ISO,bold-style=ISO]{unicode-math}
\setmathfont{FiraMath-Regular.otf}[Path=../../../fonts/]
\setmathfont{PlusJakartaSans-Medium.ttf}[Path=../../../fonts/,range={up/num,up/{Latin,latin},\mathup}]
% (icomma retiré : décimales avec point en anglais)
% Caractères mathématiques Unicode absents des polices de texte (Plus Jakarta,
% Archivo Black) : rendus par leur équivalent mathématique, en texte comme en
% formule — sinon ils s'affichent en carré vide (ex. « Arithmétique dans ℤ »).
\usepackage{newunicodechar}
\newunicodechar{ℝ}{\ensuremath{\mathbb{R}}}
\newunicodechar{ℤ}{\ensuremath{\mathbb{Z}}}
\newunicodechar{ℂ}{\ensuremath{\mathbb{C}}}
\newunicodechar{ℕ}{\ensuremath{\mathbb{N}}}
\newunicodechar{ℚ}{\ensuremath{\mathbb{Q}}}
\newunicodechar{𝔻}{\ensuremath{\mathbb{D}}}
\newunicodechar{α}{\ensuremath{\alpha}}
\newunicodechar{β}{\ensuremath{\beta}}
\newunicodechar{γ}{\ensuremath{\gamma}}
\newunicodechar{δ}{\ensuremath{\delta}}
\newunicodechar{ε}{\ensuremath{\varepsilon}}
\newunicodechar{θ}{\ensuremath{\theta}}
\newunicodechar{λ}{\ensuremath{\lambda}}
\newunicodechar{μ}{\ensuremath{\mu}}
\newunicodechar{σ}{\ensuremath{\sigma}}
\newunicodechar{τ}{\ensuremath{\tau}}
\newunicodechar{φ}{\ensuremath{\varphi}}
\newunicodechar{ψ}{\ensuremath{\psi}}
\newunicodechar{ω}{\ensuremath{\omega}}
\newunicodechar{Σ}{\ensuremath{\Sigma}}
% majuscules produites par \MakeUppercase dans les titres (α-aminés -> Α)
\newunicodechar{Α}{\ensuremath{\alpha}}
\newunicodechar{Β}{\ensuremath{\beta}}
\newunicodechar{Γ}{\ensuremath{\Gamma}}
\newunicodechar{Δ}{\ensuremath{\Delta}}
\newunicodechar{Ε}{\ensuremath{\varepsilon}}
\newunicodechar{Θ}{\ensuremath{\Theta}}
\newunicodechar{Λ}{\ensuremath{\Lambda}}
\newunicodechar{Μ}{\ensuremath{\mu}}
\newunicodechar{Τ}{\ensuremath{\tau}}
\newunicodechar{Φ}{\ensuremath{\Phi}}
\newunicodechar{Ψ}{\ensuremath{\Psi}}
\newunicodechar{Ω}{\ensuremath{\Omega}}
\newunicodechar{↦}{\ensuremath{\mapsto}}
\newunicodechar{∈}{\ensuremath{\in}}
\newunicodechar{∉}{\ensuremath{\notin}}
\newunicodechar{∧}{\ensuremath{\wedge}}
\newunicodechar{⇔}{\ensuremath{\Leftrightarrow}}
\newunicodechar{⟺}{\ensuremath{\Longleftrightarrow}}
\newunicodechar{⇒}{\ensuremath{\Rightarrow}}
\newunicodechar{⟹}{\ensuremath{\Longrightarrow}}
\newunicodechar{∘}{\ensuremath{\circ}}
\newunicodechar{∪}{\ensuremath{\cup}}
\newunicodechar{∩}{\ensuremath{\cap}}
\newunicodechar{≡}{\ensuremath{\equiv}}
\newunicodechar{⊂}{\ensuremath{\subset}}
\newunicodechar{⊥}{\ensuremath{\perp}}
\newunicodechar{∃}{\ensuremath{\exists}}
\newunicodechar{∀}{\ensuremath{\forall}}
\newunicodechar{∛}{\ensuremath{\sqrt[3]{\,}}}
\newunicodechar{∜}{\ensuremath{\sqrt[4]{\,}}}
\newunicodechar{⃗}{\ensuremath{{}^{\to}}}
\newunicodechar{ⁿ}{\ifmmode^{n}\else\textsuperscript{\ensuremath{n}}\fi}
\newunicodechar{ˣ}{\ifmmode^{x}\else\textsuperscript{\ensuremath{x}}\fi}
\newunicodechar{ᵇ}{\ifmmode^{b}\else\textsuperscript{\ensuremath{b}}\fi}
\newunicodechar{ʳ}{\ifmmode^{r}\else\textsuperscript{\ensuremath{r}}\fi}
\newunicodechar{ᵘ}{\ifmmode^{u}\else\textsuperscript{\ensuremath{u}}\fi}
\newunicodechar{ʸ}{\ifmmode^{y}\else\textsuperscript{\ensuremath{y}}\fi}
\newunicodechar{ᵃ}{\ifmmode^{a}\else\textsuperscript{\ensuremath{a}}\fi}
\newunicodechar{ᵉ}{\ifmmode^{e}\else\textsuperscript{\ensuremath{e}}\fi}
\newunicodechar{ᵏ}{\ifmmode^{k}\else\textsuperscript{\ensuremath{k}}\fi}
\newunicodechar{ᵖ}{\ifmmode^{p}\else\textsuperscript{\ensuremath{p}}\fi}
\newunicodechar{ᴺ}{\ifmmode^{N}\else\textsuperscript{\ensuremath{N}}\fi}
\newunicodechar{ᶜ}{\ifmmode^{c}\else\textsuperscript{\ensuremath{c}}\fi}
\newunicodechar{ʰ}{\ifmmode^{h}\else\textsuperscript{\ensuremath{h}}\fi}
\newunicodechar{ᵠ}{\ifmmode^{q}\else\textsuperscript{\ensuremath{q}}\fi}
\newunicodechar{ₙ}{\ifmmode_{n}\else\textsubscript{\ensuremath{n}}\fi}
\newunicodechar{ₐ}{\ifmmode_{a}\else\textsubscript{\ensuremath{a}}\fi}
\newunicodechar{ᵢ}{\ifmmode_{i}\else\textsubscript{\ensuremath{i}}\fi}
\newunicodechar{ᵣ}{\ifmmode_{r}\else\textsubscript{\ensuremath{r}}\fi}
\newunicodechar{ₖ}{\ifmmode_{k}\else\textsubscript{\ensuremath{k}}\fi}
\newunicodechar{ᵦ}{\ifmmode_{\beta}\else\textsubscript{\ensuremath{\beta}}\fi}
\newunicodechar{ₓ}{\ifmmode_{x}\else\textsubscript{\ensuremath{x}}\fi}
\newfontfamily\popdisplay{ArchivoBlack-Regular.ttf}[Path=../../../fonts/]
\newfontfamily\poplogo{SpaceGrotesk-Bold.ttf}[Path=../../../fonts/]
\newfontfamily\popsemi{PlusJakartaSans-SemiBold.ttf}[Path=../../../fonts/]
\newfontfamily\popxbold{PlusJakartaSans-ExtraBold.ttf}[Path=../../../fonts/]
\newfontfamily\popxboldls{PlusJakartaSans-ExtraBold.ttf}[Path=../../../fonts/,LetterSpace=3.0]

\setlist[enumerate]{leftmargin=1.7em,itemsep=5pt,topsep=4pt}
\setlist[itemize]{leftmargin=1.7em,itemsep=4pt,topsep=4pt,label={\color{poporange}\textbullet}}
\setlength{\parindent}{0pt}
\setlength{\parskip}{5pt}
\renewcommand{\arraystretch}{1.25}

% pgfplots : style Pop par défaut
\pgfplotsset{
  every axis/.append style={axis line style={popink,line width=1pt},tick style={popink},
    grid style={popline,line width=0.5pt},label style={font=\small},tick label style={font=\footnotesize}},
  popcurve/.style={poporange,line width=1.8pt,no markers,smooth},
}
% Même style disponible dans un \draw TikZ simple (hors pgfplots)
\tikzset{popcurve/.style={poporange,line width=1.8pt}}
\tikzset{popgrid/.style={popline,very thin}}
\colorlet{popcurve}{poporange} % le nom du style est parfois utilisé comme couleur

% ============================================================
% Pill / squiggle
% ============================================================
\newcommand{\poppill}[3]{%
  \tikz[baseline=-0.55ex]\node[draw=popink,line width=1.2pt,rounded corners=2.6pt,%
    fill=#2,inner xsep=6.5pt,inner ysep=3pt]{%
    \popxboldls\fontsize{7}{7}\selectfont\color{#3}\MakeUppercase{#1}};%
}
\newcommand{\popsquiggle}[1]{%
  \tikz[baseline=-0.4ex]{
    \draw[#1,line width=2.6pt,line cap=round]
      plot [smooth] coordinates {(0,0.09) (0.34,0.01) (0.68,0.21) (1.02,0.01) (1.36,0.10)};
  }%
}
% Mot-clé surligné (marqueur orange sous le texte)
\newcommand{\cle}[1]{{\popxbold\color{popdark}#1}}

% ============================================================
% En-tête / pied
% ============================================================
\pagestyle{fancy}
\fancyhf{}
\renewcommand{\headrulewidth}{0pt}
\renewcommand{\footrulewidth}{0pt}
\lhead{\raisebox{-0.15ex}{\poplogo\fontsize{13}{13}\selectfont\color{popink}OnBuch\color{poporange}+}}
\rhead{\poppill{\DOCMATIERE\ \textbullet\ \DOCNIVEAU\ \textbullet\ Chapter \DOCLECON}{white}{popink}}
\lfoot{\popxbold\fontsize{7.6}{7.6}\selectfont\color{popmuted}\MakeUppercase{Signed course OnBuch+}\ \textbullet\ {\popsemi\DOCTITRE}}
\rfoot{\tikz[baseline=-0.6ex]\node[rounded corners=8.5pt,fill=popink,minimum height=17pt,minimum width=17pt,inner xsep=5pt,inner sep=0]{\popxbold\fontsize{8}{8}\selectfont\color{popcream}\thepage\,/\,\pageref*{LastPage}};}

% ============================================================
% Titres
% ============================================================
\newcommand{\chaptitre}[2]{%
  \begin{tcolorbox}[enhanced,colback=poporange,colframe=popink,boxrule=2.6pt,arc=11pt,
      drop shadow=popink,left=15pt,right=15pt,top=12pt,bottom=11pt,before skip=3pt,after skip=18pt]
    \poppill{#2}{popink}{popcream}\par\vspace{9pt}
    {\raggedright\hyphenpenalty=10000\exhyphenpenalty=10000\popdisplay\fontsize{22}{24}\selectfont\color{popcream}\MakeUppercase{#1}\par}
  \end{tcolorbox}
}
% Section numérotée : pastille orange avec le numéro + titre Archivo
\newcounter{popsec}
\newcommand{\coursec}[1]{%
  \stepcounter{popsec}\setcounter{popsub}{0}%
  \par\vspace{16pt}\noindent
  \begin{minipage}{\linewidth}
  \tikz[baseline=-0.7ex]\node[draw=popink,line width=1.6pt,fill=poporange,rounded corners=4pt,minimum size=22pt,inner sep=0]{\popdisplay\fontsize{12}{12}\selectfont\color{popcream}\thepopsec};%
  \hspace{8pt}{\popdisplay\fontsize{15}{17}\selectfont\color{popink}\MakeUppercase{#1}}\par
  \vspace{4pt}\noindent{\color{poporange}\rule{36pt}{3.6pt}}
  \end{minipage}\par\nopagebreak\vspace{8pt}
}
\newcounter{popsub}
\newcommand{\courssub}[1]{%
  \stepcounter{popsub}%
  \par\vspace{9pt}\noindent{\popxbold\fontsize{11.5}{13}\selectfont\color{popink}{\color{poporange}\thepopsec.\thepopsub}\hspace{6pt}#1}\par\nopagebreak\vspace{4pt}
}

% ============================================================
% Boîtes pédagogiques — même recette : fond blanc, bordure épaisse colorée,
% ombre dure encre, badge plein. Titre optionnel [#1].
% ============================================================
\tcbset{popbase/.style={breakable,enhanced,colback=white,boxrule=2.3pt,arc=9pt,
  shadow={3pt}{-3pt}{0pt}{popink},left=14pt,right=14pt,top=11pt,bottom=11pt,before skip=11pt,after skip=15pt,
  fontupper=\popsemi\fontsize{10.6}{14.5}\selectfont\color{popink},
  fontlower=\popsemi\fontsize{10.6}{14.5}\selectfont\color{popink}}}
% Coche dessinée (le glyphe ✓ n'existe pas dans les polices du document)
\newcommand{\popcheck}[1]{\tikz[baseline=-0.5ex]\draw[#1,line width=1.6pt,line cap=round,line join=round](-0.13,0.0)--(-0.04,-0.09)--(0.13,0.1);}
\providecommand{\coche}{\popcheck{popgreen}}
\providecommand{\resultat}[1]{\par\smallskip\noindent\fcolorbox{popgreen}{popgreenL}{\parbox{\dimexpr\linewidth-2\fboxsep-2\fboxrule\relax}{\textbf{Result.} #1}}\par\smallskip}
\newcommand{\popboxhead}[3]{\poppill{#1}{#2}{white}\ifx\relax#3\relax\else\hspace{6pt}{\popxbold\fontsize{10.6}{12}\selectfont\color{popink}#3}\fi\par\vspace{7pt}}

\newtcolorbox{aretenir}[1][]{popbase,colframe=popgreen,before upper={\popboxhead{Key points}{popgreen}{#1}}}
\newtcolorbox{exemplebox}[1][]{popbase,colframe=poporange,before upper={\popboxhead{Exemple}{poporange}{#1}}}
\newtcolorbox{definition}[1][]{popbase,colframe=popblue,before upper={\popboxhead{Definition}{popblue}{#1}}}
\newtcolorbox{propriete}[1][]{popbase,colframe=poppurple,before upper={\popboxhead{Property}{poppurple}{#1}}}
\newtcolorbox{methode}[1][]{popbase,colframe=popgold,before upper={\popboxhead{Method}{popgold}{#1}}}
\newtcolorbox{attention}[1][]{popbase,colframe=poppink,before upper={\popboxhead{Warning}{poppink}{#1}}}
\newtcolorbox{experience}[1][]{popbase,colframe=popblue,colback=popblueL,before upper={\popboxhead{Experiment}{popblue}{#1}}}
% « Le savais-tu ? » : carte violette pleine (contexte, culture, Cameroun)
\newtcolorbox{savaistu}[1][]{popbase,colback=poppurple,colframe=popink,
  fontupper=\popsemi\fontsize{10.4}{14.2}\selectfont\color{white},
  before upper={\poppill{Did you know?}{white}{poppurple}\ifx\relax#1\relax\else\hspace{6pt}{\popxbold\color{white}#1}\fi\par\vspace{7pt}}}
% Exercice résolu : énoncé en haut, \tcblower puis la solution sur fond crème
\newcounter{popexo}\newcounter{popexr}
\newtcolorbox{exoresolu}[1][]{popbase,colframe=poporange,colbacklower=popcream,
  segmentation style={popink,line width=1.2pt,dashed},
  before upper={\stepcounter{popexr}\popboxhead{Worked example \thepopexr}{poporange}{#1}},
  before lower={\poppill{Solution}{popink}{popcream}\par\vspace{6pt}}}
% Exercice (non résolu en place) — la correction est dans la partie Corrigés
\newtcolorbox{exercice}[1][]{popbase,colframe=popink,
  before upper={\stepcounter{popexo}\popboxhead{Exercise \thepopexo}{popink}{#1}}}
% Corrigé
\newtcolorbox{corrige}[1][]{popbase,colframe=popgreen,colback=popgreenL,
  before upper={\popboxhead{Solution}{popgreen}{#1}}}
% Objectifs / prérequis / bilan
\newtcolorbox{objectifs}{popbase,colframe=popink,colback=poporangeL,
  before upper={\popboxhead{Chapter objectives}{popink}{}}}
\newtcolorbox{prerequis}{popbase,colframe=popink,colback=popblueL,
  before upper={\popboxhead{Prerequisites}{popblue}{}}}
\newtcolorbox{bilan}{popbase,colframe=popink,colback=white,boxrule=2.8pt,
  before upper={\popboxhead{Fiche bilan}{poporange}{L'essentiel à connaître pour l'examen}}}
% Figure encadrée avec légende
\newtcolorbox{popfigure}[1][]{enhanced,breakable=false,colback=white,colframe=popline,boxrule=1.4pt,arc=8pt,
  left=10pt,right=10pt,top=10pt,bottom=8pt,before skip=10pt,after skip=12pt,halign=center,
  lower separated=false,halign lower=center,
  fontlower=\popsemi\fontsize{9}{11.5}\selectfont\color{popmuted},
  #1}
\newcommand{\legende}[1]{\par\vspace{6pt}{\popsemi\fontsize{9}{11.5}\selectfont\color{popmuted}#1}}

% ============================================================
% Signature OnBuch+
% ============================================================
\newcommand{\poplogoinline}[1]{{\poplogo\fontsize{#1}{#1}\selectfont\color{popink}OnBuch\color{poporange}+}}
\newcommand{\signatureonbuch}{%
  \begin{tcolorbox}[enhanced,colback=popink,colframe=popink,boxrule=2.2pt,arc=10pt,
      drop shadow=poporange,left=14pt,right=14pt,top=10pt,bottom=10pt,before skip=0pt,after skip=0pt]
    \tikz[baseline=-0.7ex]\node[circle,fill=popgreen,draw=popcream,line width=1.3pt,minimum size=22pt,inner sep=0]{\popcheck{white}};%
    \hspace{9pt}\begin{minipage}[c]{0.62\linewidth}
      {\popxbold\fontsize{12}{13}\selectfont\color{popcream}Signed\ \textbullet\ {\poplogo OnBuch\color{poporange}+}}\par\vspace{2pt}
      {\raggedright\popsemi\fontsize{8}{10.5}\selectfont\color{popline}\MakeUppercase{OnBuch+ teaching team}\\\MakeUppercase{Follows the official MINESEC syllabus}\par}
    \end{minipage}\hfill
    {\poplogo\fontsize{22}{22}\selectfont\color{popcream}OnBuch\color{poporange}+}
  \end{tcolorbox}%
}

% ============================================================
% Page de garde
% #1 titre · #2 matière · #3 niveau · #4 module · #5 n° leçon · #6 durée ·
% #7 description / objectifs (texte ou itemize)
% ============================================================
\newcommand{\pagegarde}[7]{%
  \thispagestyle{empty}
  \newgeometry{a4paper,margin=1.7cm,top=1.6cm,bottom=1.4cm}%
  \begin{tikzpicture}[remember picture,overlay]
    \fill[poporange!18] ([xshift=-3.2cm,yshift=-3.6cm]current page.north east) circle (4.6cm);
    \fill[poppurple!14] ([xshift=2.4cm,yshift=7.6cm]current page.south west) circle (3.8cm);
  \end{tikzpicture}%
  {\poplogo\fontsize{30}{30}\selectfont\color{popink}OnBuch\color{poporange}+}\hfill
  \raisebox{0.6ex}{\poppill{Signed course \textbullet\ Premium}{popink}{popcream}}\par
  \vspace{1pt}\popsquiggle{poppurple}\par
  \vspace{8pt}
  \begin{tcolorbox}[enhanced,colback=poporange,colframe=popink,boxrule=2.8pt,arc=13pt,
      drop shadow=popink,left=18pt,right=18pt,top=12pt,bottom=13pt,after skip=10pt]
    \poppill{#2 \textbullet\ #3}{popink}{popcream}\hspace{5pt}\poppill{#4}{white}{popink}\par\vspace{8pt}
    {\popdisplay\fontsize{14}{14}\selectfont\color{popink}CHAPTER #5}\par\vspace{4pt}
    {\raggedright\hyphenpenalty=10000\exhyphenpenalty=10000\popdisplay\fontsize{25}{27}\selectfont\color{popcream}\MakeUppercase{#1}\par}
    \vspace{8pt}
    {\popxbold\fontsize{9.5}{9.5}\selectfont\color{popink}\MakeUppercase{Suggested time: #6}}
  \end{tcolorbox}
  \begin{tcolorbox}[enhanced,colback=white,colframe=popink,boxrule=2.2pt,arc=10pt,
      drop shadow=popink,left=16pt,right=16pt,top=10pt,bottom=10pt,after skip=8pt]
    \poppill{In this chapter}{poppurple}{white}\par\vspace{5pt}
    \popsemi\fontsize{9.3}{12.6}\selectfont\color{popink}#7
  \end{tcolorbox}
  \noindent\begin{minipage}[t]{0.487\linewidth}
    \begin{tcolorbox}[enhanced,colback=popgreen,colframe=popink,boxrule=2.2pt,arc=10pt,
        drop shadow=popink,left=11pt,right=11pt,top=10pt,bottom=10pt,height=3.1cm,valign=center]
      \tcbox[colback=white,colframe=popink,boxrule=1.3pt,arc=5pt,boxsep=0pt,nobeforeafter,
        left=3pt,right=3pt,top=3pt,bottom=3pt,valign=center]{\qrcode[height=1.9cm]{https://play.google.com/store/apps/details?id=com.luvvix.onbuch&hl=en}}%
      \hspace{8pt}\begin{minipage}[c]{0.5\linewidth}
        {\popdisplay\fontsize{11}{12.5}\selectfont\color{white}\MakeUppercase{Download OnBuch}}\par\vspace{4pt}
        {\raggedright\popsemi\fontsize{8.4}{11}\selectfont\color{white}Free \textbullet\ Google Play\\AI tutor, past papers, results\par}
      \end{minipage}
    \end{tcolorbox}
  \end{minipage}\hfill
  \begin{minipage}[t]{0.487\linewidth}
    \begin{tcolorbox}[enhanced,colback=poppurple,colframe=popink,boxrule=2.2pt,arc=10pt,
        drop shadow=popink,left=11pt,right=11pt,top=10pt,bottom=10pt,height=3.1cm,valign=center]
      \tikz[baseline=-0.7ex]\node[circle,fill=white,draw=popink,line width=1.3pt,minimum size=1.6cm,inner sep=0]{\popdisplay\fontsize{24}{24}\selectfont\color{poppurple}+};%
      \hspace{8pt}\begin{minipage}[c]{0.6\linewidth}
        {\popdisplay\fontsize{11}{12.5}\selectfont\color{white}\MakeUppercase{Go OnBuch+}}\par\vspace{4pt}
        {\raggedright\popsemi\fontsize{8.4}{11}\selectfont\color{white}All signed courses, unlimited AI tutor, PDF sheets — OnBuch+ tab in the app\par}
      \end{minipage}
    \end{tcolorbox}
  \end{minipage}
  \vspace{8pt}
  \popcontenu
  \vfill
  \signatureonbuch
  \restoregeometry
  \newpage
}

% Rangée « Ce cours contient » (page de garde)
\newcommand{\poptile}[3]{% #1 couleur #2 gros texte #3 légende
  \begin{tcolorbox}[enhanced,colback=#1,colframe=popink,boxrule=2pt,arc=9pt,shadow={2.5pt}{-2.5pt}{0pt}{popink},
      left=6pt,right=6pt,top=9pt,bottom=9pt,halign=center,valign=center,height=1.75cm,width=\linewidth,nobeforeafter]
    {\popdisplay\fontsize{12.5}{13}\selectfont\color{white}\MakeUppercase{#2}}\par\vspace{3pt}
    {\centering\popsemi\fontsize{7.8}{9.5}\selectfont\color{white}#3\par}
  \end{tcolorbox}}
\newcommand{\popcontenu}{%
  \par\noindent{\popxboldls\fontsize{7.5}{7.5}\selectfont\color{popmuted}THIS SIGNED COURSE CONTAINS}\par\vspace{7pt}
  \noindent\begin{minipage}[t]{0.235\linewidth}\poptile{popblue}{Course}{complete, illustrated, step by step}\end{minipage}\hfill
  \begin{minipage}[t]{0.235\linewidth}\poptile{poporange}{Methods}{and worked examples}\end{minipage}\hfill
  \begin{minipage}[t]{0.235\linewidth}\poptile{poppink}{Exercises}{exam style + solutions}\end{minipage}\hfill
  \begin{minipage}[t]{0.235\linewidth}\poptile{popgreen}{Summary sheet}{the essentials to remember}\end{minipage}\par}

% Sommaire
\newtcolorbox{sommaire}{enhanced,colback=white,colframe=popink,boxrule=2.2pt,arc=10pt,
  drop shadow=popink,left=18pt,right=18pt,top=13pt,bottom=13pt,before skip=6pt,after skip=20pt,
  before upper={\poppill{Contents}{poppurple}{white}\par\vspace{9pt}\popsemi\fontsize{10.6}{14.5}\selectfont\color{popink}}}

% Signature de fin de cours — poussée tout en bas de la dernière page
\newcommand{\findecours}{%
  \par\vfill\nopagebreak
  \begin{center}\popsquiggle{poporange}\end{center}\vspace{4pt}
  \signatureonbuch
}
\AtBeginDocument{\def\llbracket{\mathopen{[\mkern-3.5mu[}}\def\rrbracket{\mathclose{]\mkern-3.5mu]}}} % intervalles d'entiers (glyphe ⟦ absent de la police)
% Glyphes absents de Fira Math : redéfinis à partir de glyphes disponibles
\AtBeginDocument{%
  \def\setminus{\mathbin{\backslash}}\let\smallsetminus\setminus
  \def\vdots{\mathord{\vbox{\baselineskip3.2pt\lineskiplimit0pt\kern4pt\hbox{.}\hbox{.}\hbox{.}}}}%
  \def\ddots{\mathinner{\mkern1mu\raise6.5pt\hbox{.}\mkern2mu\raise3.6pt\hbox{.}\mkern2mu\raise0.7pt\hbox{.}\mkern1mu}}%
  \def\triangle{\mathord{\scalebox{0.9}{$\symup{\Delta}$}}}%
  \def\nabla{\mathord{\raisebox{\depth}{\scalebox{1}[-1]{$\symup{\Delta}$}}}}%
  \def\checkmark{\popcheck{popdark}}%
  \def\star{\mathbin{\ast}}\def\bigstar{\mathord{\ast}}%
  \def\diamond{\mathbin{\scalebox{0.6}{$\Diamond$}}}%
  \def\bigcup{\mathop{\mathchoice{\raisebox{-0.3em}{\scalebox{1.7}{$\cup$}}}{\scalebox{1.25}{$\cup$}}{\cup}{\cup}}\displaylimits}%
  \def\bigcap{\mathop{\mathchoice{\raisebox{-0.3em}{\scalebox{1.7}{$\cap$}}}{\scalebox{1.25}{$\cap$}}{\cap}{\cap}}\displaylimits}%
  \def\lesseqqgtr{\mathrel{\lessgtr}}\def\bigoplus{\mathop{\oplus}\displaylimits}%
}
\providecommand{\ssi}{if and only if} % abréviation fréquente
\AtBeginDocument{\providecommand{\wideparen}{\overparen}} % arc de cercle

% Fonctions usuelles que les modèles utilisent sans que amsmath les fournisse
\providecommand{\arsinh}{\operatorname{arsinh}}\providecommand{\arcosh}{\operatorname{arcosh}}
\providecommand{\artanh}{\operatorname{artanh}}\providecommand{\arsech}{\operatorname{arsech}}
\providecommand{\arcsch}{\operatorname{arcsch}}\providecommand{\arcoth}{\operatorname{arcoth}}
\providecommand{\arcsinh}{\operatorname{arsinh}}\providecommand{\arccosh}{\operatorname{arcosh}}
\providecommand{\arctanh}{\operatorname{artanh}}\providecommand{\sech}{\operatorname{sech}}
\providecommand{\csch}{\operatorname{csch}}\providecommand{\arcsec}{\operatorname{arcsec}}
\providecommand{\arccsc}{\operatorname{arccsc}}\providecommand{\arccot}{\operatorname{arccot}}
\providecommand{\sgn}{\operatorname{sgn}}\providecommand{\lcm}{\operatorname{lcm}}
\providecommand{\Var}{\operatorname{Var}}\providecommand{\Cov}{\operatorname{Cov}}

%% FIGFIX-BEGIN v4 — correctifs globaux des figures (docs/onbuch-generation/tools/figfix)
\usepackage{adjustbox}\usepackage{etoolbox}
% toute figure TikZ/pgfplots plus large que la ligne est réduite à la largeur disponible
\BeforeBeginEnvironment{tikzpicture}{\begin{adjustbox}{max width=\linewidth}}
\AfterEndEnvironment{tikzpicture}{\end{adjustbox}}
% légendes pgfplots lisibles par-dessus les courbes
\pgfplotsset{every axis/.append style={legend style={fill=white,fill opacity=0.92,text opacity=1,draw=black!15,inner sep=3pt}}}
% étiquettes d'arêtes (syntaxe edge["..."]) avec fond blanc
\tikzset{every edge quotes/.append style={fill=white,inner sep=1.2pt}}
%% FIGFIX-END
\begin{document}

\pagegarde{Regression and Correlation}{Pure Mathematics with Statistics}{Upper Sixth}{Upper Sixth — Pure mathematics with statistics}{9}{15 periods}{This chapter introduces regression and correlation, the statistical study of relationships between two variables. Students learn to draw scatter diagrams, fit regression lines by Mayer's method and by least squares, compute residual sums of squares, and measure association using the product moment, Spearman's and Kendall's correlation coefficients.
\vspace{4pt}
\begin{multicols}{2}\small
\begin{itemize}
\item By the end of this chapter I can interpret bivariate data and explain the idea of regression.
\item By the end of this chapter I can draw and interpret a scatter diagram and sketch regression lines.
\item By the end of this chapter I can calculate a regression line using Mayer's method.
\item By the end of this chapter I can compute sample variance and covariance from raw data.
\item By the end of this chapter I can determine the least squares regression lines of y on x and of x on y.
\item By the end of this chapter I can compute and interpret the residual sum of squares for each regression line.
\end{itemize}\end{multicols}}

\begin{sommaire}
\begin{multicols}{2}
\begin{enumerate}
\item Introduction to Regression and Scatter Diagrams
\item Mayer's Method for the Regression Line
\item Correlation, Sample Variance and Covariance
\item Least Squares Regression Lines
\item Residual Sum of Squares and the Line of Best Fit
\item The Product Moment Correlation Coefficient
\item Rank Correlation: Spearman and Kendall
\item Integration activity
\item Methods and worked examples
\item Exercises
\item Solutions to the exercises
\item Summary sheet
\end{enumerate}
\end{multicols}
\end{sommaire}

\begin{objectifs}
\begin{itemize}
\item By the end of this chapter I can interpret bivariate data and explain the idea of regression.
\item By the end of this chapter I can draw and interpret a scatter diagram and sketch regression lines.
\item By the end of this chapter I can calculate a regression line using Mayer's method.
\item By the end of this chapter I can compute sample variance and covariance from raw data.
\item By the end of this chapter I can determine the least squares regression lines of y on x and of x on y.
\item By the end of this chapter I can compute and interpret the residual sum of squares for each regression line.
\item By the end of this chapter I can identify and use the line of best fit.
\item By the end of this chapter I can calculate and interpret the product moment correlation coefficient.
\item By the end of this chapter I can rank data and compute Spearman's and Kendall's rank correlation coefficients, stating their uses and limitations.
\end{itemize}
\end{objectifs}

\begin{prerequis}
\begin{itemize}
\item Mean, variance and standard deviation of a data set (Lower Sixth statistics)
\item Equation of a straight line, gradient and intercept (Pure Mathematics)
\item Summation notation Σ and manipulation of sums
\item Plotting points and reading graphs in the Cartesian plane
\item Basic idea of a function and of dependent/independent variables
\end{itemize}
\end{prerequis}

\newpage

\coursec{Introduction to Regression and Scatter Diagrams}

Every market day in Bamenda, Mama Ngo notices that the more bags of Irish potatoes she sells, the more money she takes home --- but the relationship is never perfectly exact. At the Douala port, engineers observe that heavier trucks tend to consume more fuel per trip, and they want to predict fuel needs from truck weight. During the GCE mock exams, a principal wonders whether students who rank high in Mathematics also rank high in Physics. Can we measure how strongly two quantities are linked, and can we use one to predict the other? This chapter gives you the statistical tools --- \cle{regression lines} and \cle{correlation coefficients} --- to answer such questions with confidence.

The common feature of all these situations is that \emph{two} quantities are measured on the \emph{same} individual or item, and we suspect they move together. Our first task is to organise such information, display it on a diagram, and learn to read the story the diagram tells.

\courssub{Bivariate Data}

In Lower Sixth statistics you studied \emph{univariate} data: a single variable measured on each individual (the heights of students, the marks in one test). Real life, however, constantly pairs measurements.

\begin{definition}[Bivariate data]
\cle{Bivariate data} consist of pairs of observations $(x_i, y_i)$, where both values are measured on the \emph{same} individual or item. If there are $n$ individuals, the data set is
\[
(x_1, y_1),\ (x_2, y_2),\ \ldots,\ (x_n, y_n).
\]
The two variables $x$ and $y$ may be of the same nature (two exam marks) or of different natures (rainfall in mm and yield in tonnes).
\end{definition}

\begin{exemplebox}[Examples of bivariate data]
\begin{itemize}
\item For each district of the North West Region: $x = $ annual rainfall (mm), $y = $ maize yield (tonnes per hectare).
\item For each Upper Sixth student: $x = $ hours of study per week, $y = $ score in the GCE mock examination.
\item For each truck leaving the Douala port: $x = $ loaded weight (tonnes), $y = $ fuel consumed (litres).
\end{itemize}
\end{exemplebox}

\courssub{Independent and Dependent Variables}

When we suspect that one variable \emph{explains} or \emph{drives} the other, we give the two variables different roles.

\begin{definition}[Explanatory and response variables]
The \cle{independent variable} (also called the \cle{explanatory} or controlled variable) is the one we believe influences the other, or the one we can fix or observe first. The \cle{dependent variable} (also called the \cle{response} variable) is the one we wish to explain or predict. By universal convention the independent variable is denoted $x$ and the dependent variable $y$.
\end{definition}

Rainfall is the independent variable and crop yield the dependent one: the rain is not caused by the maize! Similarly, hours of study ($x$) may influence the mock score ($y$), not the reverse.

\begin{attention}[The independent variable goes on the horizontal axis]
When plotting, the \textbf{independent variable $x$ is always placed on the horizontal axis} and the dependent variable $y$ on the vertical axis. Examiners penalise reversed axes, and reversing the axes changes the whole regression analysis (we shall see that the line of $y$ on $x$ is \emph{not} the same as the line of $x$ on $y$).
\end{attention}

\courssub{The Idea of Regression}

Mama Ngo's takings are not an exact multiple of the bags sold: prices fluctuate, some customers bargain, some bags are heavier. The relationship is \emph{statistical}, not exact. We therefore ask a weaker but very useful question: \emph{on average}, what value of $y$ corresponds to a given value of $x$?

\begin{definition}[Regression]
\cle{Regression} is the study of how the \emph{average} (expected) value of one variable depends on the value of another variable. A regression model allows us to \textbf{estimate or predict} the response variable $y$ from a known value of the explanatory variable $x$.
\end{definition}

The simplest and most important model is \cle{linear regression}, in which we assume the average value of $y$ is a linear function of $x$:
\[
y = a + bx,
\]
where $b$ is the \cle{gradient} (the average change in $y$ per unit increase in $x$) and $a$ is the \cle{intercept} (the value of $y$ predicted when $x = 0$). The observed points will not lie exactly on this line; each observation differs from the line by a small \emph{error} or \emph{residual}, and a good line makes these residuals small. Later sections will show how to compute $a$ and $b$ precisely (Mayer's method, then the method of least squares). For now we learn to \emph{see} the line.

\begin{savaistu}[Why the strange word ``regression''?]
The term was introduced by Sir Francis Galton in the 19th century while studying the heights of parents and their children. He noticed that very tall parents tended to have children who were tall, but on average \emph{less} tall than themselves --- the children's heights ``regressed'' towards the mean of the population. The name stuck, and today ``regression'' simply means modelling how one variable depends on another.
\end{savaistu}

\courssub{The Scatter Diagram}

Before any calculation, \textbf{always plot the data}. A picture reveals the type and strength of the relationship instantly.

\begin{definition}[Scatter diagram]
A \cle{scatter diagram} (or scatter plot) is a graph of the $n$ points $(x_i, y_i)$ of a bivariate data set, with the independent variable $x$ on the horizontal axis and the dependent variable $y$ on the vertical axis. No lines join the points: each point represents one individual.
\end{definition}

\begin{methode}[Drawing and interpreting a scatter diagram]
\begin{enumerate}
\item \textbf{Identify the variables.} Decide which is independent ($x$) and which is dependent ($y$).
\item \textbf{Draw and label the axes.} Place $x$ horizontally, $y$ vertically; write the variable names and units.
\item \textbf{Choose scales.} Each axis should cover the range of its data (the axes need not start at 0, and the two scales need not be equal). Use as much of the page as possible.
\item \textbf{Plot the points.} Mark each pair $(x_i, y_i)$ with a small clear dot or cross.
\item \textbf{Interpret the pattern.} Look at the \emph{direction} (rising or falling), the \emph{form} (linear or curved) and the \emph{strength} (points tightly or loosely clustered around an imaginary line).
\end{enumerate}
\end{methode}

\begin{exoresolu}[Hours of study and GCE mock score]
Ten Upper Sixth students reported their weekly study hours $x$ and obtained the following mock scores $y$ (out of 100):
\[
\begin{array}{c|cccccccccc}
x & 4 & 6 & 8 & 9 & 11 & 13 & 15 & 17 & 19 & 21 \\ \hline
y & 38 & 45 & 50 & 55 & 58 & 63 & 70 & 72 & 80 & 84
\end{array}
\]
Draw a scatter diagram and comment.
\tcblower
\textbf{Solution.} Study hours ($x$) is the independent variable (horizontal axis); the score ($y$) is the response (vertical axis). We scale the $x$-axis from 0 to 22 and the $y$-axis from 30 to 90, then plot the ten points (see the figure below). The points rise steadily from bottom-left to top-right and cluster fairly closely around an imaginary straight line: there is a \textbf{strong positive linear association} between study time and mock score. A straight-line model $y = a + bx$ appears appropriate.
\end{exoresolu}

\begin{popfigure}
\begin{tikzpicture}[scale=0.42]
\draw[->, thick, popink] (0,0) -- (23,0) node[below left, popink] {\small $x$ (hours)};
\draw[->, thick, popink] (0,0) -- (0,11.5) node[below left, popink] {\small $y$ (score)};
\foreach \x in {0,4,8,12,16,20} \draw[popline] (\x,0.08) -- (\x,-0.08) node[below, popmuted] {\scriptsize \x};
\foreach \y/\lab in {1/40,3/50,5/60,7/70,9/80} \draw[popline] (0.08,\y) -- (-0.08,\y) node[left, popmuted] {\scriptsize \lab};
\foreach \x/\y in {4/0.8, 6/1.5, 8/2.0, 9/2.5, 11/2.8, 13/3.3, 15/4.0, 17/4.2, 19/5.0, 21/5.4}
  \fill[popblue] (\x,\y) circle (0.14);
\draw[dashed, poporange, thick] (2,0.35) -- (22,5.75);
\node[poporange, align=center] at (17.5,1.1) {\scriptsize trend by eye};
\end{tikzpicture}
\legende{Scatter diagram of mock score against weekly study hours for 10 students: the points suggest a strong positive linear association.}
\end{popfigure}

\courssub{Interpreting a Scatter Diagram}

Three typical patterns occur.

\begin{definition}[Types of association]
\begin{itemize}
\item \cle{Positive association}: as $x$ increases, $y$ tends to increase; the points cluster around a line with positive gradient.
\item \cle{Negative association}: as $x$ increases, $y$ tends to decrease; the points cluster around a line with negative gradient.
\item \cle{No apparent association}: the points are scattered at random; knowing $x$ gives no information about $y$.
\end{itemize}
The association is \cle{linear} if the points cluster around a straight line, and \cle{non-linear} if they follow a curve (for example a parabola or an exponential shape). Association is described as \emph{strong} when the points lie close to the trend and \emph{weak} when they are widely spread around it.
\end{definition}

\begin{popfigure}
\begin{tikzpicture}
\begin{axis}[
  name=plot1, width=5.6cm, height=4.6cm,
  title={\small Positive}, title style={color=popink},
  xtick=\empty, ytick=\empty, xmin=0, xmax=10, ymin=0, ymax=10,
  axis lines=left, axis line style={popink}]
\addplot[only marks, mark=*, mark size=1.4pt, popblue] coordinates {(1,1.6) (2,2.6) (3,2.7) (4,4.1) (5,4.6) (6,5.9) (7,6.4) (8,7.6) (9,8.2)};
\end{axis}
\begin{axis}[
  name=plot2, at={($(plot1.east)+(1.1cm,0)$)}, anchor=west,
  width=5.6cm, height=4.6cm,
  title={\small Negative}, title style={color=popink},
  xtick=\empty, ytick=\empty, xmin=0, xmax=10, ymin=0, ymax=10,
  axis lines=left, axis line style={popink}]
\addplot[only marks, mark=*, mark size=1.4pt, poporange] coordinates {(1,8.6) (2,7.9) (3,7.2) (4,6.3) (5,5.8) (6,4.9) (7,4.1) (8,3.0) (9,2.2)};
\end{axis}
\begin{axis}[
  name=plot3, at={($(plot2.east)+(1.1cm,0)$)}, anchor=west,
  width=5.6cm, height=4.6cm,
  title={\small No correlation}, title style={color=popink},
  xtick=\empty, ytick=\empty, xmin=0, xmax=10, ymin=0, ymax=10,
  axis lines=left, axis line style={popink}]
\addplot[only marks, mark=*, mark size=1.4pt, popgreen] coordinates {(1,4.2) (2,7.1) (3,2.3) (4,6.2) (5,3.5) (6,7.9) (7,2.8) (8,5.6) (9,4.8)};
\end{axis}
\end{tikzpicture}
\legende{Three scatter diagrams: positive association (left), negative association (centre), no apparent association (right).}
\end{popfigure}

\begin{attention}[Association is not causation]
A scatter diagram can show that two variables \emph{move together}, but it can \textbf{never prove} that one \emph{causes} the other. Ice-cream sales and drowning incidents may both rise in the same months --- not because ice cream causes drowning, but because a hidden third factor (hot weather) drives both. Such a hidden variable is called a \cle{lurking} (or confounding) \cle{variable}. Always ask: is there a plausible mechanism, or could a third factor explain the pattern?
\end{attention}

\courssub{Sketching Regression Lines by Eye; Outliers}

Once a linear pattern is visible, we can draw a first approximation of the regression line \emph{by eye}: a straight line that follows the trend of the cloud of points, with roughly as many points above it as below it, and passing as close as possible to the \textbf{mean point} $(\bar{x}, \bar{y})$, where
\[
\bar{x} = \frac{1}{n}\sum_{i=1}^{n} x_i, \qquad \bar{y} = \frac{1}{n}\sum_{i=1}^{n} y_i.
\]
A fundamental fact, which we shall prove when studying least squares, is that \textbf{every least squares regression line passes through $(\bar{x}, \bar{y})$}. The mean point is therefore the natural ``anchor'' of any line drawn by eye.

\begin{definition}[Outlier]
An \cle{outlier} is a point that lies far from the pattern followed by the rest of the data. Outliers may come from recording errors or from genuinely unusual individuals; they can strongly pull a regression line towards themselves, so they must always be spotted on the scatter diagram and investigated before any calculation.
\end{definition}

\textbf{Two regression lines, not one.} There are in fact \emph{two} natural straight lines on the same scatter diagram:
\begin{itemize}
\item the \cle{regression line of $y$ on $x$}, used to predict $y$ from $x$; it minimises the \emph{vertical} deviations of the points from the line;
\item the \cle{regression line of $x$ on $y$}, used to predict $x$ from $y$; it minimises the \emph{horizontal} deviations.
\end{itemize}
Both lines pass through the mean point $(\bar{x}, \bar{y})$, but they generally have \emph{different gradients}. The closer the points are to a single straight line, the closer the two lines are to each other; if the correlation is perfect, the two lines coincide.

\begin{exoresolu}[Computing the mean point and sketching both lines]
For the ten students of the previous worked example, compute the mean point and describe how to sketch the two regression lines.
\tcblower
\textbf{Solution.} Summing the values:
\[
\sum x_i = 4+6+8+9+11+13+15+17+19+21 = 123, \qquad \bar{x} = \frac{123}{10} = 12.3,
\]
\[
\sum y_i = 38+45+50+55+58+63+70+72+80+84 = 615, \qquad \bar{y} = \frac{615}{10} = 61.5.
\]
The mean point is $(12.3,\ 61.5)$. To sketch the line of $y$ on $x$ by eye, draw through $(12.3, 61.5)$ the straight line that best follows the rising cloud of points, balancing vertical deviations. The line of $x$ on $y$ also passes through $(12.3, 61.5)$ but is slightly steeper on the diagram (it balances horizontal deviations). Since the association here is strong, the two lines lie close together, as shown below.
\end{exoresolu}

\begin{popfigure}
\begin{tikzpicture}[scale=0.42]
\draw[->, thick, popink] (0,0) -- (23,0) node[below left, popink] {\small $x$ (hours)};
\draw[->, thick, popink] (0,0) -- (0,11.5) node[below left, popink] {\small $y$ (score)};
\foreach \x in {0,4,8,12,16,20} \draw[popline] (\x,0.08) -- (\x,-0.08) node[below, popmuted] {\scriptsize \x};
\foreach \y/\lab in {1/40,3/50,5/60,7/70,9/80} \draw[popline] (0.08,\y) -- (-0.08,\y) node[left, popmuted] {\scriptsize \lab};
\foreach \x/\y in {4/0.8, 6/1.5, 8/2.0, 9/2.5, 11/2.8, 13/3.3, 15/4.0, 17/4.2, 19/5.0, 21/5.4}
  \fill[popblue] (\x,\y) circle (0.14);
\draw[poporange, very thick] (2,0.5) -- (22,5.6);
\draw[poppurple, very thick, dashed] (3.5,0.2) -- (20.5,5.9);
\fill[popdark] (12.3,3.075) circle (0.2);
\node[popdark, align=center] at (12.3,4.0) {\scriptsize $(\bar{x},\bar{y})=(12.3,\ 61.5)$};
\node[poporange, align=center] at (20,3.4) {\scriptsize $y$ on $x$};
\node[poppurple, align=center] at (6.5,2.6) {\scriptsize $x$ on $y$};
\end{tikzpicture}
\legende{Both regression lines sketched on the same scatter diagram; they cross at the mean point $(\bar{x}, \bar{y})$.}
\end{popfigure}

\begin{aretenir}[Key points of this section]
\begin{itemize}
\item Bivariate data are pairs $(x_i, y_i)$ measured on the same individuals; $x$ is the independent (explanatory) variable, $y$ the dependent (response) variable.
\item Regression estimates the \emph{average} value of one variable from the other; linear regression uses a straight line $y = a + bx$.
\item Always begin with a scatter diagram: $x$ on the horizontal axis, $y$ on the vertical axis; read its direction, form and strength.
\item Association does \textbf{not} imply causation --- beware of lurking variables.
\item An outlier lies far from the general pattern and can distort the regression line.
\item There are two regression lines ($y$ on $x$ and $x$ on $y$); both pass through the mean point $(\bar{x}, \bar{y})$ and coincide only when correlation is perfect.
\end{itemize}
\end{aretenir}

\begin{exercice}[Practice]
A farmer in Buea records, for 8 plots, the quantity of fertiliser applied $x$ (kg) and the harvest $y$ (bags of cocoyams):
\[
\begin{array}{c|cccccccc}
x & 10 & 15 & 20 & 25 & 30 & 35 & 40 & 45 \\ \hline
y & 12 & 15 & 14 & 19 & 22 & 21 & 26 & 28
\end{array}
\]
\begin{enumerate}
\item Identify the independent and dependent variables.
\item Draw a scatter diagram.
\item Compute the mean point $(\bar{x}, \bar{y})$ and mark it on your diagram.
\item Sketch by eye a regression line through the mean point and describe the association.
\item The farmer claims: ``Fertiliser alone causes the harvest.'' Comment critically.
\end{enumerate}
\end{exercice}

\begin{corrige}[Practice]
\begin{enumerate}
\item $x$ (fertiliser, kg) is independent; $y$ (harvest, bags) is dependent.
\item Plot the eight points with $x$ horizontal (0 to 50) and $y$ vertical (10 to 30); the points rise from bottom-left to top-right.
\item $\sum x_i = 220$, so $\bar{x} = 27.5$; $\sum y_i = 157$, so $\bar{y} = 19.625$. Mean point $(27.5,\ 19.625)$.
\item A line through $(27.5, 19.625)$ following the rising cloud shows a strong positive linear association.
\item The diagram shows association only. Other factors (rainfall, soil quality, seed variety, weeding) may also influence the harvest; without a controlled experiment, causation cannot be claimed.
\end{enumerate}
\end{corrige}

In the next section we replace the ``by eye'' line by a first genuine calculation method: \textbf{Mayer's method}, which splits the data into two groups and uses their mean points to determine the gradient and intercept of the regression line.

\coursec{Mayer's Method for the Regression Line}

In the previous section we saw that a scatter diagram often reveals a linear trend between two variables, and that a regression line is a straight line drawn to summarise that trend. But how do we actually \emph{find} such a line? Drawing it ``by eye'' is subjective: two students looking at the same cloud of points will rarely draw the same line. We need a simple, reproducible rule. The first such rule we study is \cle{Mayer's method}: an intuitive technique, due to the German astronomer Tobias Mayer, which replaces the whole cloud of points by just \emph{two} representative points and draws the line through them. It is quick, requires only averages, and gives a good first approximation of the regression line. Later in this chapter we shall meet the more precise \emph{least squares} method; Mayer's method is the natural stepping stone towards it.

\courssub{The Intuition Behind Mayer's Method}

Imagine a farmer near Bafoussam who records, for eight plots of land, the quantity of fertiliser applied (in kg) and the maize yield obtained (in bags). The scatter diagram shows yield increasing with fertiliser, roughly along a straight line. Mayer's idea is the following:

\begin{itemize}
\item A single point is a poor summary of the data, but the \emph{mean point} (centre of gravity) of a group of points is a good summary of that group.
\item If we split the data into a ``lower'' group (small $x$-values) and an ``upper'' group (large $x$-values), each group has its own mean point.
\item A sensible trend line should pass through, or very near, both centres of gravity. Since two points determine a unique straight line, we simply draw the line through the two mean points.
\end{itemize}

That is the whole method. Let us now state it rigorously.

\begin{definition}[Group mean points]
Let $(x_1,y_1),(x_2,y_2),\dots,(x_n,y_n)$ be a set of $n$ data points, ordered so that $x_1 \le x_2 \le \dots \le x_n$. Split the data into a \cle{lower group} of the first $n_1$ points and an \cle{upper group} of the remaining $n_2$ points (with $n_1+n_2=n$ and $n_1$, $n_2$ as equal as possible). The \textbf{group mean points} are
\[
G_1=(\bar{x}_1,\bar{y}_1) \quad\text{with}\quad \bar{x}_1=\frac{1}{n_1}\sum_{i=1}^{n_1}x_i,\qquad \bar{y}_1=\frac{1}{n_1}\sum_{i=1}^{n_1}y_i,
\]
\[
G_2=(\bar{x}_2,\bar{y}_2) \quad\text{with}\quad \bar{x}_2=\frac{1}{n_2}\sum_{i=n_1+1}^{n}x_i,\qquad \bar{y}_2=\frac{1}{n_2}\sum_{i=n_1+1}^{n}y_i.
\]
The \cle{Mayer line} is the straight line through $G_1$ and $G_2$.
\end{definition}

\begin{attention}[Order by $x$ before splitting!]
The split into two groups is \textbf{always} made after ordering the data by the $x$-values (the explanatory variable), \emph{never} in the order the data happen to be written in the table, and never by ordering the $y$-values. Splitting unordered data mixes large and small $x$-values in the same group, the two mean points come out close together, and the resulting line can have a completely wrong gradient. This is the single most common mistake with Mayer's method in the examination.
\end{attention}

\courssub{Deriving the Equation of the Mayer Line}

Once the two mean points $G_1=(\bar{x}_1,\bar{y}_1)$ and $G_2=(\bar{x}_2,\bar{y}_2)$ are known, finding the equation of the line through them is pure Lower Sixth coordinate geometry.

\begin{propriete}[Equation of the Mayer regression line]
The Mayer regression line of $y$ on $x$ has gradient
\[
b=\frac{\bar{y}_2-\bar{y}_1}{\bar{x}_2-\bar{x}_1}
\]
and equation
\[
y-\bar{y}_1=b\,(x-\bar{x}_1),
\qquad\text{equivalently}\qquad
y=bx+a\ \ \text{with}\ \ a=\bar{y}_1-b\bar{x}_1.
\]
(The derivation is a direct application of the two-point form of a straight line and \textbf{is examinable}.)
\end{propriete}

\textbf{Derivation.} A non-vertical line through $(\bar{x}_1,\bar{y}_1)$ has equation $y-\bar{y}_1=b(x-\bar{x}_1)$ for some gradient $b$. Requiring it to pass also through $(\bar{x}_2,\bar{y}_2)$ gives
\[
\bar{y}_2-\bar{y}_1=b(\bar{x}_2-\bar{x}_1)
\quad\Longrightarrow\quad
b=\frac{\bar{y}_2-\bar{y}_1}{\bar{x}_2-\bar{x}_1},
\]
which is well defined because $\bar{x}_2\neq\bar{x}_1$ (the upper group contains strictly larger $x$-values than the lower group). Substituting back gives the stated equation. \hfill$\blacksquare$

\begin{methode}[Calculating the regression line by Mayer's method]
\begin{enumerate}
\item \textbf{Order} the $n$ data points by increasing values of $x$.
\item \textbf{Split} into two groups as equal as possible: if $n$ is even, each group has $n/2$ points; if $n$ is odd, either exclude the middle point or place it in both groups (state your convention).
\item \textbf{Compute} the mean point $G_1=(\bar{x}_1,\bar{y}_1)$ of the lower group and $G_2=(\bar{x}_2,\bar{y}_2)$ of the upper group.
\item \textbf{Compute the gradient} $b=\dfrac{\bar{y}_2-\bar{y}_1}{\bar{x}_2-\bar{x}_1}$.
\item \textbf{Write the equation} $y-\bar{y}_1=b(x-\bar{x}_1)$ and simplify to the form $y=bx+a$.
\item \textbf{Check}: verify that $G_2$ satisfies your equation (it must, by construction).
\end{enumerate}
\end{methode}

\courssub{Handling an Odd Number of Points}

When $n$ is even, the split is unambiguous. When $n$ is odd, say $n=2m+1$, two conventions are accepted:

\begin{itemize}
\item \textbf{Exclude the middle point}: the lower group is the first $m$ points, the upper group the last $m$ points; the $(m+1)$-th point is ignored.
\item \textbf{Share the middle point}: the middle point is included in \emph{both} groups, so each group has $m+1$ points.
\end{itemize}

Both conventions appear in textbooks and both are accepted at the GCE, provided you \emph{state clearly} which one you use and apply it consistently. In this course, unless told otherwise, we \textbf{exclude the middle point} when $n$ is odd.

\courssub{Worked Example: Fertiliser and Maize Yield}

\begin{exoresolu}[Mayer's method on a Cameroonian data set]
An agricultural extension officer in the West Region records the quantity of fertiliser $x$ (kg) applied to eight maize plots and the yield $y$ (bags) harvested:
\begin{center}
\begin{tabular}{|l|c|c|c|c|c|c|c|c|}
\hline
\rowcolor{popblueL} Plot & A & B & C & D & E & F & G & H \\
\hline
$x$ (kg) & 10 & 20 & 30 & 40 & 50 & 60 & 70 & 80 \\
\hline
$y$ (bags) & 4 & 6 & 7 & 9 & 10 & 13 & 14 & 15 \\
\hline
\end{tabular}
\end{center}
\begin{enumerate}
\item Find the equation of the Mayer regression line of $y$ on $x$.
\item Use it to predict the yield for $x=55$ kg and for $x=120$ kg, commenting on reliability.
\end{enumerate}
\tcblower
\textbf{Solution.}

\textbf{1.} The data are already ordered by increasing $x$. Since $n=8$ is even, each group has $4$ points.

\emph{Lower group} ($x=10,20,30,40$):
\[
\bar{x}_1=\frac{10+20+30+40}{4}=\frac{100}{4}=25,\qquad
\bar{y}_1=\frac{4+6+7+9}{4}=\frac{26}{4}=6.5.
\]
So $G_1=(25,\,6.5)$.

\emph{Upper group} ($x=50,60,70,80$):
\[
\bar{x}_2=\frac{50+60+70+80}{4}=\frac{260}{4}=65,\qquad
\bar{y}_2=\frac{10+13+14+15}{4}=\frac{52}{4}=13.
\]
So $G_2=(65,\,13)$.

\emph{Gradient}:
\[
b=\frac{\bar{y}_2-\bar{y}_1}{\bar{x}_2-\bar{x}_1}=\frac{13-6.5}{65-25}=\frac{6.5}{40}=0.1625.
\]

\emph{Equation}: $y-6.5=0.1625(x-25)$, i.e.
\[
\boxed{\,y=0.1625x+2.4375\,}
\]
since $a=6.5-0.1625\times25=6.5-4.0625=2.4375$.

\emph{Check with }$G_2$: $0.1625\times65+2.4375=10.5625+2.4375=13=\bar{y}_2$. \checkmark

\textbf{2.} For $x=55$: $y=0.1625\times55+2.4375=8.9375+2.4375=11.375\approx 11.4$ bags. Since $55$ kg lies \emph{inside} the observed range $[10,80]$, this is \textbf{interpolation} and the prediction is reasonably reliable (the scatter is close to linear).

For $x=120$: $y=0.1625\times120+2.4375=19.5+2.4375=21.9375\approx 21.9$ bags. Since $120$ kg lies far \emph{outside} the observed range, this is \textbf{extrapolation}: we have no evidence the linear trend continues (in reality, excess fertiliser can even reduce yield), so this prediction is \emph{unreliable}.
\end{exoresolu}

\courssub{Illustrating the Method Graphically}

The figure below shows the fertiliser--maize data, the two group mean points $G_1$ and $G_2$, and the Mayer line passing through them.

\begin{popfigure}
\begin{center}
\begin{tikzpicture}
\begin{axis}[
    width=12cm, height=8cm,
    xlabel={Fertiliser $x$ (kg)}, ylabel={Yield $y$ (bags)},
    xmin=0, xmax=95, ymin=0, ymax=18,
    xtick={0,10,20,30,40,50,60,70,80,90},
    ytick={0,2,4,6,8,10,12,14,16,18},
    grid=both, grid style={popline, very thin},
    legend style={at={(0.02,0.98)}, anchor=north west, font=\small},
]
\addplot[only marks, mark=*, mark size=2.2pt, color=popblue]
    coordinates {(10,4) (20,6) (30,7) (40,9) (50,10) (60,13) (70,14) (80,15)};
\addlegendentry{Data points}
\addplot[domain=0:95, thick, color=poporange] {0.1625*x+2.4375};
\addlegendentry{Mayer line $y=0.1625x+2.4375$}
\addplot[only marks, mark=square*, mark size=3pt, color=poppurple]
    coordinates {(25,6.5) (65,13)};
\addlegendentry{Group mean points}
\node[color=poppurple, font=\small, anchor=south west] at (axis cs:25,6.5) {$G_1(25;\,6.5)$};
\node[color=poppurple, font=\small, anchor=south east] at (axis cs:65,13) {$G_2(65;\,13)$};
\draw[dashed, color=popmuted] (axis cs:40,0) -- (axis cs:40,18);
\node[color=popmuted, font=\footnotesize, rotate=90, anchor=south] at (axis cs:40,17) {group split};
\end{axis}
\end{tikzpicture}
\end{center}
\legende{Scatter diagram of the fertiliser--maize data. The dashed line separates the lower group (left) from the upper group (right); the Mayer line passes through the two group mean points $G_1$ and $G_2$.}
\end{popfigure}

Notice how the line balances the points: it does not pass through every point, but it passes through the \emph{centres of gravity} of the two halves of the data, which is exactly the spirit of a trend line.

\courssub{Predictions and Reliability}

The Mayer line, like any regression line, is used to \textbf{predict} $y$ for a given $x$:

\begin{itemize}
\item If the given $x$ lies \emph{within} the range of observed $x$-values, the prediction is an \cle{interpolation} and is generally reliable, provided the points are reasonably close to a straight line.
\item If the given $x$ lies \emph{outside} the observed range, the prediction is an \cle{extrapolation} and must be treated with great caution: the linear relationship may not hold beyond the data.
\end{itemize}

Reliability also depends on the \emph{strength} of the linear association, which we shall measure later with the product moment correlation coefficient.

\begin{attention}[Mayer's line is only an approximation]
Mayer's method uses only \emph{two} summary points; the individual positions of all the other points play no role except through the two group means. Consequently:
\begin{itemize}
\item Different splits (odd-$n$ conventions) give slightly different lines.
\item The Mayer line is \textbf{not} the least squares regression line studied later in this chapter, and it does not minimise any sum of squared errors.
\item In the examination, do not mix the two methods: if the question says ``Mayer's method'', use group mean points; if it says ``least squares'', use the formulae $b=\dfrac{s_{xy}}{s_x^2}$ developed in the next sections.
\end{itemize}
Mayer's method should be seen as a quick, transparent first estimate --- excellent for checking the plausibility of a least squares result.
\end{attention}

\begin{aretenir}[Key points]
\begin{itemize}
\item Order the data by $x$, split into two (near-)equal groups, and compute the mean points $G_1=(\bar{x}_1,\bar{y}_1)$ and $G_2=(\bar{x}_2,\bar{y}_2)$.
\item The Mayer line passes through $G_1$ and $G_2$: gradient $b=\dfrac{\bar{y}_2-\bar{y}_1}{\bar{x}_2-\bar{x}_1}$, equation $y-\bar{y}_1=b(x-\bar{x}_1)$.
\item For odd $n$, exclude the middle point (or share it between groups) and state your convention.
\item Interpolation is reliable; extrapolation is risky.
\item Mayer's method is simple and intuitive but only approximate; the least squares method gives the ``best'' line in a precise sense.
\end{itemize}
\end{aretenir}

\begin{exercice}[Mayer's method with an odd number of points]
The table gives the number of years of experience $x$ of nine tailors in a workshop in Buea and the number $y$ of shirts each completes per day.
\begin{center}
\begin{tabular}{|l|c|c|c|c|c|c|c|c|c|}
\hline
\rowcolor{popblueL} $x$ (years) & 1 & 2 & 3 & 4 & 5 & 6 & 7 & 8 & 9 \\
\hline
\rowcolor{popblueL} $y$ (shirts) & 3 & 4 & 5 & 5 & 6 & 7 & 8 & 8 & 9 \\
\hline
\end{tabular}
\end{center}
\begin{enumerate}
\item Using Mayer's method and \emph{excluding} the middle point, find the equation of the regression line of $y$ on $x$.
\item Repeat, this time \emph{sharing} the middle point between the two groups. Comment on the difference.
\item Predict the daily output of a tailor with 6.5 years of experience.
\end{enumerate}
\end{exercice}

\begin{corrige}[Exercise 1]
\textbf{1.} Excluding the middle point $(5,6)$: lower group $x=1,2,3,4$; upper group $x=6,7,8,9$.
\[
\bar{x}_1=\frac{1+2+3+4}{4}=2.5,\quad \bar{y}_1=\frac{3+4+5+5}{4}=4.25;\qquad
\bar{x}_2=\frac{6+7+8+9}{4}=7.5,\quad \bar{y}_2=\frac{7+8+8+9}{4}=8.
\]
\[
b=\frac{8-4.25}{7.5-2.5}=\frac{3.75}{5}=0.75,\qquad
y-4.25=0.75(x-2.5)\ \Longrightarrow\ y=0.75x+2.375.
\]
\textbf{2.} Sharing the middle point: each group has 5 points. Lower: $\bar{x}_1=\frac{15}{5}=3$, $\bar{y}_1=\frac{23}{5}=4.6$. Upper: $\bar{x}_2=\frac{35}{5}=7$, $\bar{y}_2=\frac{38}{5}=7.6$.
\[
b=\frac{7.6-4.6}{7-3}=\frac{3}{4}=0.75,\qquad y-4.6=0.75(x-3)\ \Longrightarrow\ y=0.75x+2.35.
\]
The gradient is identical ($0.75$); only the intercept changes slightly ($2.375$ vs $2.35$). The two conventions give very close lines here because the data are almost perfectly linear.
\textbf{3.} $y=0.75\times6.5+2.375=7.25$ shirts (interpolation, reliable). With the second line: $7.225$ --- essentially the same prediction.
\end{corrige}

\begin{savaistu}[Tobias Mayer, astronomer of the Moon]
Tobias Mayer (1723--1762) was a German astronomer famous for his remarkably accurate tables of the Moon's motion. Faced with many imperfect observations, he had the idea of grouping them and averaging within groups to reduce observational error --- the ancestor of the method you have just studied. His averaging philosophy later inspired Legendre and Gauss, who developed the method of \emph{least squares} that we shall meet in the next sections. In Cameroon today, the same averaging spirit is used by agronomists of IRAD (Institute of Agricultural Research for Development) when summarising field trial data before fitting more sophisticated models.
\end{savaistu}

\coursec{Correlation, Sample Variance and Covariance}

In the previous section we met Mayer's method, a quick way of fitting a straight line to a scatter diagram by splitting the data into two groups. Mayer's method gives us \emph{a} line, but it tells us nothing about how well the points actually cluster around a line. Two data sets can produce almost identical Mayer lines, yet in one the points hug the line tightly while in the other they are scattered all over the plane. We therefore need a way to \emph{measure} the strength and direction of a linear relationship. That is the purpose of this section: we introduce the idea of \cle{correlation}, and we build the two fundamental tools on which everything else in this chapter rests — the \cle{sample variance} and the \cle{sample covariance}.

\courssub{Introduction to Correlation}

\begin{experience}[Observing how two variables move together]
A teacher in Buea records, for 10 students, the number of hours $x$ spent revising Mathematics each week and the mark $y$ (out of 100) obtained in the mock examination. Plotting the points $(x,y)$, she notices that the points rise from the bottom-left to the top-right of the diagram: students who revise more tend to score higher. A second teacher records the price $x$ of a second-hand car and its age $y$; the points fall from top-left to bottom-right. A third records shoe size $x$ and History mark $y$; the points form a shapeless cloud.
\end{experience}

These three situations illustrate the three basic types of correlation.

\begin{definition}[Correlation]
\cle{Correlation} describes the extent to which two variables $x$ and $y$ are \emph{linearly} related.
\begin{itemize}
\item There is \textbf{positive correlation} when large values of $x$ tend to occur with large values of $y$ (the points rise from left to right).
\item There is \textbf{negative correlation} when large values of $x$ tend to occur with small values of $y$ (the points fall from left to right).
\item There is \textbf{zero (or no) correlation} when there is no apparent linear tendency at all.
\end{itemize}
Correlation is said to be \textbf{strong} when the points lie close to a straight line, and \textbf{weak} when they are widely scattered around it.
\end{definition}

\begin{popfigure}
\begin{center}
\begin{tikzpicture}
\begin{axis}[width=0.3\linewidth,height=0.25\linewidth,title={Strong positive},xlabel={$x$},ylabel={$y$},xmin=0,xmax=10,ymin=0,ymax=10,xtick=\empty,ytick=\empty,axis lines=left,scale only axis]
\addplot[only marks,mark=*,mark size=1.5pt,color=popblue] coordinates {(1,1.5)(2,2.4)(3,3.2)(4,3.9)(5,5.1)(6,5.8)(7,7.1)(8,7.8)(9,8.9)};
\end{axis}
\end{tikzpicture}\hfill
\begin{tikzpicture}
\begin{axis}[width=0.3\linewidth,height=0.25\linewidth,title={Weak positive},xlabel={$x$},ylabel={$y$},xmin=0,xmax=10,ymin=0,ymax=10,xtick=\empty,ytick=\empty,axis lines=left,scale only axis]
\addplot[only marks,mark=*,mark size=1.5pt,color=popgreen] coordinates {(1,3)(2,1.5)(3,5)(4,3.5)(5,6.5)(6,4)(7,8)(8,6)(9,8.5)};
\end{axis}
\end{tikzpicture}\hfill
\begin{tikzpicture}
\begin{axis}[width=0.3\linewidth,height=0.25\linewidth,title={Negative},xlabel={$x$},ylabel={$y$},xmin=0,xmax=10,ymin=0,ymax=10,xtick=\empty,ytick=\empty,axis lines=left,scale only axis]
\addplot[only marks,mark=*,mark size=1.5pt,color=poporange] coordinates {(1,8.8)(2,8)(3,6.9)(4,6.2)(5,5)(6,4.4)(7,3.1)(8,2.2)(9,1.3)};
\end{axis}
\end{tikzpicture}
\end{center}
\legende{Scatter diagrams showing strong positive, weak positive and negative correlation.}
\end{popfigure}

\begin{attention}[Correlation is not causation]
A strong correlation between two variables does \emph{not} prove that one causes the other. Ice-cream sales and drowning incidents are positively correlated, but neither causes the other: both increase in hot weather. Always ask whether a hidden third variable could explain the observed association.
\end{attention}

\courssub{Sample Variance}

To measure correlation numerically we first need to measure how \emph{spread out} each variable is on its own. You already know the mean $\bar{x}=\dfrac{1}{n}\sum x$. The spread about the mean is measured by the variance.

\begin{definition}[Sample variance]
For a sample of $n$ values $x_1,x_2,\dots,x_n$ with mean $\bar{x}$, the \cle{sample variance} of $x$ is
\[ s_{xx}=\frac{1}{n}\sum_{i=1}^{n}(x_i-\bar{x})^2=\frac{1}{n}\sum_{i=1}^{n}x_i^2-\bar{x}^2. \]
Similarly, for values $y_1,\dots,y_n$ with mean $\bar{y}$,
\[ s_{yy}=\frac{1}{n}\sum_{i=1}^{n}(y_i-\bar{y})^2=\frac{1}{n}\sum_{i=1}^{n}y_i^2-\bar{y}^2. \]
The positive square roots $s_x=\sqrt{s_{xx}}$ and $s_y=\sqrt{s_{yy}}$ are the \textbf{sample standard deviations}.
\end{definition}

The second form of each formula is the \emph{computational} form; it is far quicker with a calculator. Let us prove the equivalence for $s_{xx}$ (the proof for $s_{yy}$ is identical).

\begin{propriete}[Equivalence of the two forms of the variance]
$\dfrac{1}{n}\sum (x_i-\bar{x})^2=\dfrac{1}{n}\sum x_i^2-\bar{x}^2$.
\end{propriete}

\textbf{Proof.} Expanding the square and using $\sum x_i=n\bar{x}$:
\begin{align*}
\sum_{i=1}^{n}(x_i-\bar{x})^2 &=\sum_{i=1}^{n}\left(x_i^2-2\bar{x}x_i+\bar{x}^2\right)\\
&=\sum x_i^2-2\bar{x}\sum x_i+n\bar{x}^2\\
&=\sum x_i^2-2\bar{x}(n\bar{x})+n\bar{x}^2\\
&=\sum x_i^2-n\bar{x}^2.
\end{align*}
Dividing by $n$ gives $\dfrac{1}{n}\sum(x_i-\bar{x})^2=\dfrac{1}{n}\sum x_i^2-\bar{x}^2$. $\blacksquare$

\begin{attention}[Do not confuse $\sum x^2$ with $(\sum x)^2$]
This is the single most common error in this chapter. $\sum x^2$ means: \emph{square each value, then add}. $(\sum x)^2$ means: \emph{add the values, then square the total}. For the data $1,2,3$: $\sum x^2=1+4+9=14$ but $(\sum x)^2=6^2=36$. Keep the two quantities in clearly separate columns of your table.
\end{attention}

\courssub{Sample Covariance}

Variance measures the spread of \emph{one} variable. To describe how two variables move \emph{together}, we need a mixed quantity.

\begin{experience}[The idea behind covariance]
For each data point, compute the deviations $x-\bar{x}$ and $y-\bar{y}$. If $x$ and $y$ tend to be above their means \emph{together} (and below together), the products $(x-\bar{x})(y-\bar{y})$ are mostly positive, so their average is positive. If one variable tends to be above its mean when the other is below, the products are mostly negative. If there is no pattern, positive and negative products cancel out. The average of these products is the covariance.
\end{experience}

\begin{definition}[Sample covariance]
For $n$ pairs of data $(x_i,y_i)$ with means $\bar{x}$ and $\bar{y}$, the \cle{sample covariance} of $x$ and $y$ is
\[ s_{xy}=\frac{1}{n}\sum_{i=1}^{n}(x_i-\bar{x})(y_i-\bar{y})=\frac{1}{n}\sum_{i=1}^{n}x_iy_i-\bar{x}\,\bar{y}. \]
\end{definition}

\begin{propriete}[Equivalence of the two forms of the covariance]
$\dfrac{1}{n}\sum (x_i-\bar{x})(y_i-\bar{y})=\dfrac{1}{n}\sum x_iy_i-\bar{x}\bar{y}$.
\end{propriete}

\textbf{Proof.} Expand the product and use $\sum x_i=n\bar{x}$, $\sum y_i=n\bar{y}$:
\begin{align*}
\sum_{i=1}^{n}(x_i-\bar{x})(y_i-\bar{y})&=\sum_{i=1}^{n}\left(x_iy_i-\bar{x}y_i-\bar{y}x_i+\bar{x}\bar{y}\right)\\
&=\sum x_iy_i-\bar{x}\sum y_i-\bar{y}\sum x_i+n\bar{x}\bar{y}\\
&=\sum x_iy_i-\bar{x}(n\bar{y})-\bar{y}(n\bar{x})+n\bar{x}\bar{y}\\
&=\sum x_iy_i-n\bar{x}\bar{y}.
\end{align*}
Dividing by $n$ gives the result. $\blacksquare$

\begin{aretenir}[Interpreting the sign of the covariance]
\begin{itemize}
\item $s_{xy}>0$: positive linear association (points rise left to right).
\item $s_{xy}<0$: negative linear association (points fall left to right).
\item $s_{xy}\approx 0$: little or no linear association.
\end{itemize}
Note also that $s_{xx}$ is simply the covariance of $x$ with itself, and covariance is symmetric: $s_{xy}=s_{yx}$.
\end{aretenir}

\begin{attention}[Covariance depends on the units of measurement]
The covariance is \emph{not} a standardised measure of strength. If $x$ is height measured in centimetres and $y$ is mass in kilograms, $s_{xy}$ is measured in $\mathrm{cm}\cdot\mathrm{kg}$. Measuring the same heights in metres multiplies the covariance by $\tfrac{1}{100}$, without the underlying relationship changing at all! Consequently a covariance of $50$ is not automatically ``stronger'' than a covariance of $5$ — the units must be taken into account. This is exactly why, in a later section, we shall divide by the standard deviations to obtain the unit-free \emph{product moment correlation coefficient}. For now, use the covariance mainly for its \textbf{sign} and as a building block.
\end{attention}

\courssub{Setting Out the Calculations in a Table}

\begin{methode}[Computing $s_{xx}$, $s_{yy}$ and $s_{xy}$ from a table of values]
\begin{enumerate}
\item Draw a table with columns $x$, $y$, $x^2$, $y^2$, $xy$ and complete one row per data pair.
\item Add a final row of totals: $\sum x$, $\sum y$, $\sum x^2$, $\sum y^2$, $\sum xy$.
\item Compute the means: $\bar{x}=\dfrac{\sum x}{n}$ and $\bar{y}=\dfrac{\sum y}{n}$.
\item Apply the computational formulae:
\[ s_{xx}=\frac{\sum x^2}{n}-\bar{x}^2,\qquad s_{yy}=\frac{\sum y^2}{n}-\bar{y}^2,\qquad s_{xy}=\frac{\sum xy}{n}-\bar{x}\bar{y}. \]
\item Interpret the sign of $s_{xy}$ and compare the magnitudes with the spread of the data.
\end{enumerate}
\end{methode}

\begin{exoresolu}[Revision time and mock examination marks]
Six students recorded their weekly revision time $x$ (hours) and their mock mark $y$ (out of 20):
\begin{center}
\begin{tabular}{|c|cccccc|}
\hline
\rowcolor{popblueL} $x$ & 2 & 4 & 5 & 7 & 8 & 10 \\
\hline
\rowcolor{popblueL} $y$ & 6 & 8 & 9 & 12 & 13 & 15 \\
\hline
\end{tabular}
\end{center}
Calculate $s_{xx}$, $s_{yy}$ and $s_{xy}$, and comment on the relationship.
\tcblower
\textbf{Step 1: the table of values.}
\begin{center}
\begin{tabular}{|c|c|c|c|c|}
\hline
\rowcolor{popblueL} $x$ & $y$ & $x^2$ & $y^2$ & $xy$ \\
\hline
2 & 6 & 4 & 36 & 12 \\
4 & 8 & 16 & 64 & 32 \\
5 & 9 & 25 & 81 & 45 \\
7 & 12 & 49 & 144 & 84 \\
8 & 13 & 64 & 169 & 104 \\
10 & 15 & 100 & 225 & 150 \\
\hline
\rowcolor{popblueL} $\sum x=36$ & $\sum y=63$ & $\sum x^2=258$ & $\sum y^2=719$ & $\sum xy=427$ \\
\hline
\end{tabular}
\end{center}
\textbf{Step 2: the means.} With $n=6$:
\[ \bar{x}=\frac{36}{6}=6,\qquad \bar{y}=\frac{63}{6}=10.5. \]
\textbf{Step 3: the variances and covariance.}
\begin{align*}
s_{xx}&=\frac{258}{6}-6^2=43-36=7,\\[2pt]
s_{yy}&=\frac{719}{6}-10.5^2=119.833\ldots-110.25=\frac{115}{12}\approx 9.583,\\[2pt]
s_{xy}&=\frac{427}{6}-(6)(10.5)=71.166\ldots-63=\frac{49}{6}\approx 8.167.
\end{align*}
\textbf{Step 4: interpretation.} Since $s_{xy}>0$, there is a positive linear association: more revision time tends to go with a higher mark. Note that $s_{xy}$ is measured in ``hour-marks'', so its size alone cannot tell us how strong the link is — we shall standardise it later. As a check, $s_{xy}^2=\left(\tfrac{49}{6}\right)^2\approx 66.7$ is less than $s_{xx}s_{yy}=7\times 9.583\approx 67.1$, which must always hold.
\end{exoresolu}

\begin{propriete}[The Cauchy--Schwarz inequality for data]
For any bivariate data set, $s_{xy}^2\le s_{xx}s_{yy}$, with equality exactly when all the points lie on a single straight line. This guarantees that the correlation coefficient we shall define later always lies between $-1$ and $1$.
\end{propriete}

\begin{exercice}[Practice]
Eight candidates in a Further Mathematics mock obtained the following marks in Paper 1 ($x$) and Paper 2 ($y$):
\begin{center}
\begin{tabular}{|c|cccccccc|}
\hline
\rowcolor{popblueL} $x$ & 12 & 15 & 18 & 20 & 24 & 27 & 30 & 34 \\
\hline
\rowcolor{popblueL} $y$ & 10 & 14 & 15 & 19 & 20 & 25 & 26 & 31 \\
\hline
\end{tabular}
\end{center}
\begin{enumerate}
\item Set out a full table of values and find $\sum x$, $\sum y$, $\sum x^2$, $\sum y^2$, $\sum xy$.
\item Calculate $s_{xx}$, $s_{yy}$ and $s_{xy}$.
\item Verify that $s_{xy}^2\le s_{xx}s_{yy}$ and state what the sign of $s_{xy}$ tells you.
\end{enumerate}
\end{exercice}

\begin{corrige}[Exercise 1]
\textbf{1.} The completed table gives
\[ \sum x=180,\quad \sum y=160,\quad \sum x^2=4454,\quad \sum y^2=3624,\quad \sum xy=3969. \]
\textbf{2.} With $n=8$: $\bar{x}=\tfrac{180}{8}=22.5$ and $\bar{y}=\tfrac{160}{8}=20$. Hence
\[ s_{xx}=\frac{4454}{8}-22.5^2=556.75-506.25=50.5, \]
\[ s_{yy}=\frac{3624}{8}-20^2=453-400=53, \]
\[ s_{xy}=\frac{3969}{8}-(22.5)(20)=496.125-450=46.125. \]
\textbf{3.} $s_{xy}^2\approx 2127.5$ and $s_{xx}s_{yy}=50.5\times 53=2676.5$, so indeed $s_{xy}^2\le s_{xx}s_{yy}$. Since $s_{xy}>0$, performance in the two papers is positively associated: candidates who do well in Paper 1 tend to do well in Paper 2.
\end{corrige}

\begin{aretenir}[Key points of this section]
\begin{itemize}
\item Correlation describes the strength and direction of a \emph{linear} relationship: positive, negative or zero; strong or weak.
\item Sample variances: $s_{xx}=\dfrac{\sum x^2}{n}-\bar{x}^2$, $s_{yy}=\dfrac{\sum y^2}{n}-\bar{y}^2$.
\item Sample covariance: $s_{xy}=\dfrac{\sum xy}{n}-\bar{x}\bar{y}$; its \emph{sign} gives the direction of the association.
\item Always work from a table with columns $x$, $y$, $x^2$, $y^2$, $xy$ and a row of totals.
\item Covariance carries units, so its magnitude alone cannot measure strength — this motivates the standardised correlation coefficient of the next sections.
\end{itemize}
\end{aretenir}

\coursec{Least Squares Regression Lines}

In the previous section, we introduced the sample variance $s_{xx}$, $s_{yy}$ and the sample covariance $s_{xy}$ as measures of spread and linear association. We also defined the product moment correlation coefficient $r = \dfrac{s_{xy}}{\sqrt{s_{xx}s_{yy}}}$ which quantifies the strength of a linear relationship. We now put these quantities to work: given a scatter diagram suggesting a linear trend, how do we find the \emph{best} straight line to summarise it? The answer is the \textbf{least squares regression line}. There are two such lines, depending on which variable we wish to predict.

\courssub{The Least Squares Principle}

Consider a set of $n$ data pairs $(x_i, y_i)$. If we draw a candidate line $y = a + bx$, the vertical distance from the point $(x_i, y_i)$ to the line is the \textbf{residual}
\[
e_i = y_i - (a + bx_i).
\]
Some residuals are positive (point above the line), some negative (point below). A good line should make these residuals ``small overall''. The \textbf{least squares criterion} chooses $a$ and $b$ to minimise the \textbf{Residual Sum of Squares (RSS)}:
\[
\mathrm{RSS}_{y|x} = \sum_{i=1}^n e_i^2 = \sum_{i=1}^n \bigl[y_i - (a + bx_i)\bigr]^2.
\]

\begin{popfigure}
\begin{tikzpicture}[scale=0.7]
% Axes
\draw[->] (0,0) -- (6.5,0) node[right] {$x$};
\draw[->] (0,0) -- (0,7) node[above] {$y$};
% Ticks
\foreach \x in {1,2,3,4,5} \draw (\x,0.1) -- (\x,-0.1) node[below] {\x};
\foreach \y in {1,2,3,4,5,6} \draw (0.1,\y) -- (-0.1,\y) node[left] {\y};
% Data points
\foreach \x/\y in {1/2, 2/3, 3/5, 4/4, 5/6} {
  \fill[popblue] (\x,\y) circle (3pt);
}
% Regression line y = 1.3 + 0.9x
\draw[poporange, thick] (0,1.3) -- (5.5, 1.3+0.9*5.5) node[right] {$y = a+bx$};
% Vertical residuals (dashed)
\foreach \x/\y in {1/2, 2/3, 3/5, 4/4, 5/6} {
  \pgfmathsetmacro{\ypred}{1.3 + 0.9*\x}
  \draw[popgreen, dashed] (\x,\y) -- (\x,\ypred);
  \fill[popgreen] (\x,\ypred) circle (2pt);
}
% Mean point
\fill[popred] (3,4) circle (4pt) node[above right] {$(\bar{x},\bar{y})$};
\end{tikzpicture}
\legende{Scatter diagram with vertical residuals (green dashed) for the least squares line of $y$ on $x$. The line passes through the mean point $(\bar{x},\bar{y})$.}
\end{popfigure}

\begin{definition}[Least Squares Regression Line of $y$ on $x$]
The line $y = a + bx$ that minimises the sum of squared \emph{vertical} deviations $\sum (y_i - a - bx_i)^2$ is called the \textbf{least squares regression line of $y$ on $x$}. Its coefficients are
\[
b = \frac{s_{xy}}{s_{xx}}, \qquad a = \bar{y} - b\bar{x},
\]
where $s_{xx} = \sum(x_i-\bar{x})^2$, $s_{xy} = \sum(x_i-\bar{x})(y_i-\bar{y})$, and $(\bar{x},\bar{y})$ is the mean point.
\end{definition}

\begin{propriete}[The Mean Point Property]
The least squares regression line of $y$ on $x$ \textbf{always passes through the mean point} $(\bar{x}, \bar{y})$.
\end{propriete}

\begin{proof}
Substitute $x = \bar{x}$ into the equation $y = a + bx$ with $a = \bar{y} - b\bar{x}$:
\[
y = (\bar{y} - b\bar{x}) + b\bar{x} = \bar{y}.
\]
Hence $(\bar{x}, \bar{y})$ lies on the line. This is a direct consequence of the normal equations derived from minimising RSS.
\end{proof}

\begin{propriete}[Residual Sum of Squares for $y$ on $x$]
The minimum value of the residual sum of squares (RSS) for the regression line of $y$ on $x$ is
\[
\mathrm{RSS}_{y|x} = s_{yy} - \frac{s_{xy}^2}{s_{xx}} = s_{yy}(1 - r^2).
\]
\end{propriete}

\begin{proof}
Using the normal equations $\sum y_i = na + b\sum x_i$ and $\sum x_iy_i = a\sum x_i + b\sum x_i^2$, one can show
\[
\sum (y_i - a - bx_i)^2 = \sum y_i^2 - a\sum y_i - b\sum x_iy_i = s_{yy} - b s_{xy} = s_{yy} - \frac{s_{xy}^2}{s_{xx}}.
\]
Since $r^2 = s_{xy}^2/(s_{xx}s_{yy})$, the result follows.
\end{proof}

\courssub{Regression Line of $x$ on $y$}

If our goal is to predict $x$ from a known value of $y$, we should minimise the sum of squared \emph{horizontal} deviations. This leads to the \textbf{regression line of $x$ on $y$}.

\begin{definition}[Least Squares Regression Line of $x$ on $y$]
The line $x = a' + b'y$ that minimises the sum of squared \emph{horizontal} deviations $\sum (x_i - a' - b'y_i)^2$ is called the \textbf{least squares regression line of $x$ on $y$}. Its coefficients are
\[
b' = \frac{s_{xy}}{s_{yy}}, \qquad a' = \bar{x} - b'\bar{y}.
\]
\end{definition}

\begin{propriete}[Both Lines Pass Through $(\bar{x},\bar{y})$]
Both the regression line of $y$ on $x$ and the regression line of $x$ on $y$ pass through the mean point $(\bar{x}, \bar{y})$.
\end{propriete}

\begin{proof}
For $x = a' + b'y$, substitute $y = \bar{y}$:
\[
x = (\bar{x} - b'\bar{y}) + b'\bar{y} = \bar{x}.
\]
Thus $(\bar{x}, \bar{y})$ satisfies both equations.
\end{proof}

\begin{propriete}[Residual Sum of Squares for $x$ on $y$]
The minimum value of the residual sum of squares (RSS) for the regression line of $x$ on $y$ is
\[
\mathrm{RSS}_{x|y} = s_{xx} - \frac{s_{xy}^2}{s_{yy}} = s_{xx}(1 - r^2).
\]
\end{propriete}

\courssub{Choosing the Appropriate Line}

\begin{methode}[When to Use Which Regression Line]
\begin{enumerate}
\item Identify the \textbf{independent} (explanatory) variable and the \textbf{dependent} (response) variable.
\item If $x$ is independent and we want to predict $y$, use the \textbf{regression line of $y$ on $x$}: $y = a + bx$.
\item If $y$ is independent and we want to predict $x$, use the \textbf{regression line of $x$ on $y$}: $x = a' + b'y$.
\item \textbf{Never} use the $y$ on $x$ line to predict $x$ from $y$ (or vice versa) by simply rearranging the equation. The two lines are different unless $r = \pm 1$ (perfect correlation).
\end{enumerate}
\end{methode}

\begin{attention}[Wrong Line for Prediction]
A very common error in examinations is to compute the line $y = a + bx$ and then, when asked to estimate $x$ for a given $y$, rearrange it to $x = (y-a)/b$. This is \textbf{incorrect}. You must compute the separate line $x = a' + b'y$. The two lines coincide only when all points lie exactly on a straight line ($r = \pm 1$).
\end{attention}

\courssub{Geometric Interpretation: The Two Lines and the Correlation}

Since both lines pass through $(\bar{x}, \bar{y})$, they intersect at that point. Their slopes are
\[
\text{slope of } y \text{ on } x: \quad b = \frac{s_{xy}}{s_{xx}}, \qquad
\text{slope of } x \text{ on } y: \quad \frac{1}{b'} = \frac{s_{yy}}{s_{xy}}.
\]
The product of the slopes is
\[
b \cdot \frac{1}{b'} = \frac{s_{xy}^2}{s_{xx}s_{yy}} = r^2.
\]
Thus the angle between the two lines is directly related to $|r|$:
\begin{itemize}
\item If $|r| = 1$ (perfect correlation), the two lines coincide ($b = 1/b'$).
\item If $r = 0$ (no linear correlation), the lines are perpendicular: $y$ on $x$ is horizontal ($b=0$), $x$ on $y$ is vertical ($b'=0$).
\item For $0 < |r| < 1$, the lines are distinct; the smaller $|r|$, the larger the angle between them.
\end{itemize}

\begin{popfigure}
\begin{tikzpicture}[scale=0.7]
% Axes
\draw[->] (0,0) -- (6.5,0) node[right] {$x$};
\draw[->] (0,0) -- (0,7) node[above] {$y$};
% Ticks
\foreach \x in {1,2,3,4,5} \draw (\x,0.1) -- (\x,-0.1) node[below] {\x};
\foreach \y in {1,2,3,4,5,6} \draw (0.1,\y) -- (-0.1,\y) node[left] {\y};
% Data points
\foreach \x/\y in {1/2, 2/3, 3/5, 4/4, 5/6} {
  \fill[popblue] (\x,\y) circle (3pt);
}
% Line y on x: y = 1.3 + 0.9x  -> at x=5.5, y=6.25
\draw[poporange, thick] (0,1.3) -- (5.5, 6.25) node[right] {$y$ on $x$};
% Line x on y: x = -0.6 + 0.9y  =>  y = (x+0.6)/0.9 -> at x=5.5, y=6.777... ; at x=0, y=0.667
\draw[poppurple, thick] (0,0.667) -- (5.5, 6.778) node[above right] {$x$ on $y$};
% Mean point
\fill[popred] (3,4) circle (5pt) node[above left] {$(\bar{x},\bar{y})$};
% Angle arc
\draw[popmuted, thick] (3,4) ++(0.8,0) arc (0:12:0.8) node[midway, right] {$\theta$};
\end{tikzpicture}
\legende{Both least squares regression lines intersect at the mean point $(\bar{x},\bar{y})$. The angle $\theta$ between them reflects the correlation strength: here $r=0.9$, so the lines are close.}
\end{popfigure}

\courssub{Computing the Lines from Summary Statistics}

In examinations, you are often given summary statistics $\sum x$, $\sum y$, $\sum x^2$, $\sum y^2$, $\sum xy$ and $n$. The following method is efficient and avoids rounding errors until the final step.

\begin{methode}[Finding Regression Coefficients from Summary Statistics]
Given $n$, $\sum x$, $\sum y$, $\sum x^2$, $\sum y^2$, $\sum xy$:
\begin{enumerate}
\item Compute the corrected sums of squares and products:
\[
s_{xx} = \sum x^2 - \frac{(\sum x)^2}{n}, \quad
s_{yy} = \sum y^2 - \frac{(\sum y)^2}{n}, \quad
s_{xy} = \sum xy - \frac{\sum x \sum y}{n}.
\]
\item Compute the means: $\bar{x} = \frac{\sum x}{n}$, $\bar{y} = \frac{\sum y}{n}$.
\item For $y$ on $x$: $b = \dfrac{s_{xy}}{s_{xx}}$, $a = \bar{y} - b\bar{x}$.
\item For $x$ on $y$: $b' = \dfrac{s_{xy}}{s_{yy}}$, $a' = \bar{x} - b'\bar{y}$.
\item Write the equations clearly: $y = a + bx$ and $x = a' + b'y$.
\end{enumerate}
\end{methode}

\courssub{Interpolation and Extrapolation}

Once a regression line is established, we use it to estimate values.
\begin{itemize}
\item \textbf{Interpolation}: Estimating $y$ for an $x$ value \emph{within} the range of the observed $x$ data. This is generally reliable if the linear model fits well.
\item \textbf{Extrapolation}: Estimating $y$ for an $x$ value \emph{outside} the range of the observed data. This is \textbf{dangerous} because the linear relationship may not hold beyond the observed interval.
\end{itemize}

\begin{attention}[Dangers of Extrapolation]
Extrapolation assumes the linear trend continues indefinitely. In reality, relationships often curve, saturate, or reverse outside the observed range. For example, a regression line fitted to the height of a child from age 5 to 15 would absurdly predict the child growing to 5 metres tall by age 50. \textbf{Always state the range of $x$ values used to fit the line and warn against predictions far outside this range.}
\end{attention}

\courssub{Worked Example: Computing Both Regression Lines and RSS}

\begin{exoresolu}[Finding both regression lines and RSS from raw data]
The table below shows the number of hours of revision ($x$) and the examination mark ($y$) for 6 students.

\begin{center}
\begin{tabular}{|c|cccccc|}\hline
$x$ (hours) & 10 & 12 & 15 & 18 & 20 & 25 \\ \hline
$y$ (mark)  & 45 & 52 & 60 & 68 & 72 & 80 \\ \hline
\end{tabular}
\end{center}

Find (a) the equation of the regression line of $y$ on $x$, (b) the equation of the regression line of $x$ on $y$, (c) the residual sum of squares for each line, (d) estimate the mark for a student who revises for 16 hours, (e) estimate the revision hours for a student who scored 65.
\tcblower
\textbf{Solution.}

First, compute the summary statistics. We add columns for $x^2$, $y^2$, $xy$.

\begin{center}
\begin{tabular}{|c|c|c|c|c|c|}\hline
$x$ & $y$ & $x^2$ & $y^2$ & $xy$ \\ \hline
10 & 45 & 100 & 2025 & 450 \\
12 & 52 & 144 & 2704 & 624 \\
15 & 60 & 225 & 3600 & 900 \\
18 & 68 & 324 & 4624 & 1224 \\
20 & 72 & 400 & 5184 & 1440 \\
25 & 80 & 625 & 6400 & 2000 \\ \hline
$\sum$=100 & $\sum$=377 & $\sum$=1818 & $\sum$=24537 & $\sum$=6638 \\ \hline
\end{tabular}
\end{center}

$n = 6$.

\textbf{Corrected sums:}
\begin{align*}
s_{xx} &= \sum x^2 - \frac{(\sum x)^2}{n} = 1818 - \frac{10000}{6} = 1818 - 1666.\overline{6} = 151.\overline{3} = \frac{454}{3} \approx 151.333, \\
s_{yy} &= \sum y^2 - \frac{(\sum y)^2}{n} = 24537 - \frac{142129}{6} = 24537 - 23688\frac{1}{6} = 848\frac{5}{6} = \frac{5093}{6} \approx 848.833, \\
s_{xy} &= \sum xy - \frac{\sum x \sum y}{n} = 6638 - \frac{100 \times 377}{6} = 6638 - 6283\frac{1}{3} = 354\frac{2}{3} = \frac{1064}{3} \approx 354.667.
\end{align*}

\textbf{Means:}
\[
\bar{x} = \frac{100}{6} = \frac{50}{3} \approx 16.667, \qquad
\bar{y} = \frac{377}{6} \approx 62.833.
\]

\textbf{(a) Regression line of $y$ on $x$:}
\[
b = \frac{s_{xy}}{s_{xx}} = \frac{1064/3}{454/3} = \frac{1064}{454} = \frac{532}{227} \approx 2.3436.
\]
\[
a = \bar{y} - b\bar{x} = \frac{377}{6} - \frac{532}{227} \times \frac{50}{3}
= \frac{377}{6} - \frac{26600}{681}
= \frac{377 \times 227 - 53200}{1362}.
\]
$377 \times 227 = 85579$. $85579 - 53200 = 32379$.
$a = \frac{32379}{1362} = \frac{10793}{454} \approx 23.773$.
Equation: $\boxed{y = 23.77 + 2.34x}$ (rounded to 2 d.p.).

\textbf{(b) Regression line of $x$ on $y$:}
\[
b' = \frac{s_{xy}}{s_{yy}} = \frac{1064/3}{5093/6} = \frac{1064}{3} \times \frac{6}{5093} = \frac{2128}{5093} \approx 0.4178.
\]
\[
a' = \bar{x} - b'\bar{y} = \frac{50}{3} - \frac{2128}{5093} \times \frac{377}{6}
= \frac{50}{3} - \frac{802256}{30558} \approx 16.667 - 26.253 = -9.586.
\]
Equation: $\boxed{x = -9.59 + 0.418y}$ (rounded to 2 d.p.).

\textbf{(c) Residual Sum of Squares:}
For $y$ on $x$:
\[
\mathrm{RSS}_{y|x} = s_{yy} - \frac{s_{xy}^2}{s_{xx}} = \frac{5093}{6} - \frac{(1064/3)^2}{454/3}
= \frac{5093}{6} - \frac{1064^2}{3 \times 454}
= \frac{5093}{6} - \frac{1132096}{1362}
= \frac{5093 \times 227 - 1132096}{1362}
= \frac{1156111 - 1132096}{1362} = \frac{24015}{1362} \approx 17.63.
\]
For $x$ on $y$:
\[
\mathrm{RSS}_{x|y} = s_{xx} - \frac{s_{xy}^2}{s_{yy}} = \frac{454}{3} - \frac{(1064/3)^2}{5093/6}
= \frac{454}{3} - \frac{1064^2 \times 2}{3 \times 5093}
= \frac{454}{3} - \frac{2264192}{15279}
= \frac{454 \times 5093 - 2264192}{15279}
= \frac{2312222 - 2264192}{15279} = \frac{48030}{15279} \approx 3.14.
\]

\textbf{(d) Estimate mark for $x = 16$ hours (interpolation, since $10 \le 16 \le 25$):}
Use $y$ on $x$ line:
$y = 23.77 + 2.34 \times 16 = 23.77 + 37.44 = 61.21 \approx \boxed{61}$.

\textbf{(e) Estimate revision hours for $y = 65$ (interpolation, since $45 \le 65 \le 80$):}
Use $x$ on $y$ line:
$x = -9.59 + 0.418 \times 65 = -9.59 + 27.17 = 17.58 \approx \boxed{17.6 \text{ hours}}$.

\textbf{Check:} Both lines should pass through $(\bar{x}, \bar{y}) \approx (16.67, 62.83)$.
$y$ on $x$: $23.77 + 2.34 \times 16.67 \approx 23.77 + 39.01 = 62.78 \approx 62.83$ (rounding).
$x$ on $y$: $-9.59 + 0.418 \times 62.83 \approx -9.59 + 26.26 = 16.67$. Correct.
\end{exoresolu}

\courssub{Summary of Key Formulas}

\begin{aretenir}[Least Squares Regression Lines]
\begin{itemize}
\item \textbf{Line of $y$ on $x$}: $y = a + bx$ with $b = \dfrac{s_{xy}}{s_{xx}}$, $a = \bar{y} - b\bar{x}$.
\item \textbf{Line of $x$ on $y$}: $x = a' + b'y$ with $b' = \dfrac{s_{xy}}{s_{yy}}$, $a' = \bar{x} - b'\bar{y}$.
\item \textbf{Both lines pass through} $(\bar{x}, \bar{y})$.
\item \textbf{Product of slopes}: $b \times \frac{1}{b'} = r^2$.
\item \textbf{RSS for $y$ on $x$}: $\mathrm{RSS}_{y|x} = s_{yy} - \dfrac{s_{xy}^2}{s_{xx}} = s_{yy}(1-r^2)$.
\item \textbf{RSS for $x$ on $y$}: $\mathrm{RSS}_{x|y} = s_{xx} - \dfrac{s_{xy}^2}{s_{yy}} = s_{xx}(1-r^2)$.
\item \textbf{Use $y$ on $x$} to predict $y$ from $x$; \textbf{use $x$ on $y$} to predict $x$ from $y$.
\item \textbf{Interpolation} (within data range) is reliable; \textbf{extrapolation} (outside data range) is risky.
\end{itemize}
\end{aretenir}

\begin{exercice}[Practice]
The following data gives the age ($x$ years) and the systolic blood pressure ($y$ mmHg) of 8 men.
\begin{center}
\begin{tabular}{|c|cccccccc|}\hline
$x$ & 35 & 40 & 45 & 50 & 55 & 60 & 65 & 70 \\ \hline
$y$ & 120 & 125 & 132 & 138 & 145 & 150 & 158 & 165 \\ \hline
\end{tabular}
\end{center}
1. Calculate $s_{xx}$, $s_{yy}$, $s_{xy}$.
2. Find the equation of the regression line of $y$ on $x$.
3. Find the equation of the regression line of $x$ on $y$.
4. Calculate the residual sum of squares for each line.
5. Estimate the blood pressure of a man aged 52.
6. Estimate the age of a man with blood pressure 140.
7. Would it be sensible to use the line to estimate the blood pressure of a 20-year-old? Explain.
\end{exercice}

\begin{corrige}[Exercise]
First compute the necessary sums:
$n=8$,
$\sum x = 35+40+45+50+55+60+65+70 = 420$,
$\sum y = 120+125+132+138+145+150+158+165 = 1133$,
$\sum x^2 = 1225+1600+2025+2500+3025+3600+4225+4900 = 23100$,
$\sum y^2 = 14400+15625+17424+19044+21025+22500+24964+27225 = 162207$,
$\sum xy = 4200+5000+5940+6900+7975+9000+10270+11550 = 60835$.

1. Corrected sums:
$s_{xx} = 23100 - \frac{420^2}{8} = 23100 - 22050 = 1050$.
$s_{yy} = 162207 - \frac{1133^2}{8} = 162207 - \frac{1283689}{8} = 162207 - 160461.125 = 1745.875$.
$s_{xy} = 60835 - \frac{420 \times 1133}{8} = 60835 - \frac{475860}{8} = 60835 - 59482.5 = 1352.5$.

Means: $\bar{x} = \frac{420}{8} = 52.5$, $\bar{y} = \frac{1133}{8} = 141.625$.

2. Regression line of $y$ on $x$:
$b = \frac{s_{xy}}{s_{xx}} = \frac{1352.5}{1050} = \frac{541}{420} \approx 1.2881$.
$a = \bar{y} - b\bar{x} = 141.625 - \frac{541}{420} \times 52.5 = 141.625 - \frac{541}{8} = 141.625 - 67.625 = 74.0$.
Equation: $\boxed{y = 74.0 + 1.29x}$ (to 3 s.f.).

3. Regression line of $x$ on $y$:
$b' = \frac{s_{xy}}{s_{yy}} = \frac{1352.5}{1745.875} = \frac{10820}{13967} \approx 0.7747$.
$a' = \bar{x} - b'\bar{y} = 52.5 - \frac{10820}{13967} \times 141.625 \approx 52.5 - 109.71 = -57.21$.
Equation: $\boxed{x = -57.2 + 0.775y}$ (to 3 s.f.).

4. Residual Sum of Squares:
$\mathrm{RSS}_{y|x} = s_{yy} - \frac{s_{xy}^2}{s_{xx}} = 1745.875 - \frac{1352.5^2}{1050} = 1745.875 - 1741.25 = 4.625$.
$\mathrm{RSS}_{x|y} = s_{xx} - \frac{s_{xy}^2}{s_{yy}} = 1050 - \frac{1352.5^2}{1745.875} = 1050 - 1047.18 = 2.82$ (approx).

5. $x=52$ (interpolation, within range 35--70): Use $y$ on $x$ line.
$y = 74.0 + 1.2881 \times 52 = 74.0 + 66.98 = 140.98 \approx \boxed{141 \text{ mmHg}}$.

6. $y=140$ (interpolation, within range 120--165): Use $x$ on $y$ line.
$x = -57.21 + 0.7747 \times 140 = -57.21 + 108.46 = 51.25 \approx \boxed{51 \text{ years}}$.

7. No. $x=20$ is far outside the data range (35--70). Extrapolation assumes the linear trend continues downwards, which is biologically unrealistic (blood pressure would become negative at very young ages). The model is only validated for the observed age range.
\end{corrige}

\coursec{Residual Sum of Squares and the Line of Best Fit}

In the previous section we derived the two least squares regression lines: the line of $y$ on $x$, $y = a + bx$ with $b = \dfrac{s_{xy}}{s_{xx}}$, and the line of $x$ on $y$, $x = a' + b'y$ with $b' = \dfrac{s_{xy}}{s_{yy}}$. In each case the slope and intercept were obtained by \emph{minimising a sum of squares}. In this section we make that idea precise: we define the \cle{residuals} of a fitted line, measure the total error of the fit through the \cle{residual sum of squares} (RSS), and show that the least squares line is, in a very exact sense, the \cle{line of best fit}. This gives us a practical tool for comparing different fitted lines — for example Mayer's line and the least squares line — and a first bridge towards the correlation coefficient of the next section.

\courssub{Residuals: Measuring the Error at Each Point}

Suppose we have $n$ data points $(x_1, y_1), (x_2, y_2), \dots, (x_n, y_n)$ and we have fitted a regression line of $y$ on $x$, $\hat{y} = a + bx$. For a given value $x_i$, the line \emph{predicts} the value $\hat{y}_i = a + bx_i$, while the \emph{observed} value is $y_i$. These two values generally differ, and the difference tells us how well the line performs at that particular point.

\begin{definition}[Residual for the regression of $y$ on $x$]
For the regression line of $y$ on $x$, $\hat{y} = a + bx$, the \cle{residual} at the point $(x_i, y_i)$ is the vertical signed distance from the observed point to the line:
\[ e_i = y_i - \hat{y}_i = y_i - (a + bx_i). \]
A positive residual means the point lies \emph{above} the line; a negative residual means it lies \emph{below} the line.
\end{definition}

\begin{popfigure}
\begin{tikzpicture}[scale=1.05]
\draw[->, popmuted] (0,0) -- (7.2,0) node[below left] {$x$};
\draw[->, popmuted] (0,0) -- (0,4.6) node[below left] {$y$};
\draw[popblue, thick] (0.4,0.9) -- (6.8,3.7) node[pos=0.85, below, sloped, popblue] {$\hat{y}=a+bx$};
\coordinate (P1) at (1,1.7); \coordinate (L1) at (1,1.34);
\coordinate (P2) at (2,1.35); \coordinate (L2) at (2,1.79);
\coordinate (P3) at (3,2.7); \coordinate (L3) at (3,2.23);
\coordinate (P4) at (4.2,2.35); \coordinate (L4) at (4.2,2.77);
\coordinate (P5) at (5.2,3.6); \coordinate (L5) at (5.2,3.22);
\coordinate (P6) at (6.2,3.15); \coordinate (L6) at (6.2,3.67);
\draw[poporange, very thick] (P1) -- (L1);
\draw[poporange, very thick] (P2) -- (L2);
\draw[poporange, very thick] (P3) -- (L3);
\draw[poporange, very thick] (P4) -- (L4);
\draw[poporange, very thick] (P5) -- (L5);
\draw[poporange, very thick] (P6) -- (L6);
\foreach \p in {P1,P2,P3,P4,P5,P6} \fill[popdark] (\p) circle (2pt);
\node[poporange, right] at (3.05,2.5) {$e_3$};
\node[poporange, right] at (5.25,3.45) {$e_5$};
\node[popmuted, align=center, text width=4.2cm] at (5.6,0.8) {vertical segments\\ = residuals $e_i$};
\end{tikzpicture}
\legende{Residuals $e_i = y_i - \hat{y}_i$ shown as vertical segments from each data point to the fitted line of $y$ on $x$.}
\end{popfigure}

For the regression line of $x$ on $y$, $\hat{x} = a' + b'y$, the roles are reversed: we predict $x$ from $y$, and the residual is the \emph{horizontal} signed distance
\[ d_i = x_i - \hat{x}_i = x_i - (a' + b'y_i). \]

\begin{attention}[Vertical versus horizontal residuals]
Residuals for the regression of $y$ on $x$ are \textbf{vertical} distances; residuals for the regression of $x$ on $y$ are \textbf{horizontal} distances. Never mix the two: the two lines answer two different questions ("predict $y$ given $x$" and "predict $x$ given $y$"), and their residuals measure errors in different directions. This is precisely why the two regression lines are generally different.
\end{attention}

\courssub{The Residual Sum of Squares (RSS)}

A single residual tells us about one point. To judge the \emph{global} quality of a fitted line we need to combine all the residuals into a single number. Simply adding them is useless: positive and negative residuals cancel (indeed, for the least squares line one can show that $\sum e_i = 0$ exactly). As with the definition of variance, the natural remedy is to \emph{square} each residual before adding.

\begin{definition}[Residual sum of squares]
For the regression line of $y$ on $x$, the \cle{residual sum of squares} is
\[ \mathrm{RSS}_{y\,on\,x} = \sum_{i=1}^{n} e_i^{\,2} = \sum_{i=1}^{n} \left( y_i - \hat{y}_i \right)^2. \]
For the regression line of $x$ on $y$,
\[ \mathrm{RSS}_{x\,on\,y} = \sum_{i=1}^{n} d_i^{\,2} = \sum_{i=1}^{n} \left( x_i - \hat{x}_i \right)^2. \]
\end{definition}

The RSS is always non-negative, and it equals $0$ only when every residual is zero, that is, when \emph{every point lies exactly on the line}. A small RSS means the points cluster tightly around the line; a large RSS means the points are widely scattered around it.

\courssub{Compact Formulae for the RSS of the Least Squares Lines}

Computing each residual individually is tedious for large data sets. Fortunately, for the \emph{least squares} lines there are closed formulae in terms of the quantities $s_{xx}$, $s_{yy}$ and $s_{xy}$ already computed when fitting the line. Recall our convention (used throughout this chapter):
\[ s_{xx} = \sum (x_i - \bar{x})^2, \qquad s_{yy} = \sum (y_i - \bar{y})^2, \qquad s_{xy} = \sum (x_i - \bar{x})(y_i - \bar{y}). \]

\begin{propriete}[RSS of the least squares lines]
For the least squares regression line of $y$ on $x$,
\[ \mathrm{RSS}_{y\,on\,x} = s_{yy} - \frac{(s_{xy})^2}{s_{xx}}. \]
For the least squares regression line of $x$ on $y$,
\[ \mathrm{RSS}_{x\,on\,y} = s_{xx} - \frac{(s_{xy})^2}{s_{yy}}. \]
(The proof is instructive but is \emph{not} examinable at the GCE; the formulae themselves must be known and used fluently.)
\end{propriete}

\textbf{Derivation (for the line of $y$ on $x$).}\par
The least squares line passes through $(\bar{x}, \bar{y})$, so $\hat{y}_i - \bar{y} = b(x_i - \bar{x})$ with $b = s_{xy}/s_{xx}$. Write $u_i = x_i - \bar{x}$ and $v_i = y_i - \bar{y}$. Then
\[ e_i = y_i - \hat{y}_i = (y_i - \bar{y}) - (\hat{y}_i - \bar{y}) = v_i - b u_i. \]
Hence
\begin{align*}
\mathrm{RSS} &= \sum (v_i - b u_i)^2 = \sum v_i^2 - 2b \sum u_i v_i + b^2 \sum u_i^2 \\
&= s_{yy} - 2b\, s_{xy} + b^2 s_{xx}.
\end{align*}
Substituting $b = \dfrac{s_{xy}}{s_{xx}}$:
\[ \mathrm{RSS} = s_{yy} - 2\,\frac{s_{xy}}{s_{xx}}\, s_{xy} + \frac{(s_{xy})^2}{s_{xx}^2}\, s_{xx} = s_{yy} - \frac{(s_{xy})^2}{s_{xx}}. \]
The derivation for the line of $x$ on $y$ is identical with the roles of $x$ and $y$ exchanged. $\blacksquare$

\begin{methode}[Computing the RSS without listing the residuals]
\begin{enumerate}
\item Compute $\bar{x}$, $\bar{y}$, then $s_{xx}$, $s_{yy}$ and $s_{xy}$ (the same table of sums $\sum x$, $\sum y$, $\sum x^2$, $\sum y^2$, $\sum xy$ used for the regression line).
\item For the line of $y$ on $x$, compute $\mathrm{RSS} = s_{yy} - \dfrac{(s_{xy})^2}{s_{xx}}$.
\item For the line of $x$ on $y$, compute $\mathrm{RSS} = s_{xx} - \dfrac{(s_{xy})^2}{s_{yy}}$.
\item Check: the RSS must satisfy $0 \le \mathrm{RSS}_{y\,on\,x} \le s_{yy}$. A negative value signals an arithmetic error.
\end{enumerate}
\end{methode}

\begin{attention}[Units and the meaning of "small"]
The RSS of the line of $y$ on $x$ is measured in \emph{squared units of $y$}, and the RSS of the line of $x$ on $y$ in squared units of $x$. The two values are \textbf{not} directly comparable with each other. Moreover, "small" is relative: an RSS of $12$ is excellent if $s_{yy} = 4000$ but poor if $s_{yy} = 15$. Always judge the RSS against the total variation $s_{yy}$ (or $s_{xx}$) — this idea is made precise by the product moment correlation coefficient in the next section, since $\dfrac{\mathrm{RSS}_{y\,on\,x}}{s_{yy}} = 1 - r^2$.
\end{attention}

\courssub{The Line of Best Fit}

We now come to the central result of this section, which justifies the name "least squares".

\begin{propriete}[The least squares line is the line of best fit]
Among \emph{all} straight lines $y = a + bx$, the one that makes $\mathrm{RSS} = \sum (y_i - a - bx_i)^2$ as small as possible is exactly the least squares regression line of $y$ on $x$, with $b = \dfrac{s_{xy}}{s_{xx}}$ and $a = \bar{y} - b\bar{x}$. For this reason the least squares line is called the \cle{line of best fit} (in the least squares sense). Similarly, the least squares line of $x$ on $y$ minimises $\sum (x_i - a' - b'y_i)^2$ among all lines $x = a' + b'y$.
\end{propriete}

\textbf{Why this is true.}\par
For fixed $b$, the quantity $S(a,b) = \sum (y_i - a - bx_i)^2$ is minimised with respect to $a$ when $a = \bar{y} - b\bar{x}$ (this is the same "sum of squared deviations is minimised about the mean" fact used for the variance). Substituting and writing $u_i = x_i - \bar{x}$, $v_i = y_i - \bar{y}$ gives
\[ S = \sum (v_i - b u_i)^2 = s_{xx}\, b^2 - 2 s_{xy}\, b + s_{yy}, \]
a quadratic in $b$ with positive leading coefficient $s_{xx} > 0$. Its minimum occurs at $b = \dfrac{s_{xy}}{s_{xx}}$, and the minimum value is $s_{yy} - \dfrac{(s_{xy})^2}{s_{xx}}$, which is precisely the RSS formula above. $\blacksquare$

\begin{aretenir}[Key facts about RSS and the line of best fit]
\begin{itemize}
\item Residual: $e_i = y_i - \hat{y}_i$ (vertical, $y$ on $x$); $d_i = x_i - \hat{x}_i$ (horizontal, $x$ on $y$).
\item $\mathrm{RSS}_{y\,on\,x} = s_{yy} - \dfrac{(s_{xy})^2}{s_{xx}}$ and $\mathrm{RSS}_{x\,on\,y} = s_{xx} - \dfrac{(s_{xy})^2}{s_{yy}}$.
\item The least squares line minimises the RSS: it is the line of best fit.
\item $\mathrm{RSS} = 0 \iff$ all points lie exactly on the line (perfect linear fit, $|r| = 1$).
\item Any other line (e.g.\ Mayer's line) has a \emph{larger} RSS than the least squares line.
\end{itemize}
\end{aretenir}

\courssub{Worked Example: Residuals and RSS for a Small Data Set}

\begin{exoresolu}[Computing residuals and RSS directly]
A small trader in Buea records the number of crates of plantain sold, $y$, against the midday temperature in $°$C, $x$, on five market days:
\begin{center}
\begin{tabular}{|l|ccccc|}
\hline
\rowcolor{popblueL} $x$ & 20 & 22 & 25 & 27 & 30 \\
\hline
\rowcolor{popblueL} $y$ & 14 & 15 & 18 & 17 & 21 \\
\hline
\end{tabular}
\end{center}
The least squares line of $y$ on $x$ is $\hat{y} = 1.1 + 0.7x$. Compute the five residuals, verify that they sum to zero, and deduce the RSS. Check the answer with the compact formula.
\tcblower
\textbf{Solution.} Compute $\hat{y}_i = 1.1 + 0.7 x_i$ and $e_i = y_i - \hat{y}_i$:
\begin{center}
\begin{tabular}{|l|ccccc|c|}
\hline
\rowcolor{popblueL} $x_i$ & 20 & 22 & 25 & 27 & 30 & \\
\hline
$y_i$ & 14 & 15 & 18 & 17 & 21 & \\
\hline
$\hat{y}_i$ & 15.1 & 16.5 & 18.6 & 20.0 & 22.1 & \\
\hline
$e_i$ & $-1.1$ & $-1.5$ & $-0.6$ & $-3.0$ & $-1.1$ & \\
\hline
\end{tabular}
\end{center}
Wait — let us recompute carefully: $\hat{y}_1 = 1.1 + 0.7(20) = 15.1$, so $e_1 = 14 - 15.1 = -1.1$; $\hat{y}_2 = 1.1 + 15.4 = 16.5$, $e_2 = -1.5$; $\hat{y}_3 = 1.1 + 17.5 = 18.6$, $e_3 = -0.6$; $\hat{y}_4 = 1.1 + 18.9 = 20.0$, $e_4 = -3.0$; $\hat{y}_5 = 1.1 + 21.0 = 22.1$, $e_5 = -1.1$.

Hmm, these do not sum to zero, which means the stated line is not exactly the least squares line for this data. Let us instead derive the line ourselves. We have $\bar{x} = \dfrac{124}{5} = 24.8$, $\bar{y} = \dfrac{85}{5} = 17$.
\begin{align*}
s_{xx} &= \sum x^2 - n\bar{x}^2 = (400+484+625+729+900) - 5(24.8)^2 = 3138 - 3075.2 = 62.8,\\
s_{xy} &= \sum xy - n\bar{x}\bar{y} = (280+330+450+459+630) - 5(24.8)(17) = 2149 - 2108 = 41,\\
s_{yy} &= \sum y^2 - n\bar{y}^2 = (196+225+324+289+441) - 5(289) = 1475 - 1445 = 30.
\end{align*}
Hence $b = \dfrac{41}{62.8} \approx 0.6529$ and $a = 17 - 0.6529(24.8) \approx 0.8089$. The least squares line is $\hat{y} = 0.809 + 0.653x$.

The residuals are:
\begin{center}
\begin{tabular}{|l|ccccc|}
\hline
\rowcolor{popblueL} $x_i$ & 20 & 22 & 25 & 27 & 30 \\
\hline
$\hat{y}_i$ & 13.87 & 15.17 & 17.13 & 18.44 & 20.40 \\
\hline
$e_i$ & $0.13$ & $-0.17$ & $0.87$ & $-1.44$ & $0.60$ \\
\hline
\end{tabular}
\end{center}
Sum of residuals: $0.13 - 0.17 + 0.87 - 1.44 + 0.60 \approx 0$ (up to rounding), as expected for a least squares line. Then
\[ \mathrm{RSS} = (0.13)^2 + (0.17)^2 + (0.87)^2 + (1.44)^2 + (0.60)^2 \approx 0.017 + 0.029 + 0.757 + 2.074 + 0.360 \approx 3.24. \]
Check with the compact formula:
\[ \mathrm{RSS} = s_{yy} - \frac{(s_{xy})^2}{s_{xx}} = 30 - \frac{41^2}{62.8} = 30 - \frac{1681}{62.8} \approx 30 - 26.77 = 3.23. \]
The two answers agree (the tiny difference comes from rounding the residuals). \textbf{Conclusion:} $\mathrm{RSS} \approx 3.23$, small compared with $s_{yy} = 30$, so the line fits the data well.
\end{exoresolu}

\courssub{Comparing Two Candidate Lines by their RSS}

Because the least squares line minimises the RSS, any other line fitted to the same data — for instance Mayer's line from Section 2 — must have an RSS \emph{at least as large}. This gives a concrete way to compare two candidate lines: compute the RSS of each; the smaller RSS indicates the better fit (in the least squares sense).

\begin{exoresolu}[Mayer's line versus the least squares line]
For the data of the previous worked example, Mayer's method (splitting into the groups $\{20, 22, 25\}$ and $\{27, 30\}$) gives the line $y = 0.8x - 2.73$ approximately; take $\hat{y} = 0.8x - 2.7$ for simplicity. Compute its RSS and compare with the least squares RSS of $3.23$.
\tcblower
\textbf{Solution.} Predictions from Mayer's line: $\hat{y}_1 = 13.3$, $\hat{y}_2 = 14.9$, $\hat{y}_3 = 17.3$, $\hat{y}_4 = 18.9$, $\hat{y}_5 = 21.3$. Residuals: $0.7,\ 0.1,\ 0.7,\ -1.9,\ -0.3$. Hence
\[ \mathrm{RSS}_{\text{Mayer}} = 0.49 + 0.01 + 0.49 + 3.61 + 0.09 = 4.69. \]
Since $4.69 > 3.23$, Mayer's line fits less well than the least squares line — exactly as the theory guarantees. The difference is modest here because the points are nearly collinear; for more scattered data the gap can be much larger.
\end{exoresolu}

\begin{popfigure}
\begin{tikzpicture}
\begin{axis}[
  width=12cm, height=7.5cm,
  xlabel={$x$}, ylabel={$y$},
  xmin=18, xmax=32, ymin=11, ymax=24,
  grid=both, grid style={popline!40},
  legend style={at={(0.03,0.97)}, anchor=north west, font=\small},
]
\addplot[only marks, mark=*, mark size=2pt, popdark] coordinates {(20,14) (22,15) (25,18) (27,17) (30,21)};
\addlegendentry{data}
\addplot[domain=18:32, samples=2, very thick, popblue] {0.809 + 0.6529*x};
\addlegendentry{least squares: RSS $=3.23$}
\addplot[domain=18:32, samples=2, very thick, dashed, poporange] {0.8*x - 2.7};
\addlegendentry{Mayer: RSS $=4.69$}
\end{axis}
\end{tikzpicture}
\legende{Two candidate lines through the same data. The least squares line (solid blue) has the smaller RSS; Mayer's line (dashed orange) is close but not optimal.}
\end{popfigure}

\courssub{RSS and the Quality of a Linear Model}

The RSS connects directly with correlation. Dividing the RSS formula by $s_{yy}$:
\[ \frac{\mathrm{RSS}_{y\,on\,x}}{s_{yy}} = 1 - \frac{(s_{xy})^2}{s_{xx} s_{yy}} = 1 - r^2, \]
where $r$ is the product moment correlation coefficient studied in the next section. Consequently:
\begin{itemize}
\item $\mathrm{RSS} = 0 \iff r^2 = 1$: \emph{perfect} linear fit, all points exactly on the line;
\item the closer $r^2$ is to $1$, the smaller the RSS relative to the total variation $s_{yy}$, and the better the line explains the data;
\item if $r = 0$ then $\mathrm{RSS} = s_{yy}$: the line explains nothing, and predicting $\hat{y} = \bar{y}$ (a horizontal line) is just as good.
\end{itemize}

\begin{attention}[A small RSS does not prove the model is appropriate]
The least squares line is the best \emph{straight line}, but the true relationship may not be linear at all. A curved pattern can still yield a moderately small RSS while the residuals show a clear systematic pattern (all positive in the middle, negative at the ends, for example). \textbf{Always look at the scatter diagram first}, and if possible at the pattern of the residuals, before trusting a linear model. RSS measures the quality of the fit \emph{within the family of straight lines} — nothing more.
\end{attention}

\begin{savaistu}[Why squares and not absolute values?]
Minimising $\sum |e_i|$ would seem just as natural, and such "least absolute deviations" methods do exist in modern statistics. But squared residuals lead to clean algebra (a quadratic to minimise), explicit formulae for $a$ and $b$, and the elegant identity $\mathrm{RSS} = s_{yy}(1 - r^2)$. This is why Legendre and Gauss, who introduced the least squares method around 1800 to predict the orbit of the asteroid Ceres, chose squares — and why your calculator and the GCE syllabus still use them today.
\end{savaistu}

\begin{exercice}[RSS for the line of $x$ on $y$]
Using the plantain data of the worked example ($s_{xx} = 62.8$, $s_{yy} = 30$, $s_{xy} = 41$):
\begin{enumerate}
\item Find the equation of the least squares regression line of $x$ on $y$.
\item Compute $\mathrm{RSS}_{x\,on\,y}$ using the compact formula.
\item Explain why this value cannot be compared directly with $\mathrm{RSS}_{y\,on\,x} = 3.23$.
\end{enumerate}
\end{exercice}

\begin{corrige}[Exercise 1]
\begin{enumerate}
\item $b' = \dfrac{s_{xy}}{s_{yy}} = \dfrac{41}{30} \approx 1.3667$; $a' = \bar{x} - b'\bar{y} = 24.8 - 1.3667(17) \approx 1.567$. So $\hat{x} = 1.567 + 1.367y$.
\item $\mathrm{RSS}_{x\,on\,y} = s_{xx} - \dfrac{(s_{xy})^2}{s_{yy}} = 62.8 - \dfrac{1681}{30} = 62.8 - 56.03 \approx 6.77$.
\item $\mathrm{RSS}_{x\,on\,y}$ is measured in squared units of $x$ (here $(°\mathrm{C})^2$) while $\mathrm{RSS}_{y\,on\,x}$ is in squared crates; quantities with different units cannot be compared. Each RSS must be judged against its own total variation, $s_{xx}$ or $s_{yy}$ respectively.
\end{enumerate}
\end{corrige}

\begin{exercice}[Comparing two fitted lines]
A student fits two lines to a data set with $s_{yy} = 120$: line $L_1$ gives $\mathrm{RSS} = 30$ and line $L_2$ gives $\mathrm{RSS} = 24$.
\begin{enumerate}
\item Which line is the better fit in the least squares sense? Justify.
\item Can $L_1$ be the least squares regression line of $y$ on $x$? Explain.
\item Compute $r^2$ if $L_2$ is the least squares line.
\end{enumerate}
\end{exercice}

\begin{corrige}[Exercise 2]
\begin{enumerate}
\item $L_2$, because it has the smaller residual sum of squares ($24 < 30$): on average its points lie closer to the line.
\item No. The least squares line \emph{minimises} the RSS over all lines, so no other line can have a strictly smaller RSS. Since $L_2$ achieves $24 < 30$, $L_1$ cannot be the least squares line.
\item $r^2 = 1 - \dfrac{\mathrm{RSS}}{s_{yy}} = 1 - \dfrac{24}{120} = 0.8$. Hence $|r| = \sqrt{0.8} \approx 0.894$, indicating a strong linear relationship.
\end{enumerate}
\end{corrige}

\begin{aretenir}[Summary]
The residual $e_i = y_i - \hat{y}_i$ measures the vertical error of the line of $y$ on $x$ at each point; the RSS, $\sum e_i^2$, measures the global error. The least squares line is the unique line minimising the RSS — the \cle{line of best fit} — and its RSS is given by $s_{yy} - \dfrac{(s_{xy})^2}{s_{xx}}$. Comparing RSS values ranks competing lines, and $\mathrm{RSS} = 0$ characterises a perfect linear fit. The ratio $\mathrm{RSS}/s_{yy} = 1 - r^2$ leads naturally to the product moment correlation coefficient, the subject of the next section.
\end{aretenir}

\coursec{The Product Moment Correlation Coefficient}

In the previous section we used the \cle{residual sum of squares} to decide which of the two least squares regression lines, $y$ on $x$ or $x$ on $y$, fits a given scatter of points better. The RSS measures the \emph{quality of fit} of a chosen line. But it does not answer a more basic question that every GCE examiner loves to ask: \emph{how strong is the linear relationship between the two variables?} A small RSS tells us the line fits well, but the RSS is measured in squared units of $y$, so its size alone does not allow us to compare different data sets. What we need is a \emph{dimensionless number}, lying on a fixed scale, that summarises the strength and direction of the linear association. That number is the \cle{product moment correlation coefficient}, denoted $r$, and it is the subject of this section.

\courssub{Definition of the Product Moment Correlation Coefficient}

Think of a scatter diagram. If the points cluster tightly around an upward sloping line, we say $x$ and $y$ are strongly positively correlated. If they cluster around a downward sloping line, the correlation is strongly negative. If the points form a shapeless cloud, there is little or no linear correlation. We already possess all the ingredients needed to measure this: the covariance $s_{xy}$ tells us the \emph{direction} of the association (positive when large $x$ goes with large $y$, negative otherwise), while $s_{xx}$ and $s_{yy}$ measure the spreads of $x$ and $y$ separately. Dividing the covariance by the geometric mean of the two spreads removes the units and bounds the result between $-1$ and $1$.

\begin{definition}[Product moment correlation coefficient]
For a sample of $n$ pairs $(x_1,y_1),\dots,(x_n,y_n)$, the \cle{product moment correlation coefficient} (PMCC) is
\[
r \;=\; \frac{s_{xy}}{\sqrt{s_{xx}\,s_{yy}}}
\;=\; \frac{\displaystyle\sum (x_i-\bar{x})(y_i-\bar{y})}{\sqrt{\displaystyle\sum (x_i-\bar{x})^2 \;\sum (y_i-\bar{y})^2}} .
\]
It is also called \emph{Pearson's correlation coefficient}.
\end{definition}

For hand computation at the GCE, we never work with deviations directly; we use the summary sums $\sum x$, $\sum y$, $\sum x^2$, $\sum y^2$, $\sum xy$.

\begin{propriete}[Computational formula for $r$]
Using $S_{xy}=\sum xy - n\bar{x}\bar{y}$, $S_{xx}=\sum x^2 - n\bar{x}^2$ and $S_{yy}=\sum y^2 - n\bar{y}^2$,
\[
r \;=\; \frac{\sum xy - n\bar{x}\bar{y}}{\sqrt{\bigl(\sum x^2 - n\bar{x}^2\bigr)\bigl(\sum y^2 - n\bar{y}^2\bigr)}} .
\]
This formula is examinable: candidates must be able to derive the numerator and denominator forms from the definitions of $S_{xy}$, $S_{xx}$ and $S_{yy}$ (expanding $\sum(x_i-\bar{x})(y_i-\bar{y})$ and using $\sum x_i = n\bar{x}$).
\end{propriete}

\courssub{Properties and Interpretation of $r$}

\begin{propriete}[Range and boundary values of $r$]
For every bivariate data set,
\[
-1 \;\le\; r \;\le\; 1 .
\]
Moreover:
\begin{itemize}
\item $r = 1$ if and only if all the points lie exactly on a straight line of \emph{positive} gradient (perfect positive linear correlation);
\item $r = -1$ if and only if all the points lie exactly on a straight line of \emph{negative} gradient (perfect negative linear correlation);
\item $r = 0$ when $S_{xy}=0$: there is \emph{no linear} correlation.
\end{itemize}
The inequality $-1\le r\le 1$ follows from the Cauchy--Schwarz inequality applied to the deviations $x_i-\bar{x}$ and $y_i-\bar{y}$; its proof is \emph{not} required at the GCE, but the result and the meaning of the boundary values must be known.
\end{propriete}

Intermediate values are interpreted by their \emph{sign} (direction) and their \emph{absolute value} (strength). A commonly used guide is:

\begin{center}
\begin{tabularx}{\linewidth}{|l|X|}
\hline
\rowcolor{popblueL} Value of $r$ & Interpretation \\
\hline
$0.8 \le |r| \le 1$ & Strong linear correlation \\
\hline
$0.5 \le |r| < 0.8$ & Moderate linear correlation \\
\hline
$0.2 \le |r| < 0.5$ & Weak linear correlation \\
\hline
$|r| < 0.2$ & Very weak or no linear correlation \\
\hline
\end{tabularx}
\end{center}

These bands are conventions, not laws: always interpret $r$ \emph{in the context of the question} (for example, ``there is a strong positive linear correlation between revision hours and mock examination marks'').

\begin{popfigure}
\begin{center}
\begin{tikzpicture}
\begin{axis}[width=6.2cm,height=5.2cm,title={$r \approx 1$},xmin=0,xmax=10,ymin=0,ymax=10,xtick=\empty,ytick=\empty,axis lines=left,title style={font=\small}]
\addplot[only marks,mark=*,mark size=1.4pt,color=popblue] coordinates {(1,1.1)(2,2.05)(3,2.9)(4,4.1)(5,5)(6,5.95)(7,7.1)(8,8)(9,9.05)};
\end{axis}
\end{tikzpicture}\hfill
\begin{tikzpicture}
\begin{axis}[width=6.2cm,height=5.2cm,title={$r \approx 0.6$},xmin=0,xmax=10,ymin=0,ymax=10,xtick=\empty,ytick=\empty,axis lines=left,title style={font=\small}]
\addplot[only marks,mark=*,mark size=1.4pt,color=popgreen] coordinates {(1,2.5)(2,1.2)(3,4.5)(4,3)(5,6.5)(6,4.8)(7,8)(8,6.2)(9,9)};
\end{axis}
\end{tikzpicture}

\vspace{0.4cm}
\begin{tikzpicture}
\begin{axis}[width=6.2cm,height=5.2cm,title={$r \approx 0$},xmin=0,xmax=10,ymin=0,ymax=10,xtick=\empty,ytick=\empty,axis lines=left,title style={font=\small}]
\addplot[only marks,mark=*,mark size=1.4pt,color=poporange] coordinates {(1,5)(2,8)(3,2)(4,7)(5,3)(6,6)(7,2.5)(8,7.5)(9,4)};
\end{axis}
\end{tikzpicture}\hfill
\begin{tikzpicture}
\begin{axis}[width=6.2cm,height=5.2cm,title={$r \approx -0.9$},xmin=0,xmax=10,ymin=0,ymax=10,xtick=\empty,ytick=\empty,axis lines=left,title style={font=\small}]
\addplot[only marks,mark=*,mark size=1.4pt,color=poppink] coordinates {(1,9.2)(2,8.1)(3,6.8)(4,6.2)(5,4.9)(6,4.2)(7,2.8)(8,2.1)(9,1)};
\end{axis}
\end{tikzpicture}
\end{center}
\legende{Four scatter diagrams illustrating perfect positive, moderate positive, zero, and strong negative linear correlation.}
\end{popfigure}

\courssub{Link with the Two Regression Gradients}

Recall from the least squares section that the regression line of $y$ on $x$ has gradient $b = \dfrac{S_{xy}}{S_{xx}}$, while the regression line of $x$ on $y$ has gradient $b' = \dfrac{S_{xy}}{S_{yy}}$.

\begin{propriete}[$r^2$ is the product of the two regression gradients]
\[
b \cdot b' \;=\; \frac{S_{xy}}{S_{xx}}\cdot\frac{S_{xy}}{S_{yy}} \;=\; \frac{S_{xy}^2}{S_{xx}S_{yy}} \;=\; r^2 .
\]
Consequently $r = \pm\sqrt{b\,b'}$, the sign being the common sign of $b$ and $b'$ (which is the sign of $S_{xy}$). This identity is frequently examined: given the two regression lines, one can recover $r$ without touching the raw data.
\end{propriete}

This result also explains why the two regression lines coincide exactly when $r = \pm 1$: then $|b\cdot b'| = 1$, the two lines are mutual inverses with the same slope through $(\bar{x},\bar{y})$, i.e.\ the same line.

\begin{exemplebox}[Recovering $r$ from the regression lines]
The least squares regression lines for a data set are $y = 0.8x + 2$ ($y$ on $x$) and $x = 0.45y + 1$ ($x$ on $y$). Rewriting the second as $x$ on $y$ gives gradient $b' = 0.45$. Hence
\[
r^2 = 0.8 \times 0.45 = 0.36, \qquad r = +0.6
\]
(positive because both gradients are positive). There is a moderate positive linear correlation.
\end{exemplebox}

\courssub{Effect of a Linear Change of Scale (GCE Extra)}

\begin{propriete}[Invariance of $r$ under coding]
If the data are coded by linear transformations $u = \dfrac{x-a}{p}$ and $v = \dfrac{y-c}{q}$ with $p,q>0$, then the correlation coefficient of the coded data equals that of the original data:
\[
r_{uv} = r_{xy}.
\]
Indeed, deviations satisfy $u_i-\bar{u} = (x_i-\bar{x})/p$, so $S_{uu}=S_{xx}/p^2$, $S_{vv}=S_{yy}/q^2$, $S_{uv}=S_{xy}/(pq)$, and the factors $p,q$ cancel in the ratio. This is a useful GCE extra: coding large or awkward values (e.g.\ years, or prices in thousands of FCFA) does not change $r$, and can greatly simplify the arithmetic.
\end{propriete}

\courssub{Computing $r$ from a Table of Sums}

\begin{methode}[Computing and interpreting $r$]
\begin{enumerate}
\item From the data (or the question), obtain the five sums $\sum x$, $\sum y$, $\sum x^2$, $\sum y^2$, $\sum xy$ and the sample size $n$.
\item Compute the means $\bar{x} = \sum x / n$ and $\bar{y} = \sum y / n$.
\item Compute $S_{xy} = \sum xy - n\bar{x}\bar{y}$, $S_{xx} = \sum x^2 - n\bar{x}^2$, $S_{yy} = \sum y^2 - n\bar{y}^2$.
\item Substitute into $r = \dfrac{S_{xy}}{\sqrt{S_{xx}S_{yy}}}$ and give the value to 3 or 4 significant figures.
\item Interpret the sign and magnitude of $r$ \emph{in the context} of the problem.
\end{enumerate}
\end{methode}

\begin{exoresolu}[GCE-style worked example]
A teacher in Buea records, for 8 candidates, the number of hours $x$ spent revising Pure Mathematics during the week before a mock examination and the mark $y$ (out of 100) obtained. The results are summarised by
\[
\sum x = 96,\quad \sum y = 424,\quad \sum x^2 = 1392,\quad \sum y^2 = 23938,\quad \sum xy = 5802.
\]
Calculate the product moment correlation coefficient between revision time and mark, and interpret your result.
\tcblower
\textbf{Solution.} Here $n = 8$, so
\[
\bar{x} = \frac{96}{8} = 12, \qquad \bar{y} = \frac{424}{8} = 53.
\]
Compute the three summary quantities:
\begin{align*}
S_{xy} &= \sum xy - n\bar{x}\bar{y} = 5802 - 8(12)(53) = 5802 - 5088 = 714,\\
S_{xx} &= \sum x^2 - n\bar{x}^2 = 1392 - 8(144) = 1392 - 1152 = 240,\\
S_{yy} &= \sum y^2 - n\bar{y}^2 = 23938 - 8(2809) = 23938 - 22472 = 1466.
\end{align*}
Therefore
\[
r = \frac{714}{\sqrt{240 \times 1466}} = \frac{714}{\sqrt{351840}} = \frac{714}{593.16\ldots} \approx 0.904.
\]
\textbf{Interpretation.} Since $r \approx 0.90$ is close to $1$, there is a \emph{strong positive linear correlation} between revision time and mock mark: candidates who revised more tended to obtain higher marks.
\end{exoresolu}

\courssub{Limitations of the Product Moment Correlation Coefficient}

\begin{attention}[$r \approx 0$ does not mean ``no relationship'']
The PMCC measures \emph{linear} association only. A data set can have $r \approx 0$ while $x$ and $y$ are perfectly related by a curve. For example, points lying exactly on the parabola $y = x^2$ for $x$ symmetric about $0$ give $S_{xy} = 0$, hence $r = 0$, yet $y$ is completely determined by $x$. Always look at the scatter diagram before concluding that two variables are unrelated.
\end{attention}

\begin{attention}[Correlation is not causation]
A value of $r$ close to $\pm 1$ shows that $x$ and $y$ vary together linearly; it does \emph{not} prove that changes in $x$ \emph{cause} changes in $y$. Both variables may be driven by a hidden third factor (a \emph{confounding variable}). For instance, ice cream sales and cases of sunburn are strongly correlated in many towns, not because one causes the other, but because both increase in hot sunny weather. In GCE answers, never write ``$x$ causes $y$'' on the basis of $r$ alone.
\end{attention}

Two further cautions deserve mention. First, $r$ is \emph{sensitive to outliers}: a single point far from the main cloud can drag $r$ towards $+1$ or $-1$ and completely distort the conclusion. Second, $r$ computed on a restricted range of $x$ may badly understate the true association. Both issues are visible only on the scatter diagram, which is why examiners reward candidates who sketch or refer to it.

\begin{aretenir}[Key points on the PMCC]
\begin{itemize}
\item $r = \dfrac{S_{xy}}{\sqrt{S_{xx}S_{yy}}}$ with $S_{xy}=\sum xy - n\bar{x}\bar{y}$, etc.; always compute the five sums first.
\item $-1 \le r \le 1$; $r=\pm1$ means a perfect linear relationship, $r=0$ means no \emph{linear} relationship.
\item $r^2 = b\,b'$, the product of the two least squares gradients; the sign of $r$ is the common sign of $b$ and $b'$.
\item $r$ is unchanged by linear coding of the data.
\item $r$ measures linear association only, is sensitive to outliers, and never by itself proves causation.
\end{itemize}
\end{aretenir}

\begin{exercice}[Computing and using $r$]
For six shops in Douala, the daily number $x$ of customers and the daily revenue $y$ (in thousands of FCFA) give
\[
\sum x = 330,\quad \sum y = 612,\quad \sum x^2 = 20100,\quad \sum y^2 = 64452,\quad \sum xy = 36420.
\]
\begin{enumerate}
\item Calculate $r$ and interpret it.
\item Given that the regression line of $y$ on $x$ has gradient $b = 1.5$, find the gradient $b'$ of the regression line of $x$ on $y$.
\end{enumerate}
\end{exercice}

\begin{corrige}[Exercise 1]
\textbf{1.} $n=6$, $\bar{x} = 55$, $\bar{y} = 102$.
\begin{align*}
S_{xy} &= 36420 - 6(55)(102) = 36420 - 33660 = 2760,\\
S_{xx} &= 20100 - 6(55)^2 = 20100 - 18150 = 1950,\\
S_{yy} &= 64452 - 6(102)^2 = 64452 - 62424 = 2028.
\end{align*}
\[
r = \frac{2760}{\sqrt{1950 \times 2028}} = \frac{2760}{\sqrt{3954600}} \approx \frac{2760}{1988.6} \approx 0.888.
\]
There is a strong positive linear correlation between the number of customers and daily revenue.
\textbf{2.} From $r^2 = b\,b'$: $b' = \dfrac{r^2}{b} = \dfrac{0.888^2}{1.5} \approx \dfrac{0.788}{1.5} \approx 0.526$.
\end{corrige}

\coursec{Rank Correlation: Spearman and Kendall}

In the previous section we studied the \cle{product moment correlation coefficient} \(r\), which measures the strength of a \emph{linear} relationship between two quantitative variables. However, real data often do not satisfy the assumptions required for Pearson's \(r\): the variables may be \emph{ordinal} (only ranks are available), the relationship may be \emph{non‑linear but monotonic}, or the data may contain outliers that distort \(r\). In such situations we turn to \cle{rank correlation} coefficients, which are based solely on the \emph{order} of the observations. This section introduces the two most widely used rank correlation measures: \textbf{Spearman's rank correlation coefficient} (\(R_s\)) and \textbf{Kendall's rank correlation coefficient} (\(\tau\)).

\courssub{Ranking Data and Tied Ranks}

Before we can compute any rank correlation, we must assign \cle{ranks} to the observations. The smallest value receives rank 1, the next smallest rank 2, and so on up to \(n\). Alternatively, one may rank in descending order (largest = 1) provided the \emph{same convention} is used for both variables.

\begin{definition}[Ranks]
Given a sample of \(n\) observations \(x_1, x_2, \dots, x_n\), the \emph{rank} of \(x_i\) is its position when the data are sorted in ascending order. If several observations have the same value (a \emph{tie}), each of them receives the \emph{average} of the ranks they would have occupied had they been slightly different.
\end{definition}

\begin{attention}[Tied Ranks]
When ties occur, \textbf{always assign the mean rank} to each tied item. For example, if the 3rd and 4th smallest values are equal, they both receive rank \(\frac{3+4}{2}=3.5\). The next distinct value then receives rank 5. Failing to do this systematically is a common error that invalidates both Spearman's and Kendall's formulas.
\end{attention}

\begin{exemplebox}[Ranking with Ties]
Consider the marks of 7 students in a Mathematics test:
\[
52,\; 58,\; 58,\; 61,\; 65,\; 65,\; 65.
\]
Sorting gives the same order. The ranks are assigned as follows:
\begin{itemize}
\item 52 is smallest → rank 1.
\item The two 58's occupy positions 2 and 3 → each gets rank \(\frac{2+3}{2}=2.5\).
\item 61 is next → rank 4.
\item The three 65's occupy positions 5, 6, 7 → each gets rank \(\frac{5+6+7}{3}=6\).
\end{itemize}
The rank vector is \((1,\; 2.5,\; 2.5,\; 4,\; 6,\; 6,\; 6)\).
\end{exemplebox}

\begin{popfigure}
\begin{tikzpicture}[scale=0.9]
% Table for tied ranks example
\draw[popline] (0,0) grid (7,2);
\node[anchor=south] at (3.5,2.2) {\textbf{Ranking with Tied Values}};
% Header
\node[anchor=center] at (0.5,1.5) {\textbf{Value}};
\node[anchor=center] at (1.5,1.5) {52};
\node[anchor=center] at (2.5,1.5) {58};
\node[anchor=center] at (3.5,1.5) {58};
\node[anchor=center] at (4.5,1.5) {61};
\node[anchor=center] at (5.5,1.5) {65};
\node[anchor=center] at (6.5,1.5) {65};
\node[anchor=center] at (7.5,1.5) {65};
% Rank row
\node[anchor=center] at (0.5,0.5) {\textbf{Rank}};
\node[anchor=center] at (1.5,0.5) {1};
\node[anchor=center, poporange] at (2.5,0.5) {2.5};
\node[anchor=center, poporange] at (3.5,0.5) {2.5};
\node[anchor=center] at (4.5,0.5) {4};
\node[anchor=center, popgreen] at (5.5,0.5) {6};
\node[anchor=center, popgreen] at (6.5,0.5) {6};
\node[anchor=center, popgreen] at (7.5,0.5) {6};
% Highlight tied groups
\draw[poporange, thick] (2,1) rectangle (4,2);
\draw[popgreen, thick] (5,1) rectangle (8,2);
\node[poporange, anchor=south] at (3,2.1) {Tied pair};
\node[popgreen, anchor=south] at (6.5,2.1) {Tied triple};
\end{tikzpicture}
\legende{Illustration of tied ranks: the two 58's share rank 2.5 (orange), the three 65's share rank 6 (green).}
\end{popfigure}

\courssub{Spearman's Rank Correlation Coefficient}

Spearman's coefficient, denoted \(R_s\) (or sometimes \(\rho\)), is simply the Pearson correlation coefficient computed on the \emph{ranks} of the two variables. When there are no ties, it reduces to a very simple formula.

\begin{definition}[Spearman's Rank Correlation Coefficient]
Let \((x_i, y_i),\; i=1,\dots,n\) be paired observations. Let \(R(x_i)\) and \(R(y_i)\) be their ranks. Spearman's rank correlation coefficient is
\[
R_s = \frac{\sum_{i=1}^n \bigl(R(x_i)-\overline{R_x}\bigr)\bigl(R(y_i)-\overline{R_y}\bigr)}
{\sqrt{\sum_{i=1}^n \bigl(R(x_i)-\overline{R_x}\bigr)^2 \sum_{i=1}^n \bigl(R(y_i)-\overline{R_y}\bigr)^2}}.
\]
If there are \emph{no tied ranks}, this simplifies to
\[
\boxed{R_s = 1 - \frac{6\sum_{i=1}^n d_i^2}{n(n^2-1)},\qquad d_i = R(x_i)-R(y_i).}
\]
If ties are present, one \textbf{must} compute the Pearson correlation coefficient on the ranked data (using the averaged ranks). The simplified formula \(1 - \frac{6\sum d_i^2}{n(n^2-1)}\) is \textbf{not valid} when ties exist. In GCE examinations, the simplified formula is used only when \emph{explicitly stated} or \emph{clearly evident} that there are no ties; otherwise compute the Pearson correlation on the ranked data.
\end{definition}

\begin{propriete}[Properties of \(R_s\)]
\begin{itemize}
\item \(-1 \le R_s \le 1\).
\item \(R_s = 1\) \(\iff\) the two rankings are in perfect agreement (monotonically increasing).
\item \(R_s = -1\) \(\iff\) the two rankings are in perfect disagreement (one is the reverse of the other).
\item \(R_s = 0\) indicates no monotonic association.
\item \(R_s\) measures the strength of a \emph{monotonic} relationship, not necessarily linear.
\end{itemize}
\end{propriete}

\begin{methode}[Computing Spearman's \(R_s\) (Step‑by‑Step)]
\begin{enumerate}
\item List the paired observations \((x_i, y_i)\).
\item Assign ranks \(R(x_i)\) to the \(x\)'s and \(R(y_i)\) to the \(y\)'s. \textbf{Handle ties by averaging ranks.}
\item Check for ties:
\begin{itemize}
\item \textbf{No ties:} compute \(d_i = R(x_i)-R(y_i)\) for each pair, then \(\sum d_i^2\) and apply \(R_s = 1 - \frac{6\sum d_i^2}{n(n^2-1)}\).
\item \textbf{Ties present:} compute the Pearson correlation coefficient on the ranked data (using the averaged ranks).
\end{itemize}
\item Interpret the value: close to 1 → strong positive monotonic agreement; close to -1 → strong negative monotonic agreement; near 0 → little or no monotonic association.
\end{enumerate}
\end{methode}

\begin{exoresolu}[Spearman's \(R_s\) with No Ties]
Two judges rank 8 candidates in a singing competition. The ranks given by Judge A and Judge B are:
\[
\begin{array}{c|cccccccc}
\text{Candidate} & 1 & 2 & 3 & 4 & 5 & 6 & 7 & 8 \\ \hline
\text{Judge A} & 1 & 2 & 3 & 4 & 5 & 6 & 7 & 8 \\
\text{Judge B} & 2 & 1 & 4 & 3 & 6 & 5 & 8 & 7
\end{array}
\]
Compute Spearman's rank correlation coefficient.
\tcblower
\textbf{Solution.}
Here the ranks are already given and there are \textbf{no ties}. We compute \(d_i = R_A - R_B\):
\[
\begin{array}{c|cccccccc}
\text{Candidate} & 1 & 2 & 3 & 4 & 5 & 6 & 7 & 8 \\ \hline
R_A & 1 & 2 & 3 & 4 & 5 & 6 & 7 & 8 \\
R_B & 2 & 1 & 4 & 3 & 6 & 5 & 8 & 7 \\
d_i & -1 & 1 & -1 & 1 & -1 & 1 & -1 & 1 \\
d_i^2 & 1 & 1 & 1 & 1 & 1 & 1 & 1 & 1
\end{array}
\]
\(\sum d_i^2 = 8\). With \(n=8\):
\[
R_s = 1 - \frac{6 \times 8}{8(8^2-1)} = 1 - \frac{48}{8 \times 63} = 1 - \frac{48}{504} = 1 - 0.09524 = 0.90476.
\]
The judges show a very strong positive agreement in their rankings.
\end{exoresolu}

\begin{exoresolu}[Spearman's \(R_s\) with Tied Ranks]
Eight students are ranked by their Mathematics teacher (Mr. T) and their Physics teacher (Ms. K) based on test scores. The raw scores and ranks are:
\[
\begin{array}{c|cccccccc}
\text{Student} & A & B & C & D & E & F & G & H \\ \hline
\text{Maths score} & 72 & 65 & 65 & 58 & 58 & 58 & 50 & 45 \\
\text{Physics score} & 80 & 75 & 70 & 70 & 65 & 60 & 60 & 55
\end{array}
\]
Compute Spearman's rank correlation coefficient.
\tcblower
\textbf{Solution.}

\noindent\textbf{Step 1: Assign ranks to Maths scores (with ties).}
Sorted Maths: 45, 50, 58, 58, 58, 65, 65, 72.
\begin{itemize}
\item 45 (H) → rank 1
\item 50 (G) → rank 2
\item 58 (D,E,F) occupy positions 3,4,5 → each gets rank \(\frac{3+4+5}{3}=4\)
\item 65 (B,C) occupy positions 6,7 → each gets rank \(\frac{6+7}{2}=6.5\)
\item 72 (A) → rank 8
\end{itemize}
Maths ranks \(R_M\): A=8, B=6.5, C=6.5, D=4, E=4, F=4, G=2, H=1.

\noindent\textbf{Step 2: Assign ranks to Physics scores (with ties).}
Sorted Physics: 55, 60, 60, 65, 70, 70, 75, 80.
\begin{itemize}
\item 55 (H) → rank 1
\item 60 (F,G) occupy positions 2,3 → each gets rank \(\frac{2+3}{2}=2.5\)
\item 65 (E) → rank 4
\item 70 (C,D) occupy positions 5,6 → each gets rank \(\frac{5+6}{2}=5.5\)
\item 75 (B) → rank 7
\item 80 (A) → rank 8
\end{itemize}
Physics ranks \(R_P\): A=8, B=7, C=5.5, D=5.5, E=4, F=2.5, G=2.5, H=1.

\noindent\textbf{Step 3: Compute Pearson correlation on the ranks.}
\[
\begin{array}{c|cccccccc}
\text{Student} & A & B & C & D & E & F & G & H \\ \hline
R_M & 8 & 6.5 & 6.5 & 4 & 4 & 4 & 2 & 1 \\
R_P & 8 & 7 & 5.5 & 5.5 & 4 & 2.5 & 2.5 & 1
\end{array}
\]
\(\overline{R_M} = \overline{R_P} = \frac{8+7+6+5+4+3+2+1}{8} = 4.5\) (sum of ranks 1 to 8 is always 36).
Compute deviations and products:
\[
\begin{array}{c|cccccccc}
\text{Student} & A & B & C & D & E & F & G & H \\ \hline
R_M - 4.5 & 3.5 & 2.0 & 2.0 & -0.5 & -0.5 & -0.5 & -2.5 & -3.5 \\
R_P - 4.5 & 3.5 & 2.5 & 1.0 & 1.0 & -0.5 & -2.0 & -2.0 & -3.5 \\
\text{Product} & 12.25 & 5.00 & 2.00 & -0.50 & 0.25 & 1.00 & 5.00 & 12.25
\end{array}
\]
\(\sum (R_M - \overline{R_M})(R_P - \overline{R_P}) = 37.25\).
\(\sum (R_M - \overline{R_M})^2 = 12.25+4+4+0.25+0.25+0.25+6.25+12.25 = 39.5\).
\(\sum (R_P - \overline{R_P})^2 = 12.25+6.25+1+1+0.25+4+4+12.25 = 41\).
\[
R_s = \frac{37.25}{\sqrt{39.5 \times 41}} = \frac{37.25}{\sqrt{1619.5}} = \frac{37.25}{40.243} \approx 0.9256.
\]
There is a very strong positive monotonic association between the two teachers' rankings.
\end{exoresolu}

\courssub{Kendall's Rank Correlation Coefficient}

Kendall's \(\tau\) (tau) takes a different approach: instead of comparing rank differences, it compares the \emph{order} of every pair of observations. It counts how many pairs are \emph{concordant} (same order in both rankings) and how many are \emph{discordant} (opposite order).

\begin{definition}[Concordant and Discordant Pairs]
Consider two distinct observations \(i\) and \(j\) (\(i<j\)). Let their ranks in the two variables be \((R_x(i), R_y(i))\) and \((R_x(j), R_y(j))\).
\begin{itemize}
\item The pair is \textbf{concordant} if \((R_x(i)-R_x(j))(R_y(i)-R_y(j)) > 0\) (both differences have the same sign).
\item The pair is \textbf{discordant} if \((R_x(i)-R_x(j))(R_y(i)-R_y(j)) < 0\) (differences have opposite signs).
\item If either difference is zero (a tie in one variable), the pair is \textbf{neither concordant nor discordant} and is excluded from the count in the basic \(\tau\) formula. (Tie-corrected versions \(\tau_b\) and \(\tau_c\) exist but are beyond the GCE syllabus.)
\end{itemize}
\end{definition}

Let \(C\) be the number of concordant pairs and \(D\) the number of discordant pairs. The total number of pairs is \(\binom{n}{2} = \frac{n(n-1)}{2}\). Kendall's \(\tau\) is defined as
\[
\boxed{\tau = \frac{C - D}{\frac{1}{2}n(n-1)} = \frac{2(C-D)}{n(n-1)}.}
\]
Often the \emph{score} \(S = C - D\) is used, giving \(\tau = \frac{2S}{n(n-1)}\).

\begin{propriete}[Properties of Kendall's \(\tau\)]
\begin{itemize}
\item \(-1 \le \tau \le 1\).
\item \(\tau = 1\) \(\iff\) perfect agreement (all pairs concordant).
\item \(\tau = -1\) \(\iff\) perfect disagreement (all pairs discordant).
\item \(\tau = 0\) \(\iff\) number of concordant pairs equals number of discordant pairs (no monotonic association).
\item \(\tau\) is more robust to outliers than \(R_s\) because it only uses the \emph{sign} of the difference, not its magnitude.
\end{itemize}
\end{propriete}

\begin{methode}[Counting Concordant and Discordant Pairs (Kendall's \(\tau\))]
\begin{enumerate}
\item Arrange the data so that one variable (say \(X\)) is in \textbf{natural order} \(1,2,\dots,n\). Reorder the \(Y\) ranks accordingly.
\item For each \(i\) from 1 to \(n-1\), compare \(Y_i\) with all \(Y_j\) for \(j>i\).
\item Count:
\begin{itemize}
\item \textbf{Concordant} if \(Y_j > Y_i\) (since \(X_j > X_i\) by construction).
\item \textbf{Discordant} if \(Y_j < Y_i\).
\item Ignore if \(Y_j = Y_i\) (tie in \(Y\)).
\end{itemize}
\item Sum \(C\) and \(D\) over all \(i\). Compute \(S = C - D\) and \(\tau = \frac{2S}{n(n-1)}\).
\end{enumerate}
\end{methode}

\begin{popfigure}
\begin{tikzpicture}[scale=0.85]
% Diagram for Kendall's concordant/discordant pairs
\def\n{5}
% Left ranking (X) - natural order 1..5
\foreach \i in {1,...,\n} {
  \node[draw, rectangle, minimum width=1.2cm, minimum height=0.7cm, fill=popblueL] (L\i) at (0, -\i*1.2) {\i};
}
% Right ranking (Y) - some permutation
% Let Y ranks be: 2, 1, 4, 3, 5 (for items 1..5)
\foreach \i/\y in {1/2, 2/1, 3/4, 4/3, 5/5} {
  \node[draw, rectangle, minimum width=1.2cm, minimum height=0.7cm, fill=popgreenL] (R\i) at (5, -\y*1.2) {\i};
}
% Arrows for concordant pairs (green) and discordant (red)
% Pair (1,2): X: 1<2, Y: 2>1 -> discordant
\draw[poporange, thick, ->] (L1.east) .. controls (2.5,-1.2) and (2.5,-2.4) .. (R2.west);
% Pair (1,3): X:1<3, Y:2<4 -> concordant
\draw[popgreen, thick, ->] (L1.east) .. controls (2.5,-1.2) and (2.5,-4.8) .. (R3.west);
% Pair (2,3): X:2<3, Y:1<4 -> concordant
\draw[popgreen, thick, ->] (L2.east) .. controls (2.5,-2.4) and (2.5,-4.8) .. (R3.west);
% Pair (3,4): X:3<4, Y:4>3 -> discordant
\draw[poporange, thick, ->] (L3.east) .. controls (2.5,-3.6) and (2.5,-4.8) .. (R4.west);
% Labels
\node[anchor=east] at (-0.8, -0.6) {\textbf{Ranking X (natural order)}};
\node[anchor=west] at (5.8, -0.6) {\textbf{Ranking Y}};
\node[popgreen, anchor=south] at (2.5, -0.2) {Concordant};
\node[poporange, anchor=north] at (2.5, -6.2) {Discordant};
\end{tikzpicture}
\legende{Illustration of concordant (green) and discordant (orange) pairs for Kendall's \(\tau\). The left column is in natural order; arrows connect to the right column. Pair (1,2) is discordant because item 1 is above item 2 in X but below in Y.}
\end{popfigure}

\begin{exoresolu}[Kendall's \(\tau\) on the Same Data]
Using the same judges' rankings as in the Spearman example (8 candidates, no ties), compute Kendall's \(\tau\).
\tcblower
\textbf{Solution.}
We have \(n=8\). Judge A's ranks are already in natural order \(1,2,\dots,8\). Judge B's ranks in that order are:
\[
Y = (2, 1, 4, 3, 6, 5, 8, 7).
\]
We count concordant and discordant pairs by scanning each \(i\) and comparing with later \(j\).

\begin{itemize}
\item \(i=1, Y_1=2\): later values \(1,4,3,6,5,8,7\). Values \(<2\): only 1 → 1 discordant. Values \(>2\): 4,3,6,5,8,7 → 6 concordant.
\item \(i=2, Y_2=1\): later values \(4,3,6,5,8,7\). All are \(>1\) → 6 concordant, 0 discordant.
\item \(i=3, Y_3=4\): later values \(3,6,5,8,7\). \(<4\): 3 → 1 discordant. \(>4\): 6,5,8,7 → 4 concordant.
\item \(i=4, Y_4=3\): later values \(6,5,8,7\). All \(>3\) → 4 concordant.
\item \(i=5, Y_5=6\): later values \(5,8,7\). \(<6\): 5 → 1 discordant. \(>6\): 8,7 → 2 concordant.
\item \(i=6, Y_6=5\): later values \(8,7\). Both \(>5\) → 2 concordant.
\item \(i=7, Y_7=8\): later value \(7\). \(<8\) → 1 discordant.
\end{itemize}
Total concordant \(C = 6+6+4+4+2+2 = 24\).
Total discordant \(D = 1+1+1+1 = 4\).
Check: total pairs = \(\binom{8}{2}=28\), \(C+D=28\) ✓.
\(S = C-D = 20\).
\[
\tau = \frac{2S}{n(n-1)} = \frac{40}{8\times 7} = \frac{40}{56} = \frac{5}{7} \approx 0.7143.
\]
\end{exoresolu}

\courssub{Uses and Limitations of Kendall's Coefficient}

\begin{aretenir}[Uses of Kendall's \(\tau\)]
\begin{itemize}
\item \textbf{Small samples:} The exact distribution of \(\tau\) under the null hypothesis of independence is known for small \(n\), making it ideal for exact hypothesis tests.
\item \textbf{Ordinal data:} When data are only available as ranks (e.g., preference rankings, competition results), \(\tau\) is a natural choice.
\item \textbf{Agreement between judges/raters:} \(\tau\) directly measures the degree of concordance between two rankings, which is exactly what is needed in inter‑rater reliability studies.
\item \textbf{Robustness:} Because it uses only the sign of rank differences, \(\tau\) is less affected by outliers than Spearman's \(R_s\) or Pearson's \(r\).
\end{itemize}
\end{aretenir}

\begin{attention}[Limitations of Kendall's \(\tau\)]
\begin{itemize}
\item \textbf{Computational burden for large \(n\):} Counting concordant/discordant pairs requires \(O(n^2)\) comparisons. For \(n > 30\) it becomes tedious by hand (though software handles it easily).
\item \textbf{Ties complicate the count:} Ties in either variable produce pairs that are neither concordant nor discordant. The standard formula \(\tau = \frac{2(C-D)}{n(n-1)}\) assumes no ties. With ties, one must use \emph{Kendall's \(\tau_b\)} or \(\tau_c\) which adjust the denominator. This is beyond the GCE syllabus but you should be aware that ties reduce the maximum possible \(|\tau|\).
\item \textbf{Only measures monotonic agreement:} Like \(R_s\), \(\tau\) captures only monotonic relationships. A perfect U‑shaped relationship would give \(\tau \approx 0\).
\end{itemize}
\end{attention}

\courssub{Comparison of \(r\), \(R_s\) and Kendall's \(\tau\)}

\begin{center}
\begin{tabularx}{\linewidth}{|l|X|X|X|}
\hline
\rowcolor{popblueL}
\textbf{Feature} & \textbf{Pearson's \(r\)} & \textbf{Spearman's \(R_s\)} & \textbf{Kendall's \(\tau\)} \\ \hline
\textbf{Data type} & Quantitative (interval/ratio) & Quantitative or ordinal & Ordinal or quantitative \\ \hline
\textbf{Relationship measured} & Linear & Monotonic & Monotonic (concordance) \\ \hline
\textbf{Sensitivity to outliers} & High & Moderate (uses rank differences) & Low (uses only signs) \\ \hline
\textbf{Assumptions} & Bivariate normality, linearity, homoscedasticity & None (non‑parametric) & None (non‑parametric) \\ \hline
\textbf{Interpretation} & Strength of linear fit & Strength of monotonic trend & Probability concordance minus discordance \\ \hline
\textbf{Computation} & Simple formula & Simple if no ties; otherwise Pearson on ranks & Pair counting (\(O(n^2)\)) \\ \hline
\textbf{When to use} & Linear relationship, quantitative data, no outliers & Monotonic trend, ordinal data, outliers present & Small samples, ordinal data, judge agreement, many ties \\ \hline
\end{tabularx}
\end{center}

\courssub{Worked Example: Two Teachers Ranking Six Candidates}

\begin{exoresolu}[Comparing \(R_s\) and \(\tau\) on a Small Sample]
Two teachers, Mr. Ngwa and Ms. Fon, independently rank 6 students (A–F) based on their overall performance. The rankings are:
\[
\begin{array}{c|cccccc}
\text{Student} & A & B & C & D & E & F \\ \hline
\text{Mr. Ngwa} & 1 & 2 & 3 & 4 & 5 & 6 \\
\text{Ms. Fon} & 2 & 1 & 4 & 3 & 6 & 5
\end{array}
\]
\begin{enumerate}
\item Compute Spearman's rank correlation coefficient \(R_s\).
\item Compute Kendall's rank correlation coefficient \(\tau\).
\item Compare the two values and comment on the agreement between the teachers.
\end{enumerate}
\tcblower
\textbf{Solution.}

\noindent\textbf{1. Spearman's \(R_s\).}
There are no ties. Differences \(d_i = R_{\text{Ngwa}} - R_{\text{Fon}}\):
\[
\begin{array}{c|cccccc}
\text{Student} & A & B & C & D & E & F \\ \hline
R_N & 1 & 2 & 3 & 4 & 5 & 6 \\
R_F & 2 & 1 & 4 & 3 & 6 & 5 \\
d_i & -1 & 1 & -1 & 1 & -1 & 1 \\
d_i^2 & 1 & 1 & 1 & 1 & 1 & 1
\end{array}
\]
\(\sum d_i^2 = 6\). With \(n=6\):
\[
R_s = 1 - \frac{6 \times 6}{6(6^2-1)} = 1 - \frac{36}{6 \times 35} = 1 - \frac{36}{210} = 1 - 0.17143 = 0.82857.
\]

\noindent\textbf{2. Kendall's \(\tau\).}
Arrange in Mr. Ngwa's order (already natural). Ms. Fon's ranks in that order: \(Y = (2, 1, 4, 3, 6, 5)\).
Count concordant/discordant pairs:
\begin{itemize}
\item \(i=1, Y_1=2\): later \(1,4,3,6,5\) → discordant: 1 (value 1); concordant: 4 (4.3,6.5).
\item \(i=2, Y_2=1\): later \(4,3,6,5\) → all concordant (4).
\item \(i=3, Y_3=4\): later \(3,6,5\) → discordant: 1 (3); concordant: 2 (6,5).
\item \(i=4, Y_4=3\): later \(6,5\) → both concordant (2).
\item \(i=5, Y_5=6\): later \(5\) → discordant: 1 (5).
\end{itemize}
\(C = 4+4+2+2 = 12\), \(D = 1+1+1 = 3\). Total pairs = \(\binom{6}{2}=15\), \(C+D=15\) ✓.
\(S = C-D = 9\).
\[
\tau = \frac{2S}{n(n-1)} = \frac{18}{6 \times 5} = \frac{18}{30} = 0.6.
\]

\noindent\textbf{3. Comparison and Comment.}
Both coefficients indicate \textbf{strong positive agreement} between the two teachers. \(R_s \approx 0.83\) is higher than \(\tau = 0.60\). This is typical: Spearman's coefficient tends to be larger in magnitude than Kendall's for the same data because \(R_s\) uses the squared rank differences, while \(\tau\) uses only the sign. Kendall's \(\tau\) has a more direct probabilistic interpretation: the probability that a randomly chosen pair is concordant minus the probability it is discordant. Here, \(\tau = 0.6\) means concordant pairs exceed discordant pairs by 60\% of all pairs. The teachers' rankings are very similar; the only disagreements are the swaps of adjacent pairs (A↔B, C↔D, E↔F).
\end{exoresolu}

\begin{savaistu}[Historical Note]
Charles Spearman (1863–1945) developed his rank correlation while studying intelligence testing. Maurice Kendall (1907–1983) introduced \(\tau\) in 1938 as a measure of rank correlation with a simpler sampling distribution. Both are British statisticians. In Cameroon, rank correlation is often used in educational research to compare rankings of schools or students by different examiners, and in agricultural trials where farmers rank crop varieties by preference.
\end{savaistu}

\begin{exercice}[Practice: Rank Correlation]
\begin{enumerate}
\item Two coaches rank 7 athletes in a 100m sprint trial. The ranks are:
\[
\begin{array}{c|ccccccc}
\text{Athlete} & 1 & 2 & 3 & 4 & 5 & 6 & 7 \\ \hline
\text{Coach X} & 1 & 2 & 3 & 4 & 5 & 6 & 7 \\
\text{Coach Y} & 3 & 1 & 2 & 5 & 4 & 7 & 6
\end{array}
\]
Calculate (a) Spearman's \(R_s\), (b) Kendall's \(\tau\). Interpret.
\item In a beauty contest, 5 judges rank 8 contestants. Explain why Kendall's \(\tau\) might be preferred over Spearman's \(R_s\) for measuring agreement between two specific judges.
\item A researcher collects data on the ranking of 10 villages by \emph{access to healthcare} and \emph{literacy rate}. The data contain several tied ranks. Describe the correct procedure to compute Spearman's coefficient in this case.
\end{enumerate}
\end{exercice}

\begin{corrige}[Exercise 1]
(a) \(d_i = R_X - R_Y\): \((-2, 1, 1, -1, 1, -1, 1)\); \(\sum d_i^2 = 4+1+1+1+1+1+1 = 10\). \(R_s = 1 - \frac{6\times10}{7(49-1)} = 1 - \frac{60}{336} = 1 - 0.1786 = 0.8214\). \\
(b) \(Y = (3,1,2,5,4,7,6)\). Count pairs: \(C=15, D=6\) (details omitted). \(\tau = \frac{2(15-6)}{7\times6} = \frac{18}{42} = 0.4286\). \\
Interpretation: Strong positive agreement by Spearman; moderate agreement by Kendall. The coaches generally agree on the top athletes but differ on exact ordering.
\end{corrige}

\begin{corrige}[Exercise 2]
Kendall's \(\tau\) is based on pairwise concordance, which directly reflects whether two judges agree on the relative order of each pair of contestants. It is less affected by the magnitude of rank differences and handles small samples well. With 5 judges, one might compute \(\tau\) for each pair of judges to assess inter‑rater reliability.
\end{corrige}

\begin{corrige}[Exercise 3]
When ties are present, assign the average rank to each tied observation. Then compute the Pearson correlation coefficient on these averaged ranks (i.e., use the definition formula for \(R_s\) with the ranked data). Do \emph{not} use the simplified formula \(1 - \frac{6\sum d^2}{n(n^2-1)}\) because it assumes no ties.
\end{corrige}

\newpage

\coursec{Integration activity}

\coursec{Regression and Correlation}

\courssub{Introduction to Regression and Scatter Diagrams}

In many scientific, economic, and educational contexts, we observe two variables simultaneously and wish to understand how they relate. For instance, a teacher may suspect that the time a student spends revising influences their examination score. \emph{Regression analysis} provides tools to model the relationship between a dependent variable (response) and one or more independent variables (explanatory). \emph{Correlation} measures the strength and direction of the association.

\begin{definition}[Bivariate Data]
A set of paired observations $(x_i, y_i)$, $i = 1, 2, \dots, n$, where $x$ is the independent (explanatory) variable and $y$ is the dependent (response) variable.
\end{definition}

\begin{definition}[Scatter Diagram]
A graphical display of bivariate data where each pair $(x_i, y_i)$ is plotted as a point on a Cartesian plane with $x$ on the horizontal axis and $y$ on the vertical axis. It gives a visual impression of the relationship: linear, curved, clustered, or no pattern.
\end{definition}

\begin{exemplebox}[Scatter Diagram for Revision Hours vs Mock Scores]
Consider the data collected by Group A in a Yaoundé bilingual high school:

\begin{center}
\begin{tabularx}{\linewidth}{|c|X|X|}\hline
\rowcolor{popblueL} \textbf{Student} & \textbf{Hours of revision per week ($x$)} & \textbf{Mock score ($y$)} \\ \hline
1 & 2 & 42 \\ \hline
2 & 3 & 48 \\ \hline
3 & 4 & 55 \\ \hline
4 & 5 & 52 \\ \hline
5 & 6 & 65 \\ \hline
6 & 7 & 62 \\ \hline
7 & 8 & 75 \\ \hline
8 & 9 & 70 \\ \hline
9 & 10 & 85 \\ \hline
10 & 11 & 82 \\ \hline
11 & 12 & 90 \\ \hline
\end{tabularx}
\end{center}

\begin{popfigure}
\begin{tikzpicture}
\begin{axis}[
    width=0.9\linewidth,
    height=8cm,
    xlabel={Hours of revision per week ($x$)},
    ylabel={Mock score ($y$)},
    xmin=0, xmax=13,
    ymin=35, ymax=95,
    grid=major,
    grid style={popline},
    tick style={popdark},
    axis line style={popdark},
    scatter/classes={a={mark=*, popblue}},
]
\addplot[scatter, only marks, scatter src=explicit symbolic] coordinates {
    (2,42) [a]
    (3,48) [a]
    (4,55) [a]
    (5,52) [a]
    (6,65) [a]
    (7,62) [a]
    (8,75) [a]
    (9,70) [a]
    (10,85) [a]
    (11,82) [a]
    (12,90) [a]
};
\end{axis}
\end{tikzpicture}
\legende{Scatter diagram of mock scores against revision hours.}
\end{popfigure}

\textbf{Comment:} The points show a clear upward trend. As revision hours increase, mock scores tend to increase. The points are roughly clustered around a straight line, suggesting a strong positive linear relationship.
\end{exemplebox}

\courssub{Mayer's Method for the Regression Line}

Mayer's method (also called the \emph{double mean} or \emph{median-median} method) is a simple graphical technique to estimate a regression line by using the mean points of two halves of the data.

\begin{methode}[Mayer's Method -- Steps]
\begin{enumerate}
    \item Order the data by increasing $x$ (if not already ordered).
    \item Split the $n$ observations into two groups of equal size (or as equal as possible). If $n$ is odd, the middle observation is usually omitted or the groups are split as $\lfloor n/2 \rfloor$ and $\lceil n/2 \rceil$.
    \item Calculate the mean point $(\bar{x}_1, \bar{y}_1)$ of the first group and $(\bar{x}_2, \bar{y}_2)$ of the second group.
    \item Draw the line passing through these two mean points. Its equation is $y = a_M + b_M x$ where
    \[
    b_M = \frac{\bar{y}_2 - \bar{y}_1}{\bar{x}_2 - \bar{x}_1}, \qquad
    a_M = \bar{y}_1 - b_M \bar{x}_1 \quad (\text{or using } \bar{y}_2, \bar{x}_2).
    \]
\end{enumerate}
\end{methode}

\begin{attention}[Limitations of Mayer's Method]
Mayer's line depends on the arbitrary split of the data and does not use all points optimally. It is not the ``line of best fit'' in the least squares sense. It is mainly a pedagogical tool or a quick approximation.
\end{attention}

\courssub{Correlation, Sample Variance and Covariance}

To quantify the linear relationship, we need measures of spread and joint variability.

\begin{definition}[Sample Variance and Covariance]
For a sample $(x_i, y_i)$, $i=1,\dots,n$, with means $\bar{x}, \bar{y}$:
\begin{itemize}
    \item \textbf{Corrected sum of squares for $x$:} $S_{xx} = \sum (x_i - \bar{x})^2 = \sum x_i^2 - \frac{(\sum x_i)^2}{n}$.
    \item \textbf{Corrected sum of squares for $y$:} $S_{yy} = \sum (y_i - \bar{y})^2 = \sum y_i^2 - \frac{(\sum y_i)^2}{n}$.
    \item \textbf{Corrected sum of cross-products:} $S_{xy} = \sum (x_i - \bar{x})(y_i - \bar{y}) = \sum x_i y_i - \frac{\sum x_i \sum y_i}{n}$.
\end{itemize}
The \emph{sample covariance} is $\frac{S_{xy}}{n-1}$ (or $\frac{S_{xy}}{n}$ for the population analogue). The \emph{sample variances} are $\frac{S_{xx}}{n-1}$ and $\frac{S_{yy}}{n-1}$.
\end{definition}

\begin{propriete}[Properties of $S_{xx}, S_{yy}, S_{xy}$]
\begin{itemize}
    \item $S_{xx} \ge 0$, $S_{yy} \ge 0$; equality holds iff all $x$ (resp. $y$) are equal.
    \item $S_{xy}$ can be positive, negative, or zero.
    \item By the Cauchy-Schwarz inequality: $S_{xy}^2 \le S_{xx} S_{yy}$.
\end{itemize}
\end{propriete}

\courssub{Least Squares Regression Lines}

The \emph{least squares} criterion chooses the line that minimises the sum of squared vertical distances (residuals) between the observed $y_i$ and the fitted values $\hat{y}_i$.

\begin{propriete}[Least Squares Regression Line of $y$ on $x$]
The line $y = a + bx$ that minimises $\sum (y_i - a - b x_i)^2$ has slope and intercept:
\[
b = \frac{S_{xy}}{S_{xx}}, \qquad a = \bar{y} - b \bar{x}.
\]
The line always passes through the mean point $(\bar{x}, \bar{y})$.
\end{propriete}

\begin{propriete}[Least Squares Regression Line of $x$ on $y$]
If we instead minimise horizontal distances (treat $x$ as response), the line $x = a' + b' y$ has:
\[
b' = \frac{S_{xy}}{S_{yy}}, \qquad a' = \bar{x} - b' \bar{y}.
\]
This line also passes through $(\bar{x}, \bar{y})$. In general, the two lines are different unless $|r|=1$.
\end{propriete}

\begin{aretenir}[Interpretation of Slopes]
\begin{itemize}
    \item $b = \frac{S_{xy}}{S_{xx}}$: estimated change in $y$ for a one-unit increase in $x$.
    \item $b' = \frac{S_{xy}}{S_{yy}}$: estimated change in $x$ for a one-unit increase in $y$.
    \item $b \cdot b' = r^2$ (the coefficient of determination).
\end{itemize}
\end{aretenir}

\courssub{Residual Sum of Squares and the Line of Best Fit}

\begin{definition}[Residual]
For the regression line of $y$ on $x$, the residual for observation $i$ is $e_i = y_i - \hat{y}_i = y_i - (a + b x_i)$.
\end{definition}

\begin{propriete}[Residual Sum of Squares (RSS)]
\begin{itemize}
    \item For $y$ on $x$: $\text{RSS}_{y|x} = \sum e_i^2 = S_{yy} - b S_{xy} = S_{yy} - \frac{S_{xy}^2}{S_{xx}}$.
    \item For $x$ on $y$: $\text{RSS}_{x|y} = \sum (x_i - \hat{x}_i)^2 = S_{xx} - b' S_{xy} = S_{xx} - \frac{S_{xy}^2}{S_{yy}}$.
\end{itemize}
The least squares line is the \emph{line of best fit} because it minimises the respective RSS among all straight lines.
\end{propriete}

\begin{attention}[Residuals vs Errors]
Residuals are computed from the sample; errors are the unobservable true deviations. Residuals sum to zero ($\sum e_i = 0$) and are uncorrelated with $x$ ($\sum x_i e_i = 0$) for the least squares line.
\end{attention}

\courssub{The Product Moment Correlation Coefficient}

\begin{definition}[Pearson's Product Moment Correlation Coefficient ($r$)]
\[
r = \frac{S_{xy}}{\sqrt{S_{xx} S_{yy}}}.
\]
It measures the strength and direction of the \emph{linear} relationship between $x$ and $y$.
\end{definition}

\begin{propriete}[Properties of $r$]
\begin{itemize}
    \item $-1 \le r \le 1$.
    \item $r = 1$: perfect positive linear relationship (all points on a line with positive slope).
    \item $r = -1$: perfect negative linear relationship.
    \item $r = 0$: no linear relationship (but there could be a non-linear one).
    \item $r$ is invariant under linear changes of scale: $r(ax+b, cy+d) = \text{sign}(ac) \cdot r(x,y)$.
    \item $r^2 = \frac{\text{Explained variation}}{\text{Total variation}} = 1 - \frac{\text{RSS}}{S_{yy}}$ is the \emph{coefficient of determination}.
\end{itemize}
\end{propriete}

\begin{attention}[Correlation $\neq$ Causation]
A high $|r|$ does not imply that $x$ causes $y$. There may be confounding variables, reverse causality, or coincidence. In our example, revision hours may correlate with scores, but motivation, prior ability, and teaching quality also play roles.
\end{attention}

\courssub{Rank Correlation: Spearman and Kendall}

When data are not interval-scaled, contain outliers, or the relationship is monotonic but not linear, rank correlation coefficients are preferred.

\begin{definition}[Spearman's Rank Correlation Coefficient ($R_s$)]
Replace each $x_i$ and $y_i$ by their ranks $R(x_i), R(y_i)$ (average ranks for ties). Then
\[
R_s = 1 - \frac{6 \sum d_i^2}{n(n^2-1)}, \quad \text{where } d_i = R(x_i) - R(y_i),
\]
provided there are no ties. With ties, use the Pearson formula on the ranks.
\end{definition}

\begin{definition}[Kendall's Rank Correlation Coefficient ($\tau$)]
Consider all $\binom{n}{2}$ pairs of observations $(i,j)$ with $i<j$. A pair is \emph{concordant} if $(x_i - x_j)(y_i - y_j) > 0$ and \emph{discordant} if $< 0$. Let $C$ and $D$ be the counts. Then
\[
\tau = \frac{C - D}{C + D} = \frac{C - D}{\binom{n}{2}}.
\]
(If ties exist, a modified denominator is used: $\tau_b$ or $\tau_c$.)
\end{definition}

\begin{aretenir}[Comparison of Correlation Measures]
\begin{center}
\begin{tabularx}{\linewidth}{|l|X|X|X|}\hline
\rowcolor{popblueL} \textbf{Measure} & \textbf{Assesses} & \textbf{Assumptions} & \textbf{Robustness} \\ \hline
Pearson $r$ & Linear relationship & Bivariate normality, linearity, interval data & Sensitive to outliers \\ \hline
Spearman $R_s$ & Monotonic relationship & Ordinal or continuous data; ranks used & Less sensitive to outliers \\ \hline
Kendall $\tau$ & Concordance (pairwise agreement) & Ordinal data; ranks used & Most robust; smaller magnitude typically \\ \hline
\end{tabularx}
\end{center}
\end{aretenir}

\begin{attention}[Uses and Limitations of Kendall's $\tau$]
\begin{itemize}
    \item \textbf{Uses:} Small samples, ordinal data, many tied ranks, robust alternative to Spearman.
    \item \textbf{Limitations:} Computationally heavier for large $n$ (though $O(n \log n)$ algorithms exist); $\tau$ tends to be smaller in magnitude than $R_s$ for the same data; interpretation as ``probability of concordance minus probability of discordance'' is less intuitive for non-statisticians.
\end{itemize}
\end{attention}

\courssub{Integration Activity: Revision Hours and Mock Scores}

\courssub{Aims of the Activity}
\begin{itemize}
    \item Consolidate the understanding of regression and correlation techniques covered in the chapter.
    \item Apply scatter diagrams, Mayer's method, least squares regression, residual analysis, product moment correlation, Spearman's rank correlation, and Kendall's rank correlation to a real dataset.
    \item Develop skills in data collection, organisation, computation, interpretation, and report writing.
    \item Critically evaluate the reliability of predictions and compare different correlation measures.
    \item Foster collaborative work and communication of statistical findings in a Cameroonian educational context.
\end{itemize}

\courssub{Situation}
In many Cameroonian secondary schools, Upper Sixth students prepare intensely for the GCE Advanced Level Mathematics examination. A common belief among students and teachers is that the number of hours devoted to personal revision each week strongly influences the mock examination scores. To investigate this claim, the Mathematics teacher of a bilingual high school in Yaoundé proposes an integration activity. The class is divided into groups of four to five students. Each group must select 10–12 volunteer classmates (different from the group members) and record two variables for each volunteer:
\begin{itemize}
    \item $x$: average number of hours of personal revision per week (over the last month),
    \item $y$: score obtained in the recent Mathematics mock examination (out of 100).
\end{itemize}
The data collected by one group (Group A) is shown in the table below. This data will be used throughout the activity to illustrate the required computations. Each group will use their own data, but the steps are identical.

\begin{center}
\begin{tabularx}{\linewidth}{|c|X|X|}\hline
\rowcolor{popblueL} \textbf{Student} & \textbf{Hours of revision per week ($x$)} & \textbf{Mock score ($y$)} \\ \hline
1 & 2 & 42 \\ \hline
2 & 3 & 48 \\ \hline
3 & 4 & 55 \\ \hline
4 & 5 & 52 \\ \hline
5 & 6 & 65 \\ \hline
6 & 7 & 62 \\ \hline
7 & 8 & 75 \\ \hline
8 & 9 & 70 \\ \hline
9 & 10 & 85 \\ \hline
10 & 11 & 82 \\ \hline
11 & 12 & 90 \\ \hline
\end{tabularx}
\end{center}

The teacher asks each group to carry out a complete statistical analysis and write a short report. The activity culminates in a plenary session where groups compare their results and discuss the strengths and limitations of each method.

\courssub{Tasks}
\begin{enumerate}
    \item \textbf{Scatter diagram:} Plot the data on a scatter diagram with $x$ on the horizontal axis and $y$ on the vertical axis. Comment on the pattern observed.
    \item \textbf{Mayer's method:} 
    \begin{enumerate}
        \item Order the data by $x$ (already done). Split the 11 observations into two groups: the first 5 and the last 6 (since $n=11$ is odd).
        \item Calculate the mean point $(\bar{x}_1,\bar{y}_1)$ of the first group and $(\bar{x}_2,\bar{y}_2)$ of the second group.
        \item Draw the line passing through these two mean points on the scatter diagram. Find its equation in the form $y = a + bx$.
    \end{enumerate}
    \item \textbf{Least squares regression lines:}
    \begin{enumerate}
        \item Compute the corrected sums of squares and cross-products: $S_{xx}$, $S_{yy}$, $S_{xy}$.
        \item Determine the least squares estimates of the slope $b$ and intercept $a$ for the regression line $y = a + bx$ ($y$ on $x$).
        \item Determine the least squares estimates for the regression line $x = a' + b'y$ ($x$ on $y$).
        \item Write both equations and plot the $y$ on $x$ line on the scatter diagram.
    \end{enumerate}
    \item \textbf{Residual Sum of Squares (RSS) and Product Moment Correlation Coefficient:}
    \begin{enumerate}
        \item Calculate the Residual Sum of Squares (RSS) for the least squares line of $y$ on $x$ and for $x$ on $y$.
        \item Compute the product moment correlation coefficient $r$.
        \item Interpret the value of $r$ and $r^2$ in the context of the problem.
    \end{enumerate}
    \item \textbf{Rank correlation:}
    \begin{enumerate}
        \item Assign ranks to the $x$ values (already 1 to 11) and to the $y$ values (handle ties if any).
        \item Compute Spearman's rank correlation coefficient $R_s$.
        \item Compute Kendall's rank correlation coefficient $\tau$.
        \item Compare the three correlation coefficients ($r$, $R_s$, $\tau$) and comment on their similarities and differences.
    \end{enumerate}
    \item \textbf{Prediction and report:}
    \begin{enumerate}
        \item Use the least squares regression line of $y$ on $x$ to predict the mock score of a student who revises 9 hours per week.
        \item Discuss the reliability of this prediction, considering the value of $r$, the scatter diagram, the sample size, and whether 9 hours lies within the range of the data.
        \item Write a concise report (about 200 words) summarising the findings, the prediction, and a comparison of the three correlation measures. State any limitations of the study.
    \end{enumerate}
\end{enumerate}

\courssub{Model Answer for Group A's Data}
\begin{exoresolu}[Model Answer]
\tcblower
\textbf{1. Scatter Diagram}
The scatter diagram below shows the 11 data points. There is a clear upward trend: as revision hours increase, mock scores tend to increase. The points are roughly clustered around a straight line, suggesting a strong positive linear relationship.

\begin{popfigure}
\begin{tikzpicture}
\begin{axis}[
    width=0.9\linewidth,
    height=8cm,
    xlabel={Hours of revision per week ($x$)},
    ylabel={Mock score ($y$)},
    xmin=0, xmax=13,
    ymin=35, ymax=95,
    grid=major,
    grid style={popline},
    tick style={popdark},
    axis line style={popdark},
    legend style={at={(0.02,0.98)}, anchor=north west, draw=popdark, fill=white},
    scatter/classes={a={mark=*, popblue}},
]
\addplot[scatter, only marks, scatter src=explicit symbolic] coordinates {
    (2,42) [a]
    (3,48) [a]
    (4,55) [a]
    (5,52) [a]
    (6,65) [a]
    (7,62) [a]
    (8,75) [a]
    (9,70) [a]
    (10,85) [a]
    (11,82) [a]
    (12,90) [a]
};
\addlegendentry{Data points}
% Least squares line y on x: y = 33.418 + 4.655x
\addplot[domain=0:13, samples=2, poporange, thick] {33.418 + 4.655*x};
\addlegendentry{LS line $y$ on $x$}
% Mayer's line: y = 34.267 + 4.533x
\addplot[domain=0:13, samples=2, popgreen, thick, dashed] {34.267 + 4.533*x};
\addlegendentry{Mayer's line}
% LS line x on y: x = -5.236 + 0.203y  ->  y = 25.79 + 4.926x (rearranged for plotting)
\addplot[domain=0:13, samples=2, poppurple, thick, dotted] {25.79 + 4.926*x};
\addlegendentry{LS line $x$ on $y$ (rearranged)}
\end{axis}
\end{tikzpicture}
\legende{Scatter diagram with least squares regression lines and Mayer's line.}
\end{popfigure}

\textbf{2. Mayer's Method}
The data is already ordered by $x$. Split into two groups:
\begin{itemize}
    \item Group 1 (first 5): $(2,42), (3,48), (4,55), (5,52), (6,65)$
    \item Group 2 (last 6): $(7,62), (8,75), (9,70), (10,85), (11,82), (12,90)$
\end{itemize}
Mean points:
\[
\bar{x}_1 = \frac{2+3+4+5+6}{5} = 4,\quad \bar{y}_1 = \frac{42+48+55+52+65}{5} = \frac{262}{5} = 52.4
\]
\[
\bar{x}_2 = \frac{7+8+9+10+11+12}{6} = \frac{57}{6} = 9.5,\quad \bar{y}_2 = \frac{62+75+70+85+82+90}{6} = \frac{464}{6} \approx 77.333
\]
Slope of Mayer's line:
\[
b_M = \frac{\bar{y}_2 - \bar{y}_1}{\bar{x}_2 - \bar{x}_1} = \frac{77.333 - 52.4}{9.5 - 4} = \frac{24.933}{5.5} \approx 4.533
\]
Intercept:
\[
a_M = \bar{y}_1 - b_M \bar{x}_1 = 52.4 - 4.533 \times 4 = 52.4 - 18.132 = 34.268
\]
Equation of Mayer's line: $\boxed{y = 34.27 + 4.53x}$ (rounded to two decimals).

\textbf{3. Least Squares Regression Lines}
We compute the necessary sums:
\[
\begin{aligned}
n &= 11 \\
\sum x &= 77,\quad \sum y = 726 \\
\sum x^2 &= 649,\quad \sum y^2 = 50440 \\
\sum xy &= 5594
\end{aligned}
\]
Means:
\[
\bar{x} = \frac{77}{11} = 7,\quad \bar{y} = \frac{726}{11} = 66
\]
Corrected sums:
\[
\begin{aligned}
S_{xx} &= \sum x^2 - \frac{(\sum x)^2}{n} = 649 - \frac{5929}{11} = 649 - 539 = 110 \\
S_{yy} &= \sum y^2 - \frac{(\sum y)^2}{n} = 50440 - \frac{527076}{11} = 50440 - 47916 = 2524 \\
S_{xy} &= \sum xy - \frac{\sum x \sum y}{n} = 5594 - \frac{77 \times 726}{11} = 5594 - 5082 = 512
\end{aligned}
\]

\textbf{Line of $y$ on $x$:}
\[
b = \frac{S_{xy}}{S_{xx}} = \frac{512}{110} = \frac{256}{55} \approx 4.6545
\]
\[
a = \bar{y} - b\bar{x} = 66 - \frac{256}{55} \times 7 = 66 - \frac{1792}{55} = \frac{3630 - 1792}{55} = \frac{1838}{55} \approx 33.4182
\]
Equation: $\boxed{y = 33.42 + 4.65x}$ (rounded to two decimals).

\textbf{Line of $x$ on $y$:}
\[
b' = \frac{S_{xy}}{S_{yy}} = \frac{512}{2524} = \frac{128}{631} \approx 0.2029
\]
\[
a' = \bar{x} - b'\bar{y} = 7 - \frac{128}{631} \times 66 = 7 - \frac{8448}{631} = \frac{4417 - 8448}{631} = -\frac{4031}{631} \approx -6.388
\]
Equation: $\boxed{x = -6.39 + 0.203y}$ (or $y = 31.48 + 4.93x$ when rearranged for plotting).

Both lines pass through $(\bar{x}, \bar{y}) = (7, 66)$. The $y$ on $x$ line is plotted in solid orange; the $x$ on $y$ line (rearranged) in dotted purple. They differ because they minimise different distances (vertical vs horizontal).

\textbf{4. Residual Sum of Squares (RSS) and Product Moment Correlation Coefficient}
\[
\text{RSS}_{y|x} = S_{yy} - b S_{xy} = 2524 - \frac{512}{110} \times 512 = 2524 - \frac{262144}{110} = 2524 - 2383.127 = 140.873
\]
Alternatively, $\text{RSS}_{y|x} = S_{yy} - \frac{S_{xy}^2}{S_{xx}} = 2524 - \frac{512^2}{110} = 140.873$.

\[
\text{RSS}_{x|y} = S_{xx} - b' S_{xy} = 110 - \frac{512}{2524} \times 512 = 110 - \frac{262144}{2524} = 110 - 103.861 = 6.139
\]

Product moment correlation coefficient:
\[
r = \frac{S_{xy}}{\sqrt{S_{xx} S_{yy}}} = \frac{512}{\sqrt{110 \times 2524}} = \frac{512}{\sqrt{277640}} = \frac{512}{526.915} \approx 0.9717
\]
\[
r^2 = (0.9717)^2 \approx 0.9442
\]

\textbf{Interpretation:} $r = 0.972$ indicates a very strong positive linear correlation between hours of revision and mock scores. About $94.4\%$ of the variation in mock scores is explained by the linear relationship with revision hours (coefficient of determination $r^2$). The remaining $5.6\%$ is due to other factors or random variation.

\textbf{5. Rank Correlation}
Ranks of $x$ are simply $1,2,\dots,11$ because the $x$ values are distinct and already sorted.
Ranks of $y$: sort the $y$ values and assign ranks 1 to 11.
\[
\begin{array}{c|ccccccccccc}
\text{Student} & 1 & 2 & 3 & 4 & 5 & 6 & 7 & 8 & 9 & 10 & 11 \\
\hline
y & 42 & 48 & 55 & 52 & 65 & 62 & 75 & 70 & 85 & 82 & 90 \\
\text{Rank of } y & 1 & 2 & 4 & 3 & 6 & 5 & 8 & 7 & 10 & 9 & 11
\end{array}
\]
Differences $d = \text{rank}_x - \text{rank}_y$:
\[
d: 0, 0, -1, 1, -1, 1, -1, 1, -1, 1, 0
\]
\[
\sum d^2 = 0+0+1+1+1+1+1+1+1+1+0 = 8
\]
Spearman's rank correlation coefficient (no ties):
\[
R_s = 1 - \frac{6\sum d^2}{n(n^2-1)} = 1 - \frac{6 \times 8}{11 \times 120} = 1 - \frac{48}{1320} = 1 - 0.03636 = 0.9636
\]

Kendall's $\tau$: Count concordant ($C$) and discordant ($D$) pairs. Since $x$ ranks are in natural order, discordant pairs correspond to inversions in the sequence of $y$ ranks: $[1,2,4,3,6,5,8,7,10,9,11]$. Inversions: $(4,3), (6,5), (8,7), (10,9)$ $\to$ 4 discordant pairs. Total pairs = $\binom{11}{2} = 55$. So $C = 51, D = 4$.
\[
\tau = \frac{C - D}{C + D} = \frac{51 - 4}{55} = \frac{47}{55} \approx 0.8545
\]

\textbf{Comparison:}
\begin{itemize}
    \item $r = 0.972$ (Pearson) measures linear correlation; assumes bivariate normality and linearity.
    \item $R_s = 0.964$ (Spearman) measures monotonic correlation; based on ranks, less sensitive to outliers.
    \item $\tau = 0.855$ (Kendall) measures concordance; more robust but generally smaller in magnitude.
\end{itemize}
All three indicate a very strong positive association. $R_s$ is closer to $r$ because the relationship is nearly linear. $\tau$ is lower because it counts pairwise concordances, giving less weight to the overall trend.

\textbf{6. Prediction and Report}
Prediction for $x = 9$ hours using the least squares line of $y$ on $x$:
\[
\hat{y} = 33.418 + 4.655 \times 9 = 33.418 + 41.895 = 75.313 \approx 75.3
\]
A student revising 9 hours per week is predicted to score about 75 out of 100 in the mock examination.

\textbf{Reliability Discussion:}
\begin{itemize}
    \item The correlation is very strong ($r = 0.972$), and the scatter diagram shows points tightly clustered around the line.
    \item The prediction is within the range of the data ($x=9$ is between 2 and 12), so it is an interpolation, not an extrapolation.
    \item However, the sample size is only 11, which is small. The RSS is about 141, giving a residual standard deviation of $s = \sqrt{\text{RSS}/(n-2)} \approx \sqrt{141/9} \approx 3.95$. A rough 95\% prediction interval would be $\hat{y} \pm 2 \times 3.95 \approx 75.3 \pm 7.9$, i.e., between 67 and 83.
    \item The relationship might not be perfectly linear for extreme values, and other factors (teaching quality, prior knowledge, health) also affect scores.
\end{itemize}
Overall, the prediction is reasonably reliable but should be treated as an estimate with a margin of error.

\textbf{Report}
\begin{quote}
\textit{Our investigation on 11 Upper Sixth students reveals a very strong positive relationship between weekly revision hours and Mathematics mock scores. The scatter diagram shows a clear linear trend. Mayer's method gave the line $y = 34.27 + 4.53x$, while the least squares line of $y$ on $x$ is $y = 33.42 + 4.65x$ and the line of $x$ on $y$ is $x = -6.39 + 0.203y$. The product moment correlation coefficient is $r = 0.972$, indicating that 94.4\% of the variation in scores is explained by revision hours. Rank correlations confirm the strong association: Spearman's $R_s = 0.964$ and Kendall's $\tau = 0.855$. Using the least squares line, a student revising 9 hours per week is predicted to score about 75. This prediction is fairly reliable because the correlation is strong and 9 hours lies within our data range. However, the small sample size (11) limits the precision; a larger study across multiple schools would improve generalisability. All three correlation measures agree on a strong positive link, though Kendall's $\tau$ is more conservative. We recommend that students aim for at least 9 hours of weekly revision to target scores above 70, but they should also focus on effective study techniques.}
\end{quote}
\end{exoresolu}

\courssub{Exercises}

\begin{exercice}[Exercise 1: Least Squares Calculations]
For the data set: $(1,2), (2,3), (3,5), (4,4), (5,6)$.
\begin{enumerate}
    \item Compute $\bar{x}, \bar{y}, S_{xx}, S_{yy}, S_{xy}$.
    \item Find the least squares regression line of $y$ on $x$ and of $x$ on $y$.
    \item Calculate $r$, $r^2$, $\text{RSS}_{y|x}$, $\text{RSS}_{x|y}$.
    \item Plot the data and both lines. Comment on their difference.
\end{enumerate}
\end{exercice}

\begin{exercice}[Exercise 2: Rank Correlation with Ties]
The following data show the ranks of 8 students in Mathematics ($x$) and Physics ($y$):
\[
\begin{array}{c|cccccccc}
\text{Student} & A & B & C & D & E & F & G & H \\
\hline
\text{Math rank} & 1 & 2 & 3 & 4 & 5 & 6 & 7 & 8 \\
\text{Physics rank} & 2 & 1 & 4 & 3 & 6 & 5 & 8 & 7
\end{array}
\]
\begin{enumerate}
    \item Compute Spearman's $R_s$ using the $d^2$ formula.
    \item Compute Kendall's $\tau$ by counting concordant/discordant pairs.
    \item If the Physics ranks had a tie (e.g., students C and D both ranked 3.5), how would you compute $R_s$?
\end{enumerate}
\end{exercice}

\begin{exercice}[Exercise 3: Interpretation and Limitations]
A study finds $r = 0.85$ between the number of ice creams sold daily and the number of drowning incidents at a beach.
\begin{enumerate}
    \item Does this mean eating ice cream causes drowning? Explain.
    \item Suggest a confounding variable.
    \item Why might Spearman's $R_s$ be preferred over Pearson's $r$ if the relationship is curved but monotonic?
\end{enumerate}
\end{exercice}

\begin{corrige}[Exercise 1]
\begin{enumerate}
    \item $\sum x = 15$, $\sum y = 20$, $\sum x^2 = 55$, $\sum y^2 = 90$, $\sum xy = 69$, $n=5$.
    $\bar{x}=3$, $\bar{y}=4$.
    $S_{xx} = 55 - 225/5 = 10$, $S_{yy} = 90 - 400/5 = 10$, $S_{xy} = 69 - 300/5 = 9$.
    \item $b = 9/10 = 0.9$, $a = 4 - 0.9\times3 = 1.3$ $\to$ $y = 1.3 + 0.9x$.
    $b' = 9/10 = 0.9$, $a' = 3 - 0.9\times4 = -0.6$ $\to$ $x = -0.6 + 0.9y$.
    \item $r = 9/\sqrt{100} = 0.9$, $r^2 = 0.81$.
    $\text{RSS}_{y|x} = 10 - 0.9\times9 = 1.9$.
    $\text{RSS}_{x|y} = 10 - 0.9\times9 = 1.9$ (symmetric because $S_{xx}=S_{yy}$).
    \item The lines cross at $(3,4)$. They differ because one minimises vertical, the other horizontal distances.
\end{enumerate}
\end{corrige}

\begin{corrige}[Exercise 2]
\begin{enumerate}
    \item $d$: $-1, 1, -1, 1, -1, 1, -1, 1$ $\to$ $\sum d^2 = 8$.
    $R_s = 1 - \frac{6\times8}{8\times63} = 1 - \frac{48}{504} = 1 - 0.0952 = 0.9048$.
    \item Physics ranks in Math order: $[2,1,4,3,6,5,8,7]$. Inversions: $(2,1), (4,3), (6,5), (8,7)$ $\to$ $D=4$. Total pairs = 28. $C=24$. $\tau = (24-4)/28 = 20/28 = 0.714$.
    \item With ties, use the Pearson formula on the ranks (or the adjusted formula $R_s = 1 - \frac{6\sum d^2 + \frac{1}{2}\sum(t^3-t)}{n(n^2-1)}$ where $t$ are tie sizes).
\end{enumerate}
\end{corrige}

\begin{corrige}[Exercise 3]
\begin{enumerate}
    \item No. Correlation does not imply causation. It is a spurious correlation.
    \item Temperature (or season): hot weather increases both ice cream sales and swimming (hence drowning risk).
    \item Pearson's $r$ measures only linear association; a curved monotonic relationship would give $|r| < 1$ even if perfectly monotonic. Spearman's $R_s$ uses ranks and captures any monotonic trend.
\end{enumerate}
\end{corrige}

\newpage

\coursec{Methods and worked examples}

\courssub{Introduction to Regression and Scatter Diagrams}

\begin{definition}[Bivariate Data and Variables]
When we collect data on two variables for each individual or observation, we have \cle{bivariate data}. The variable we think might explain or influence the other is called the \cle{independent variable} (or \cle{explanatory variable}), usually denoted $x$. The variable we think might respond to changes in $x$ is the \cle{dependent variable} (or \cle{response variable}), usually denoted $y$. In the GCE Advanced Level examination, $x$ is always plotted on the horizontal axis and $y$ on the vertical axis.
\end{definition}

\begin{methode}[Drawing and interpreting a scatter diagram]
\begin{enumerate}
\item \textbf{Identify the variables:} Decide which variable is $x$ (independent) and which is $y$ (dependent). Label axes clearly with variable names and units.
\item \textbf{Choose appropriate scales:} Select scales so that the data points occupy most of the graph paper. The origin $(0,0)$ need not be included unless it is meaningful for the context. Use a broken axis symbol ($\sim$) if the scale does not start at zero.
\item \textbf{Plot the points:} For each pair $(x_i, y_i)$, mark a point using a sharp pencil and small crosses ($\times$) or dots ($\bullet$). Do not join the points.
\item \textbf{Calculate and plot the mean point:} Compute $\bar{x} = \frac{1}{n}\sum x_i$ and $\bar{y} = \frac{1}{n}\sum y_i$. Plot $(\bar{x}, \bar{y})$ clearly (e.g., with a ringed dot $\circledcirc$). \textbf{Key reflex:} Any reasonable regression line must pass through $(\bar{x}, \bar{y})$.
\item \textbf{Observe the pattern:} Look for a general trend:
\begin{itemize}
\item \cle{Positive linear correlation}: points tend to rise from left to right, roughly following a straight line.
\item \cle{Negative linear correlation}: points tend to fall from left to right, roughly following a straight line.
\item \cle{No linear correlation}: points appear randomly scattered with no discernible linear trend.
\item \cle{Non-linear pattern}: points follow a curve (e.g., quadratic, exponential). In this case, linear regression is not appropriate unless the data is transformed.
\end{itemize}
\item \textbf{Draw a line of best fit by eye (optional):} If the pattern appears roughly linear, draw a straight line that passes as close as possible to all points, with roughly equal numbers of points above and below the line, and passing through $(\bar{x}, \bar{y})$.
\item \textbf{Interpret the line:} The slope (gradient) indicates the rate of change of $y$ with respect to $x$; the intercept gives the value of $y$ when $x=0$ (only meaningful if $x=0$ is within or near the data range).
\end{enumerate}
\end{methode}

\begin{attention}[Common mistakes with scatter diagrams]
\begin{itemize}
\item Joining the points with line segments (this is a time-series plot, not a scatter diagram).
\item Choosing scales that squash all points into a tiny corner.
\item Forgetting to label axes with variable names and units.
\item Drawing a line of best fit that does not pass through $(\bar{x}, \bar{y})$.
\item Interpreting the $y$-intercept when $x=0$ is far outside the data range (extrapolation).
\end{itemize}
\end{attention}

\begin{exoresolu}[Scatter diagram for height and weight of students]
The table below shows the heights (in cm) and weights (in kg) of 10 Upper Sixth students in a college in Yaoundé.

\begin{center}
\begin{tabular}{|c|c|c|c|c|c|c|c|c|c|c|}
\hline
Height $x$ & 160 & 165 & 170 & 175 & 180 & 185 & 190 & 195 & 200 & 205 \\
\hline
Weight $y$ & 55  & 58  & 62  & 65  & 70  & 73  & 78  & 82  & 85  & 90  \\
\hline
\end{tabular}
\end{center}

\begin{enumerate}
\item Draw a scatter diagram for this data.
\item Comment on the correlation.
\item Draw a line of best fit by eye and estimate the weight of a student of height 178 cm.
\end{enumerate}
\tcblower
\textbf{Solution.}

\begin{enumerate}
\item \textbf{Scatter diagram:} We choose the horizontal axis for height (160–205 cm) and vertical axis for weight (55–90 kg). Plot each pair $(x_i, y_i)$. The points are:
$(160,55), (165,58), (170,62), (175,65), (180,70), (185,73), (190,78), (195,82), (200,85), (205,90)$.

\begin{popfigure}
\begin{tikzpicture}[scale=0.55]
\draw[->] (0,0) -- (10.5,0) node[right] {Height (cm)};
\draw[->] (0,0) -- (0,9.5) node[above] {Weight (kg)};
% x-axis ticks: 160 to 205 step 5 -> 10 units for 45 cm, 1 unit = 4.5 cm? Let's map 160->1, 205->10
\foreach \i/\lbl in {1/160,2/165,3/170,4/175,5/180,6/185,7/190,8/195,9/200,10/205}
  \draw (\i,0.1) -- (\i,-0.1) node[below,font=\small] {\lbl};
% y-axis ticks: 55 to 90 step 5 -> 8 units for 35 kg, 1 unit = 4.375 kg? Map 55->1, 90->8
\foreach \j/\lbl in {1/55,2/60,3/65,4/70,5/75,6/80,7/85,8/90}
  \draw (0.1,\j) -- (-0.1,\j) node[left,font=\small] {\lbl};
% Points: x index = (x-160)/5 + 1, y index = (y-55)/5 + 1
% (160,55)->(1,1); (165,58)->(2,1.6); (170,62)->(3,2.4); (175,65)->(4,3); (180,70)->(5,4);
% (185,73)->(6,4.6); (190,78)->(7,5.6); (195,82)->(8,6.4); (200,85)->(9,7); (205,90)->(10,8)
\foreach \x/\y in {1/1, 2/1.6, 3/2.4, 4/3, 5/4, 6/4.6, 7/5.6, 8/6.4, 9/7, 10/8}
  \draw[fill=popblue] (\x,\y) circle (0.1);
% Mean point: xbar=182.5 -> (182.5-160)/5+1 = 5.5; ybar=71.8 -> (71.8-55)/5+1 = 4.36
\draw[fill=popred] (5.5,4.36) circle (0.12) node[above right,font=\small] {$(\bar{x},\bar{y})$};
\end{tikzpicture}
\legende{Scatter diagram of height vs weight for 10 students}
\end{popfigure}

\item \textbf{Comment:} The points lie very close to a straight line sloping upwards. There is a strong \cle{positive linear correlation} between height and weight. As height increases, weight tends to increase in an approximately linear fashion.

\item \textbf{Line of best fit by eye:} The mean height $\bar{x} = \frac{160+165+\cdots+205}{10} = 182.5$ cm. The mean weight $\bar{y} = \frac{55+58+\cdots+90}{10} = 71.8$ kg. The mean point is $(182.5, 71.8)$. Draw a straight line through this point that follows the trend of the points. Using the graph, at height 178 cm (which corresponds to $x=178$), the line gives a weight of approximately 67 kg. (Reading from the graph: between 175 cm (65 kg) and 180 cm (70 kg), 178 cm is about 67 kg.)
\end{enumerate}
\end{exoresolu}

\courssub{Mayer's Method for the Regression Line (Double Mean Method)}

\begin{definition}[Mayer's Method]
\cle{Mayer's method} (also called the \cle{double mean method}) is a simple arithmetic technique to find an approximate regression line when the scatter diagram suggests a linear relationship. It uses the mean points of two halves of the data ordered by $x$. It is less accurate than the least squares method but is useful for quick estimation and is examinable in the GCE Advanced Level syllabus.
\end{definition}

\begin{methode}[Mayer's method (step by step)]
\begin{enumerate}
\item \textbf{Order the data} by the independent variable $x$ in ascending order.
\item \textbf{Split the data} into two groups of equal size.
\begin{itemize}
\item If $n$ is even: first $n/2$ points form Group 1 (lower $x$), last $n/2$ points form Group 2 (higher $x$).
\item If $n$ is odd: the usual GCE approach is to omit the middle observation, or include it in both groups. We will adopt: omit the middle observation, so Group 1 has the first $(n-1)/2$ points, Group 2 has the last $(n-1)/2$ points.
\end{itemize}
\item \textbf{Compute the mean point of each group:} $(\bar{x}_1, \bar{y}_1)$ for Group 1 and $(\bar{x}_2, \bar{y}_2)$ for Group 2.
\item \textbf{Calculate the gradient:} $m = \dfrac{\bar{y}_2 - \bar{y}_1}{\bar{x}_2 - \bar{x}_1}$. (If $\bar{x}_2 = \bar{x}_1$, the method fails — vertical line.)
\item \textbf{Find the equation:} Using point-slope form: $y - \bar{y}_1 = m(x - \bar{x}_1)$, or $y = mx + c$ with $c = \bar{y}_1 - m\bar{x}_1$.
\item \textbf{Verify:} The line should pass through (or very close to) the overall mean point $(\bar{x}, \bar{y})$.
\end{enumerate}
\textbf{Key formula:} $m = \dfrac{\bar{y}_2 - \bar{y}_1}{\bar{x}_2 - \bar{x}_1}$.
\end{methode}

\begin{attention}[When Mayer's method is inappropriate]
\begin{itemize}
\item If the scatter diagram shows a clear non-linear pattern.
\item If there are extreme outliers that distort the group means.
\item If the two group mean points are very close together (small denominator), making the gradient unstable.
\item For precise work, the least squares method (next section) is always preferred.
\end{itemize}
\end{attention}

\begin{exoresolu}[Mayer's method for advertising and sales]
A company records its monthly advertising expenditure (in thousands of CFA francs) and the corresponding sales (in millions of CFA francs) for 8 months:

\begin{center}
\begin{tabular}{|c|c|c|c|c|c|c|c|c|}
\hline
Advertising $x$ & 2 & 3 & 4 & 5 & 6 & 7 & 8 & 9 \\
\hline
Sales $y$       & 10 & 12 & 14 & 15 & 18 & 20 & 22 & 24 \\
\hline
\end{tabular}
\end{center}

Use Mayer's method to find the regression line of sales on advertising expenditure. Hence estimate the sales when advertising expenditure is 6.5 thousand CFA francs.
\tcblower
\textbf{Solution.}

\begin{enumerate}
\item The data is already ordered by $x$. $n=8$ (even), so we split into two groups of 4.
Group 1 (first 4): $(2,10), (3,12), (4,14), (5,15)$.
Group 2 (last 4): $(6,18), (7,20), (8,22), (9,24)$.

\item \textbf{Group means:}
Group 1: $\bar{x}_1 = \frac{2+3+4+5}{4} = \frac{14}{4} = 3.5$,
$\bar{y}_1 = \frac{10+12+14+15}{4} = \frac{51}{4} = 12.75$.
Group 2: $\bar{x}_2 = \frac{6+7+8+9}{4} = \frac{30}{4} = 7.5$,
$\bar{y}_2 = \frac{18+20+22+24}{4} = \frac{84}{4} = 21$.

\item \textbf{Gradient:}
$m = \dfrac{\bar{y}_2 - \bar{y}_1}{\bar{x}_2 - \bar{x}_1} = \dfrac{21 - 12.75}{7.5 - 3.5} = \dfrac{8.25}{4} = 2.0625$.

\item \textbf{Equation:} Using point $(\bar{x}_1, \bar{y}_1) = (3.5, 12.75)$:
$y - 12.75 = 2.0625(x - 3.5)$
$y = 2.0625x - 7.21875 + 12.75 = 2.0625x + 5.53125$.
So the Mayer regression line is $y = 2.0625x + 5.53125$.
(Rounded to 3 significant figures: $y = 2.06x + 5.53$.)

\item \textbf{Check overall mean:}
$\bar{x} = \frac{2+3+\cdots+9}{8} = \frac{44}{8} = 5.5$,
$\bar{y} = \frac{10+12+\cdots+24}{8} = \frac{135}{8} = 16.875$.
Plug $x=5.5$ into line: $y = 2.0625(5.5) + 5.53125 = 11.34375 + 5.53125 = 16.875$. Exactly the overall mean. Good.

\item \textbf{Estimate for $x=6.5$:}
$y = 2.0625(6.5) + 5.53125 = 13.40625 + 5.53125 = 18.9375$ million CFA francs.
(Rounded: 18.9 million CFA francs.)
\end{enumerate}
\end{exoresolu}

\courssub{Sample Variance, Covariance and Product Moment Correlation Coefficient}

\begin{definition}[Corrected Sums of Squares and Products]
For a bivariate sample $(x_i, y_i), i=1,\ldots,n$, with means $\bar{x}, \bar{y}$, we define the \cle{corrected sums of squares and products}:
\begin{align*}
S_{xx} &= \sum_{i=1}^n (x_i - \bar{x})^2 = \sum x_i^2 - \frac{(\sum x_i)^2}{n}, \\
S_{yy} &= \sum_{i=1}^n (y_i - \bar{y})^2 = \sum y_i^2 - \frac{(\sum y_i)^2}{n}, \\
S_{xy} &= \sum_{i=1}^n (x_i - \bar{x})(y_i - \bar{y}) = \sum x_i y_i - \frac{(\sum x_i)(\sum y_i)}{n}.
\end{align*}
These quantities measure the spread of $x$, spread of $y$, and joint variability of $x$ and $y$ respectively. They are the building blocks for all regression and correlation calculations.
\end{definition}

\begin{definition}[Sample Variance and Covariance]
The \cle{sample variances} and \cle{sample covariance} are:
\begin{align*}
s_x^2 &= \frac{S_{xx}}{n-1}, \quad s_y^2 = \frac{S_{yy}}{n-1}, \quad s_{xy} = \frac{S_{xy}}{n-1}.
\end{align*}
The denominator $n-1$ makes them unbiased estimators of the population variance and covariance.
\end{definition}

\begin{definition}[Product Moment Correlation Coefficient (Pearson's $r$)]
The \cle{product moment correlation coefficient} (Pearson's $r$) measures the strength and direction of the \cle{linear} relationship between $x$ and $y$:
\[
r = \frac{S_{xy}}{\sqrt{S_{xx} S_{yy}}} = \frac{s_{xy}}{s_x s_y}.
\]
Properties:
\begin{itemize}
\item $-1 \le r \le 1$ always.
\item $r = 1$: perfect positive linear correlation (all points lie exactly on a line with positive slope).
\item $r = -1$: perfect negative linear correlation (all points lie exactly on a line with negative slope).
\item $r = 0$: no linear correlation (but there could be a non-linear relationship).
\item $r$ is invariant under linear changes of scale: $r(ax+b, cy+d) = \text{sign}(ac) \cdot r(x,y)$.
\end{itemize}
\end{definition}

\begin{methode}[Calculating $S_{xx}, S_{yy}, S_{xy}$, variances, covariance and $r$]
\begin{enumerate}
\item Compute the sums: $\sum x_i$, $\sum y_i$, $\sum x_i^2$, $\sum y_i^2$, $\sum x_i y_i$.
\item Compute the corrected sums using the shortcut formulas:
$S_{xx} = \sum x^2 - \frac{(\sum x)^2}{n}$,
$S_{yy} = \sum y^2 - \frac{(\sum y)^2}{n}$,
$S_{xy} = \sum xy - \frac{(\sum x)(\sum y)}{n}$.
\item Sample variances: $s_x^2 = \frac{S_{xx}}{n-1}$, $s_y^2 = \frac{S_{yy}}{n-1}$.
\item Sample covariance: $s_{xy} = \frac{S_{xy}}{n-1}$.
\item Correlation coefficient: $r = \frac{S_{xy}}{\sqrt{S_{xx} S_{yy}}}$.
\item Interpret $r$ in context.
\end{enumerate}
\textbf{Reflex:} Always compute $S_{xx}, S_{yy}, S_{xy}$ first; they are used in regression lines as well. Keep extra decimal places during calculation; round only the final answer.
\end{methode}

\begin{exoresolu}[Variance, covariance and correlation for a small dataset]
For the following data:
\begin{center}
\begin{tabular}{|c|c|c|c|c|c|}
\hline
$x$ & 2 & 4 & 6 & 8 & 10 \\
\hline
$y$ & 3 & 5 & 7 & 9 & 11 \\
\hline
\end{tabular}
\end{center}
Calculate:
\begin{enumerate}
\item $S_{xx}, S_{yy}, S_{xy}$.
\item The sample variances $s_x^2, s_y^2$ and covariance $s_{xy}$.
\item The product moment correlation coefficient $r$.
\item Comment on the result.
\end{enumerate}
\tcblower
\textbf{Solution.}

$n=5$.

\begin{enumerate}
\item \textbf{Sums:}
$\sum x = 2+4+6+8+10 = 30$, $\sum y = 3+5+7+9+11 = 35$.
$\sum x^2 = 4+16+36+64+100 = 220$.
$\sum y^2 = 9+25+49+81+121 = 285$.
$\sum xy = 2\cdot3 + 4\cdot5 + 6\cdot7 + 8\cdot9 + 10\cdot11 = 6+20+42+72+110 = 250$.

\item \textbf{Corrected sums:}
$S_{xx} = \sum x^2 - \frac{(\sum x)^2}{n} = 220 - \frac{900}{5} = 220 - 180 = 40$.
$S_{yy} = \sum y^2 - \frac{(\sum y)^2}{n} = 285 - \frac{1225}{5} = 285 - 245 = 40$.
$S_{xy} = \sum xy - \frac{(\sum x)(\sum y)}{n} = 250 - \frac{30\times35}{5} = 250 - 210 = 40$.

\item \textbf{Sample variances and covariance:}
$s_x^2 = \frac{S_{xx}}{n-1} = \frac{40}{4} = 10$.
$s_y^2 = \frac{S_{yy}}{n-1} = \frac{40}{4} = 10$.
$s_{xy} = \frac{S_{xy}}{n-1} = \frac{40}{4} = 10$.

\item \textbf{Correlation coefficient:}
$r = \frac{S_{xy}}{\sqrt{S_{xx}S_{yy}}} = \frac{40}{\sqrt{40\times40}} = \frac{40}{40} = 1$.

\item \textbf{Comment:} $r=1$ indicates a perfect positive linear correlation. Indeed, the data satisfies $y = x+1$ exactly; all points lie on a straight line with gradient 1.
\end{enumerate}
\end{exoresolu}

\begin{savaistu}[Correlation vs Causation]
A high correlation (close to $\pm 1$) does \emph{not} imply causation. For example, ice cream sales and drowning incidents are positively correlated, but eating ice cream does not cause drowning; both are influenced by a third variable: hot weather. In the GCE exam, always state "there is evidence of a linear association" rather than "$x$ causes $y$".
\end{savaistu}

\courssub{Least Squares Regression Line of $y$ on $x$}

\begin{definition}[Least Squares Regression Line of $y$ on $x$]
The \cle{least squares regression line of $y$ on $x$} is the line $y = a + bx$ that minimises the sum of squared vertical distances (residuals) between the observed $y_i$ and the predicted values $\hat{y}_i = a + bx_i$:
\[
\text{RSS} = \sum_{i=1}^n (y_i - \hat{y}_i)^2 = \sum_{i=1}^n (y_i - a - bx_i)^2.
\]
This line is also called the \cle{line of best fit} for predicting $y$ from $x$.
\end{definition}

\begin{propriete}[Equation of the least squares line of $y$ on $x$]
The coefficients are given by:
\[
b = \frac{S_{xy}}{S_{xx}}, \qquad a = \bar{y} - b\bar{x}.
\]
The equation can be written as $y = a + bx$ or equivalently $y - \bar{y} = b(x - \bar{x})$.
\textbf{Examinable proof:} Derivation via the normal equations:
\begin{align*}
\sum y &= na + b\sum x, \\
\sum xy &= a\sum x + b\sum x^2.
\end{align*}
Solving these simultaneously yields the formulas above. You must be able to reproduce this derivation.
\end{propriete}

\begin{methode}[Finding the least squares regression line of $y$ on $x$]
\begin{enumerate}
\item Calculate $\sum x, \sum y, \sum x^2, \sum xy$ (and $\sum y^2$ if you also need $r$ or $S_{yy}$).
\item Compute $S_{xx} = \sum x^2 - \frac{(\sum x)^2}{n}$ and $S_{xy} = \sum xy - \frac{(\sum x)(\sum y)}{n}$.
\item Compute the slope: $b = \dfrac{S_{xy}}{S_{xx}}$.
\item Compute the intercept: $a = \bar{y} - b\bar{x}$.
\item Write the equation: $y = a + bx$ (round $a$ and $b$ to 3 significant figures unless instructed otherwise).
\item \textbf{Check:} Verify that the line passes through $(\bar{x}, \bar{y})$.
\end{enumerate}
\textbf{Reflex:} Never rearrange the $y$ on $x$ line to get the $x$ on $y$ line; they are different (unless $r = \pm 1$).
\end{methode}

\begin{exoresolu}[Least squares line of $y$ on $x$ for agricultural data]
An agricultural researcher collects data on the amount of fertilizer applied (in kg/ha, $x$) and the resulting crop yield (in tonnes/ha, $y$) for 6 plots:

\begin{center}
\begin{tabular}{|c|c|c|c|c|c|c|}
\hline
$x$ & 10 & 20 & 30 & 40 & 50 & 60 \\
\hline
$y$ & 2.1 & 2.9 & 3.8 & 4.5 & 5.2 & 6.0 \\
\hline
\end{tabular}
\end{center}

Find the equation of the least squares regression line of $y$ on $x$. Use it to estimate the yield when 35 kg/ha of fertilizer is applied.
\tcblower
\textbf{Solution.}

$n=6$.

\begin{enumerate}
\item \textbf{Sums:}
$\sum x = 10+20+30+40+50+60 = 210$.
$\sum y = 2.1+2.9+3.8+4.5+5.2+6.0 = 24.5$.
$\sum x^2 = 100+400+900+1600+2500+3600 = 9100$.
$\sum y^2 = 4.41+8.41+14.44+20.25+27.04+36.00 = 110.55$.
$\sum xy = 10(2.1)+20(2.9)+30(3.8)+40(4.5)+50(5.2)+60(6.0) = 21+58+114+180+260+360 = 993$.

\item \textbf{Means:}
$\bar{x} = 210/6 = 35$.
$\bar{y} = 24.5/6 = 4.08333\ldots$

\item \textbf{Corrected sums:}
$S_{xx} = \sum x^2 - \frac{(\sum x)^2}{n} = 9100 - \frac{44100}{6} = 9100 - 7350 = 1750$.
$S_{xy} = \sum xy - \frac{(\sum x)(\sum y)}{n} = 993 - \frac{210\times24.5}{6} = 993 - \frac{5145}{6} = 993 - 857.5 = 135.5$.

\item \textbf{Slope and intercept:}
$b = \frac{S_{xy}}{S_{xx}} = \frac{135.5}{1750} = 0.07742857\ldots \approx 0.0774$ (4 s.f. or 3 s.f. as required).
$a = \bar{y} - b\bar{x} = 4.08333\ldots - 0.07742857\ldots \times 35 = 4.08333\ldots - 2.71 = 1.37333\ldots \approx 1.37$ (3 s.f.).

\item \textbf{Regression line:} $y = 1.37 + 0.0774x$ (or $y = 1.373 + 0.0774x$ keeping more precision).

\item \textbf{Estimate for $x=35$:} Since $\bar{x}=35$, the estimate is exactly $\bar{y} = 4.08$ tonnes/ha. Using the equation: $y = 1.37 + 0.0774(35) = 1.37 + 2.709 = 4.079 \approx 4.08$ tonnes/ha.
\end{enumerate}
\end{exoresolu}

\courssub{Least Squares Regression Line of $x$ on $y$}

\begin{definition}[Least Squares Regression Line of $x$ on $y$]
When we want to predict $x$ from $y$, we minimise the sum of squared \cle{horizontal} distances. The line is written as $x = c + dy$ (or $x = a' + b'y$). The formulas are symmetric to the $y$ on $x$ case.
\end{definition}

\begin{propriete}[Equation of the least squares line of $x$ on $y$]
The coefficients are given by:
\[
d = \frac{S_{xy}}{S_{yy}}, \qquad c = \bar{x} - d\bar{y}.
\]
The equation is $x = c + dy$ or $x - \bar{x} = d(y - \bar{y})$.
\textbf{Relationship between the two lines:}
\begin{itemize}
\item Both lines pass through the mean point $(\bar{x}, \bar{y})$.
\item The product of the slopes is $b \times d = \frac{S_{xy}}{S_{xx}} \cdot \frac{S_{xy}}{S_{yy}} = \frac{S_{xy}^2}{S_{xx}S_{yy}} = r^2$.
\item The slopes are reciprocals ($d = 1/b$) \textbf{only if} $r = \pm 1$ (perfect correlation).
\item The angle between the two lines decreases as $|r|$ increases; they coincide when $r = \pm 1$.
\end{itemize}
\end{propriete}

\begin{methode}[Finding the least squares regression line of $x$ on $y$]
\begin{enumerate}
\item Compute $S_{yy} = \sum y^2 - \frac{(\sum y)^2}{n}$ and $S_{xy}$ (same as before).
\item Compute the slope: $d = \dfrac{S_{xy}}{S_{yy}}$.
\item Compute the intercept: $c = \bar{x} - d\bar{y}$.
\item Write the equation: $x = c + dy$.
\item \textbf{Check:} Verify that the line passes through $(\bar{x}, \bar{y})$.
\end{enumerate}
\textbf{Reflex:} Do not simply rearrange the $y$ on $x$ line; that gives a different line (unless $r=\pm1$). Always compute $d$ using $S_{yy}$.
\end{methode}

\begin{exoresolu}[Least squares line of $x$ on $y$ for the same agricultural data]
Using the data from the previous example (fertilizer $x$, yield $y$), find the least squares regression line of $x$ on $y$. Hence estimate the amount of fertilizer needed to achieve a yield of 4.5 tonnes/ha.
\tcblower
\textbf{Solution.}

We already have: $n=6$, $\sum x=210$, $\sum y=24.5$, $\sum x^2=9100$, $\sum y^2=110.55$, $\sum xy=993$, $\bar{x}=35$, $\bar{y}=4.08333\ldots$, $S_{xy}=135.5$.

\begin{enumerate}
\item \textbf{Compute $S_{yy}$:}
$S_{yy} = \sum y^2 - \frac{(\sum y)^2}{n} = 110.55 - \frac{600.25}{6} = 110.55 - 100.04166\ldots = 10.50833\ldots$

\item \textbf{Slope $d$:}
$d = \frac{S_{xy}}{S_{yy}} = \frac{135.5}{10.50833\ldots} = 12.894\ldots \approx 12.9$ (3 s.f.).

\item \textbf{Intercept $c$:}
$c = \bar{x} - d\bar{y} = 35 - 12.894\ldots \times 4.08333\ldots = 35 - 52.65\ldots = -17.65\ldots \approx -17.7$ (3 s.f.).

\item \textbf{Regression line of $x$ on $y$:}
$x = -17.7 + 12.9 y$ (or $x = 12.9y - 17.7$).

\item \textbf{Estimate for $y=4.5$:}
$x = -17.65 + 12.894(4.5) = -17.65 + 58.023 = 40.373 \approx 40.4$ kg/ha.

\item \textbf{Check:} The two lines intersect at $(\bar{x}, \bar{y}) = (35, 4.0833)$. For $y=4.0833$, $x = -17.65 + 12.894(4.0833) \approx 35$. Correct.
\end{enumerate}
\end{exoresolu}

\courssub{Residual Sum of Squares (RSS) and the Line of Best Fit}

\begin{definition}[Residuals and Residual Sum of Squares]
For a fitted regression line:
\begin{itemize}
\item For the line $y = a + bx$ ($y$ on $x$), the \cle{residual} for point $i$ is the vertical distance: $e_i = y_i - \hat{y}_i = y_i - (a + bx_i)$.
\item The \cle{Residual Sum of Squares} for $y$ on $x$ is $\text{RSS}_{y|x} = \sum e_i^2 = \sum (y_i - a - bx_i)^2$.
\item For the line $x = c + dy$ ($x$ on $y$), the residual is the horizontal distance: $f_i = x_i - \hat{x}_i = x_i - (c + dy_i)$.
\item The RSS for $x$ on $y$ is $\text{RSS}_{x|y} = \sum f_i^2 = \sum (x_i - c - dy_i)^2$.
\end{itemize}
\end{definition}

\begin{propriete}[Formulas for RSS]
The RSS can be computed without calculating individual residuals:
\begin{align*}
\text{RSS}_{y|x} &= S_{yy} - b S_{xy} = S_{yy} - \frac{S_{xy}^2}{S_{xx}} = S_{yy}(1 - r^2), \\
\text{RSS}_{x|y} &= S_{xx} - d S_{xy} = S_{xx} - \frac{S_{xy}^2}{S_{yy}} = S_{xx}(1 - r^2).
\end{align*}
The \cle{coefficient of determination} is $R^2 = r^2 = 1 - \frac{\text{RSS}_{y|x}}{S_{yy}}$. It represents the proportion of the variation in $y$ explained by the linear regression on $x$.
\end{propriete}

\begin{methode}[Calculating RSS and $R^2$]
\begin{enumerate}
\item Compute $S_{xx}, S_{yy}, S_{xy}$ and $r$ (or $b$ and $d$).
\item For $y$ on $x$: $\text{RSS}_{y|x} = S_{yy} - \frac{S_{xy}^2}{S_{xx}}$ or $S_{yy}(1-r^2)$.
\item For $x$ on $y$: $\text{RSS}_{x|y} = S_{xx} - \frac{S_{xy}^2}{S_{yy}}$ or $S_{xx}(1-r^2)$.
\item Compute $R^2 = r^2 = 1 - \frac{\text{RSS}_{y|x}}{S_{yy}}$.
\item Interpret $R^2$ as a percentage: "$R^2 \times 100\%$ of the variation in $y$ is explained by the linear relationship with $x$."
\end{enumerate}
\textbf{Reflex:} RSS is always non-negative. A smaller RSS indicates a better fit. $R^2$ close to 1 means a good linear fit.
\end{methode}

\begin{exoresolu}[RSS for both lines and interpretation]
For the agricultural data (fertilizer $x$, yield $y$) used in the previous two examples, calculate:
\begin{enumerate}
\item The Residual Sum of Squares for the regression line of $y$ on $x$.
\item The Residual Sum of Squares for the regression line of $x$ on $y$.
\item The coefficient of determination $R^2$ and interpret it.
\end{enumerate}
\tcblower
\textbf{Solution.}

We have: $S_{xx}=1750$, $S_{yy}=10.50833\ldots$, $S_{xy}=135.5$.
$r = \frac{S_{xy}}{\sqrt{S_{xx}S_{yy}}} = \frac{135.5}{\sqrt{1750 \times 10.50833\ldots}} = \frac{135.5}{\sqrt{18389.58\ldots}} = \frac{135.5}{135.608\ldots} \approx 0.9992$.
$r^2 \approx 0.9984$.

\begin{enumerate}
\item \textbf{RSS for $y$ on $x$:}
$\text{RSS}_{y|x} = S_{yy} - \frac{S_{xy}^2}{S_{xx}} = 10.50833\ldots - \frac{135.5^2}{1750} = 10.50833\ldots - \frac{18360.25}{1750}$.
$\frac{18360.25}{1750} = 10.4915714\ldots$
$\text{RSS}_{y|x} = 10.50833\ldots - 10.4915714\ldots = 0.01676\ldots \approx 0.0168$.

\item \textbf{RSS for $x$ on $y$:}
$\text{RSS}_{x|y} = S_{xx} - \frac{S_{xy}^2}{S_{yy}} = 1750 - \frac{18360.25}{10.50833\ldots} = 1750 - 1747.18\ldots = 2.82\ldots \approx 2.82$.

\item \textbf{Coefficient of determination:}
$R^2 = r^2 = 0.9984$ (or $1 - \frac{\text{RSS}_{y|x}}{S_{yy}} = 1 - \frac{0.01676}{10.50833} = 0.9984$).
\textbf{Interpretation:} 99.84\% of the variation in crop yield ($y$) is explained by the linear relationship with fertilizer amount ($x$). The fit is extremely good; the linear model is highly appropriate for this data.
\end{enumerate}
\end{exoresolu}

\courssub{Rank Correlation: Spearman and Kendall}

\begin{definition}[Rank Correlation]
When data is ordinal (ranks) or when the relationship is monotonic but not necessarily linear, or when outliers make Pearson's $r$ unreliable, we use \cle{rank correlation coefficients}. The two main ones are \cle{Spearman's rank correlation coefficient} ($r_s$ or $\rho$) and \cle{Kendall's rank correlation coefficient} ($\tau$).
\end{definition}

\begin{methode}[Spearman's rank correlation coefficient ($r_s$)]
\begin{enumerate}
\item \textbf{Rank the data:} Assign ranks to $x$ values (1 for smallest, 2 for next, etc.) and to $y$ values. If there are ties (equal values), assign the \cle{average rank} to each tied value. Example: values 10, 12, 12, 15 get ranks 1, 2.5, 2.5, 4.
\item \textbf{Compute the difference in ranks:} $d_i = \text{rank}(x_i) - \text{rank}(y_i)$.
\item \textbf{Calculate $r_s$:}
\begin{itemize}
\item If there are \textbf{no ties} (or very few ties), use the shortcut formula:
$r_s = 1 - \frac{6\sum d_i^2}{n(n^2-1)}$.
\item If there are \textbf{many ties}, the shortcut formula is biased. Use the Pearson formula on the ranks:
$r_s = \frac{\sum (R_{x_i} - \bar{R}_x)(R_{y_i} - \bar{R}_y)}{\sqrt{\sum (R_{x_i} - \bar{R}_x)^2 \sum (R_{y_i} - \bar{R}_y)^2}}$,
where $\bar{R}_x = \bar{R}_y = \frac{n+1}{2}$.
\end{itemize}
\item \textbf{Interpret:} Same as Pearson's $r$: $+1$ perfect positive monotonic association, $-1$ perfect negative monotonic association, $0$ no monotonic association.
\item \textbf{Hypothesis test:} For $n \ge 10$, $r_s$ is approximately normally distributed; for small $n$, use critical value tables provided in the exam.
\end{enumerate}
\textbf{Key formula (no ties):} $r_s = 1 - \frac{6\sum d^2}{n(n^2-1)}$.
\textbf{Reflex:} Always check for ties. If ties exist, mention them and decide which formula to use. In GCE exams, if ties are present, the Pearson-on-ranks formula is safer.
\end{methode}

\begin{methode}[Kendall's rank correlation coefficient ($\tau$)]
\begin{enumerate}
\item \textbf{Rank both variables} (same as Spearman, handling ties by average ranks).
\item \textbf{Arrange the pairs in order of one variable} (say $x$ ranks ascending). Then look at the sequence of $y$ ranks.
\item \textbf{Count concordant and discordant pairs:} For each pair of observations $(i,j)$ with $i < j$ (in the ordered $x$ list):
\begin{itemize}
\item The pair is \cle{concordant} if $(R_{y_i} - R_{y_j})$ has the same sign as $(R_{x_i} - R_{x_j})$ (which is positive since $x$ ranks are in ascending order). So concordant means $R_{y_i} < R_{y_j}$.
\item The pair is \cle{discordant} if $R_{y_i} > R_{y_j}$ (an inversion).
\item If $R_{y_i} = R_{y_j}$ (tie in $y$), the pair is neither concordant nor discordant (handled in $\tau_b$ or $\tau_c$).
\end{itemize}
\item \textbf{Compute $\tau$ (simple version, no ties or few ties):}
$\tau = \frac{N_c - N_d}{N_0}$, where $N_c$ = number of concordant pairs, $N_d$ = number of discordant pairs, $N_0 = \frac{1}{2}n(n-1)$ = total pairs.
\item \textbf{With ties (Kendall's $\tau_b$):}
$\tau_b = \frac{N_c - N_d}{\sqrt{(N_0 - N_1)(N_0 - N_2)}}$,
where $N_1 = \sum \frac{t_k(t_k-1)}{2}$ (ties in $x$), $N_2 = \sum \frac{u_k(u_k-1)}{2}$ (ties in $y$), $t_k, u_k$ are sizes of tie groups.
For GCE, the simple version $\tau = \frac{N_c - N_d}{N_0}$ is often used when ties are few or absent.
\end{enumerate}
\end{methode}

\begin{aretenir}[Uses and Limitations of Kendall's $\tau$]
\begin{itemize}
\item \textbf{Uses:} More robust to outliers than Spearman; better for small samples; directly interpretable as the difference between the probability of concordance and discordance; $\tau$ estimates a population parameter (Kendall's $\tau$) naturally.
\item \textbf{Limitations:} Computationally heavier for large $n$ (though $O(n \log n)$ algorithms exist); less intuitive than Spearman's $r_s$; not as widely tabulated for hypothesis testing in exam tables; the simple formula ignores ties, so $\tau_b$ is needed for tied data.
\end{itemize}
\end{aretenir}

\begin{exoresolu}[Spearman and Kendall for examination ranks]
Two examiners rank 8 candidates in an oral examination:

\begin{center}
\begin{tabular}{|c|c|c|c|c|c|c|c|c|}
\hline
Candidate & A & B & C & D & E & F & G & H \\
\hline
Examiner 1 ($x$) & 1 & 2 & 3 & 4 & 5 & 6 & 7 & 8 \\
\hline
Examiner 2 ($y$) & 2 & 1 & 4 & 3 & 6 & 5 & 8 & 7 \\
\hline
\end{tabular}
\end{center}

\begin{enumerate}
\item Calculate Spearman's rank correlation coefficient $r_s$.
\item Calculate Kendall's $\tau$.
\item Comment on the agreement between the examiners.
\end{enumerate}
\tcblower
\textbf{Solution.}

\begin{enumerate}
\item \textbf{Spearman's $r_s$:} The data are already ranks (1 to 8). No ties.
Compute $d_i = R_{x_i} - R_{y_i}$:
\begin{center}
\begin{tabular}{|c|c|c|c|c|}
\hline
Candidate & $R_x$ & $R_y$ & $d$ & $d^2$ \\
\hline
A & 1 & 2 & -1 & 1 \\
B & 2 & 1 & 1 & 1 \\
C & 3 & 4 & -1 & 1 \\
D & 4 & 3 & 1 & 1 \\
E & 5 & 6 & -1 & 1 \\
F & 6 & 5 & 1 & 1 \\
G & 7 & 8 & -1 & 1 \\
H & 8 & 7 & 1 & 1 \\
\hline
Total & & & & 8 \\
\hline
\end{tabular}
\end{center}
$n=8$, $\sum d^2 = 8$.
$r_s = 1 - \frac{6 \times 8}{8(64-1)} = 1 - \frac{48}{8 \times 63} = 1 - \frac{48}{504} = 1 - \frac{2}{21} = \frac{19}{21} \approx 0.9048$.

\item \textbf{Kendall's $\tau$:} Arrange in order of Examiner 1 (already ordered). Look at Examiner 2's ranks: $2, 1, 4, 3, 6, 5, 8, 7$.
Count concordant and discordant pairs among the $y$ ranks. Total pairs $N_0 = \frac{1}{2} \times 8 \times 7 = 28$.
We count inversions (discordant pairs) in the sequence $2, 1, 4, 3, 6, 5, 8, 7$:
- (2,1): inversion $\to$ discordant
- (2,4): concordant
- (2,3): concordant
- (2,6): concordant
- (2,5): concordant
- (2,8): concordant
- (2,7): concordant
- (1,4): concordant
- (1,3): concordant
- (1,6): concordant
- (1,5): concordant
- (1,8): concordant
- (1,7): concordant
- (4,3): inversion $\to$ discordant
- (4,6): concordant
- (4,5): concordant
- (4,8): concordant
- (4,7): concordant
- (3,6): concordant
- (3,5): concordant
- (3,8): concordant
- (3,7): concordant
- (6,5): inversion $\to$ discordant
- (6,8): concordant
- (6,7): concordant
- (5,8): concordant
- (5,7): concordant
- (8,7): inversion $\to$ discordant
Total discordant $N_d = 4$ (the inversions: (2,1), (4,3), (6,5), (8,7)).
Concordant $N_c = 28 - 4 = 24$.
$\tau = \frac{N_c - N_d}{N_0} = \frac{24 - 4}{28} = \frac{20}{28} = \frac{5}{7} \approx 0.7143$.

\item \textbf{Comment:} Both coefficients indicate strong positive agreement between the examiners. Spearman's $r_s = 0.905$ is higher than Kendall's $\tau = 0.714$; this is typical because $\tau$ is more conservative (it counts pairwise agreements/disagreements directly). The examiners consistently rank candidates in a very similar order, with only minor swaps of adjacent ranks (each adjacent pair is swapped). At the 5\% significance level, both coefficients would indicate statistically significant agreement (critical values for $n=8$: $r_s \approx 0.738$, $\tau \approx 0.619$).
\end{enumerate}
\end{exoresolu}

\begin{exercice}[Practice: Complete Regression and Correlation Analysis]
The table below shows the number of hours of revision ($x$) and the examination mark ($y$) for 8 students.

\begin{center}
\begin{tabular}{|c|c|c|c|c|c|c|c|c|}
\hline
Hours $x$ & 10 & 12 & 14 & 16 & 18 & 20 & 22 & 24 \\
\hline
Mark $y$  & 45 & 50 & 55 & 60 & 65 & 70 & 75 & 80 \\
\hline
\end{tabular}
\end{center}

\begin{enumerate}
\item Draw a scatter diagram and comment.
\item Calculate $S_{xx}, S_{yy}, S_{xy}$ and the product moment correlation coefficient $r$.
\item Find the least squares regression line of $y$ on $x$.
\item Find the least squares regression line of $x$ on $y$.
\item Calculate $\text{RSS}_{y|x}$ and $\text{RSS}_{x|y}$ and the coefficient of determination $R^2$.
\item Use the appropriate line to estimate the mark for a student who revises for 15 hours, and the hours needed to achieve a mark of 72.
\item Calculate Spearman's $r_s$ and Kendall's $\tau$ for this data (treating the raw data as ranks after verifying they are already in rank order). Compare with $r$.
\end{enumerate}
\end{exercice}

\begin{corrige}[Exercise: Complete Regression and Correlation Analysis]
\begin{enumerate}
\item \textbf{Scatter diagram:} Points lie exactly on a straight line. Perfect positive linear correlation.
\item \textbf{Sums:} $n=8$, $\sum x = 136$, $\sum y = 500$, $\sum x^2 = 2520$, $\sum y^2 = 32400$, $\sum xy = 9160$.
$\bar{x} = 17$, $\bar{y} = 62.5$.
$S_{xx} = 2520 - 136^2/8 = 2520 - 2312 = 208$.
$S_{yy} = 32400 - 500^2/8 = 32400 - 31250 = 1150$.
$S_{xy} = 9160 - (136\times500)/8 = 9160 - 8500 = 660$.
$r = 660/\sqrt{208\times1150} = 660/\sqrt{239200} = 660/489.08 = 1$ (exactly, since $y = 2.5x + 20$).
\item \textbf{$y$ on $x$:} $b = 660/208 = 3.17308\ldots$? Wait: $y = 2.5x + 20$ exactly. Let's recalc: $S_{xy}=660$, $S_{xx}=208$, $b=660/208 = 3.173$? But $y = 2.5x + 20$ gives $b=2.5$. Check sums: $x$: 10,12,14,16,18,20,22,24. $y$: 45,50,55,60,65,70,75,80. $\sum x = 136$, $\sum y = 500$. $\sum xy = 10*45+12*50+14*55+16*60+18*65+20*70+22*75+24*80 = 450+600+770+960+1170+1400+1650+1920 = 8920$. Not 9160. My previous sum was wrong. Let's recompute properly.
$\sum xy = 450+600+770+960+1170+1400+1650+1920 = 8920$.
$S_{xy} = 8920 - (136*500)/8 = 8920 - 8500 = 420$.
$S_{xx} = 2520 - 2312 = 208$.
$b = 420/208 = 2.019$? No, $y = 2.5x + 20$ exactly. $\sum x^2 = 100+144+196+256+324+400+484+576 = 2480$. Not 2520.
Let's do it cleanly:
$x$: 10,12,14,16,18,20,22,24. $\sum x = 136$. $\sum x^2 = 100+144+196+256+324+400+484+576 = 2480$.
$y$: 45,50,55,60,65,70,75,80. $\sum y = 500$. $\sum y^2 = 2025+2500+3025+3600+4225+4900+5625+6400 = 32300$.
$\sum xy = 450+600+770+960+1170+1400+1650+1920 = 8920$.
$S_{xx} = 2480 - 136^2/8 = 2480 - 2312 = 168$.
$S_{yy} = 32300 - 500^2/8 = 32300 - 31250 = 1050$.
$S_{xy} = 8920 - 136*500/8 = 8920 - 8500 = 420$.
$b = 420/168 = 2.5$. $a = 62.5 - 2.5*17 = 62.5 - 42.5 = 20$. $y = 20 + 2.5x$. Correct.
$r = 420/\sqrt{168*1050} = 420/\sqrt{176400} = 420/420 = 1$.
\item \textbf{$x$ on $y$:} $d = 420/1050 = 0.4$. $c = 17 - 0.4*62.5 = 17 - 25 = -8$. $x = -8 + 0.4y$. (Equivalent to $y = 20 + 2.5x$).
\item \textbf{RSS:} $r=1 \Rightarrow \text{RSS}_{y|x}=0$, $\text{RSS}_{x|y}=0$, $R^2=1$.
\item \textbf{Estimates:} For $x=15$, $y = 20 + 2.5(15) = 57.5$. For $y=72$, $x = -8 + 0.4(72) = 20.8$ hours.
\item \textbf{Rank correlation:} Data is perfectly linear, so ranks are perfectly correlated. $r_s = 1$, $\tau = 1$. Same as $r$.
\end{enumerate}
\end{corrige}

\newpage

\coursec{Exercises}

\courssub{I check my knowledge}

\begin{exercice}[MCQ: Regression Line Properties]
Which of the following statements is \textbf{always true} for the least squares regression line of $y$ on $x$?
\begin{enumerate}
    \item It passes through the origin $(0,0)$.
    \item It passes through the mean point $(\bar{x}, \bar{y})$.
    \item Its slope is equal to the product moment correlation coefficient $r$.
    \item The sum of the residuals is equal to the sum of the $y$-values.
\end{enumerate}
\end{exercice}

\begin{exercice}[True/False: Correlation vs Causation]
State whether the following statements are True or False. Justify your answer briefly.
\begin{enumerate}
    \item If the product moment correlation coefficient $r = 0$, then there is no relationship whatsoever between the variables $x$ and $y$.
    \item The regression line of $x$ on $y$ can be obtained by simply rearranging the equation of the regression line of $y$ on $x$ to make $x$ the subject.
    \item Spearman's rank correlation coefficient $R_s$ can only be calculated if the data are quantitative (numerical).
    \item For a given data set, the residual sum of squares (RSS) for the least squares line of $y$ on $x$ is always less than or equal to the RSS for any other line.
\end{enumerate}
\end{exercice}

\begin{exercice}[Definitions and Direct Formulas]
\begin{enumerate}
    \item Define the \emph{residual} for a data point $(x_i, y_i)$ with respect to the regression line $y = a + bx$.
    \item Write down the formula for the slope $b$ of the least squares regression line of $y$ on $x$ in terms of $S_{xy}$ and $S_{xx}$.
    \item State the relationship between the product moment correlation coefficient $r$, the slope $b$ of the regression line of $y$ on $x$, and the slope $b'$ of the regression line of $x$ on $y$.
    \item Given the ranks of 5 items by two judges, explain how to handle \emph{tied ranks} when calculating Spearman's rank correlation coefficient.
\end{enumerate}
\end{exercice}

\begin{exercice}[Direct Calculation: Covariance and Correlation]
For a bivariate data set of $n=6$ pairs, the following summary statistics are given:
\[
\sum x = 42,\quad \sum y = 60,\quad \sum x^2 = 320,\quad \sum y^2 = 650,\quad \sum xy = 440.
\]
Calculate:
\begin{enumerate}
    \item The sample covariance $S_{xy}$.
    \item The sample variances $S_{xx}$ and $S_{yy}$.
    \item The product moment correlation coefficient $r$.
    \item Comment on the strength and direction of the linear relationship.
\end{enumerate}
\end{exercice}

\courssub{I practise}

\begin{exercice}[Least Squares Regression Line: Rainfall and Crop Yield]
The table below shows the annual rainfall (in cm) and the corresponding maize yield (in tonnes per hectare) for a farm in Buea over 8 years.

\begin{center}
\begin{tabular}{|c|c|c|c|c|c|c|c|c|}
\hline
Rainfall $x$ (cm) & 120 & 145 & 130 & 160 & 155 & 140 & 170 & 150 \\
\hline
Yield $y$ (t/ha) & 2.1 & 3.0 & 2.5 & 3.8 & 3.5 & 2.8 & 4.0 & 3.2 \\
\hline
\end{tabular}
\end{center}

\begin{enumerate}
    \item Plot a scatter diagram of yield against rainfall. \emph{(Graph paper not required; describe the trend)}.
    \item Calculate the equation of the least squares regression line of $y$ on $x$ in the form $y = a + bx$.
    \item Use your equation to estimate the maize yield when the rainfall is $158$ cm.
    \item Calculate the residual for the year when rainfall was $145$ cm and yield was $3.0$ t/ha. Interpret its sign.
\end{enumerate}
\end{exercice}

\begin{exercice}[Regression Line of $x$ on $y$ and Comparison]
Using the same data from the previous exercise (Rainfall $x$, Yield $y$):
\begin{enumerate}
    \item Calculate the equation of the least squares regression line of $x$ on $y$ in the form $x = c + dy$.
    \item Using the line of $x$ on $y$, estimate the rainfall required to achieve a yield of $3.6$ t/ha.
    \item Rearrange the regression line of $y$ on $x$ (from the previous exercise) to make $x$ the subject. Compare this equation with the line of $x$ on $y$ obtained in (i). Explain why they are different.
\end{enumerate}
\end{exercice}

\begin{exercice}[Mayer's Method vs Least Squares]
A small business in Douala records the advertising expenditure (in 1000 FCFA) and the resulting sales (in 1000 FCFA) for 9 months.

\begin{center}
\begin{tabular}{|c|c|c|c|c|c|c|c|c|c|}
\hline
Advertising $x$ & 2 & 4 & 5 & 7 & 8 & 10 & 11 & 12 & 14 \\
\hline
Sales $y$ & 15 & 22 & 25 & 30 & 33 & 38 & 40 & 44 & 48 \\
\hline
\end{tabular}
\end{center}

\begin{enumerate}
    \item Divide the data into two groups of equal size (or as equal as possible) to apply Mayer's method. Find the mean points $(\bar{x}_1, \bar{y}_1)$ and $(\bar{x}_2, \bar{y}_2)$.
    \item Determine the equation of the regression line by Mayer's method (line joining the two mean points).
    \item Calculate the equation of the least squares regression line of $y$ on $x$.
    \item Compute the Residual Sum of Squares (RSS) for both lines. Which line provides a better fit? Justify.
\end{enumerate}
\end{exercice}

\begin{exercice}[Rank Correlation: Singing Competition]
Two judges rank 7 singers in a "Yaoundé Talent Show" as follows:

\begin{center}
\begin{tabular}{|c|c|c|c|c|c|c|c|}
\hline
Singer & A & B & C & D & E & F & G \\
\hline
Judge 1 Rank & 1 & 2 & 3 & 4 & 5 & 6 & 7 \\
\hline
Judge 2 Rank & 2 & 1 & 4 & 3 & 6 & 7 & 5 \\
\hline
\end{tabular}
\end{center}

\begin{enumerate}
    \item Calculate Spearman's rank correlation coefficient $R_s$.
    \item Interpret the value of $R_s$ in the context of the competition.
    \item If Singer H is added and both judges rank H as 8th, explain how $R_s$ would change without recalculating fully.
\end{enumerate}
\end{exercice}

\courssub{I prepare for the GCE}

\begin{exercice}[GCE Style: Mixed Regression and Correlation Problem -- 18 marks]
The ages ($x$ years) and prices ($y$ thousand FCFA) of 10 second-hand taxis in Douala are summarised below:
\[
n=10,\quad \sum x = 65,\quad \sum y = 820,\quad \sum x^2 = 485,\quad \sum y^2 = 72500,\quad \sum xy = 5020.
\]

\begin{enumerate}
    \item Calculate the product moment correlation coefficient $r$ between age and price. \hfill [3 marks]
    \item Interpret your value of $r$ in the context of the question. \hfill [2 marks]
    \item Find the equation of the least squares regression line of $y$ on $x$ in the form $y = a + bx$. \hfill [4 marks]
    \item A buyer wants to purchase a 6-year-old taxi. Use your regression line to estimate a fair price. \hfill [2 marks]
    \item The oldest taxi in the sample is 9 years old. Explain why it would be unreliable to use your regression line to estimate the price of a 12-year-old taxi. \hfill [2 marks]
    \item Calculate the Residual Sum of Squares (RSS) for the regression line of $y$ on $x$. \hfill [3 marks]
    \item Given that the regression line of $x$ on $y$ is $x = 7.2 - 0.0045y$, verify that $r^2 = b \times b'$, where $b$ and $b'$ are the slopes of the two regression lines. \hfill [2 marks]
\end{enumerate}
\end{exercice}

\begin{exercice}[GCE Style: Kendall's Rank Correlation and Interpretation -- 14 marks]
Five students are ranked by their Mathematics teacher and their Physics teacher according to their performance in the mock exams.

\begin{center}
\begin{tabular}{|c|c|c|c|c|c|}
\hline
Student & P & Q & R & S & T \\
\hline
Maths Rank & 1 & 2 & 3 & 4 & 5 \\
\hline
Physics Rank & 2 & 1 & 5 & 3 & 4 \\
\hline
\end{tabular}
\end{center}

\begin{enumerate}
    \item List all $\binom{5}{2} = 10$ pairs of students. For each pair, determine whether the rankings are \emph{concordant} or \emph{discordant}. \hfill [5 marks]
    \item Calculate Kendall's rank correlation coefficient $\tau$. \hfill [3 marks]
    \item Calculate Spearman's rank correlation coefficient $R_s$ for the same data. \hfill [3 marks]
    \item Compare the values of $\tau$ and $R_s$. State one \emph{use} and one \emph{limitation} of Kendall's $\tau$ compared to Spearman's $R_s$. \hfill [3 marks]
\end{enumerate}
\end{exercice}

\begin{exercice}[GCE Style: Regression Lines and Mean Point -- 12 marks]
For a certain bivariate data set, the least squares regression line of $y$ on $x$ is $y = 0.8x + 3$ and the regression line of $x$ on $y$ is $x = 0.45y + 1$.

\begin{enumerate}
    \item Find the coordinates of the mean point $(\bar{x}, \bar{y})$. \hfill [3 marks]
    \item Calculate the product moment correlation coefficient $r$. \hfill [3 marks]
    \item Given that $S_{xx} = 50$, find the value of $S_{yy}$. \hfill [3 marks]
    \item If a new data point $(\bar{x}, \bar{y} + 2)$ is added to the data set, describe qualitatively what would happen to the slope of the regression line of $y$ on $x$ and to the value of $r$. \hfill [3 marks]
\end{enumerate}
\end{exercice}

\newpage

\coursec{Solutions to the exercises}

\begin{corrige}[Exercise 1 — MCQ: Regression Line Properties]
The correct answer is \textbf{(b)}: the least squares regression line of $y$ on $x$ always passes through the mean point $(\bar{x}, \bar{y})$.

\textbf{Justification.} The least squares line is $y = a + bx$ with $b = \dfrac{S_{xy}}{S_{xx}}$ and $a = \bar{y} - b\bar{x}$. Substituting $x = \bar{x}$:
\[
a + b\bar{x} = (\bar{y} - b\bar{x}) + b\bar{x} = \bar{y},
\]
so the point $(\bar{x}, \bar{y})$ lies on the line, for \emph{any} data set.

\textbf{Why the others are false:}
\begin{itemize}
\item (a) The line passes through $(0,0)$ only if $a = \bar{y} - b\bar{x} = 0$, which is not true in general.
\item (c) The slope is $b = \dfrac{S_{xy}}{S_{xx}} = r\dfrac{s_y}{s_x}$, which equals $r$ only when $s_x = s_y$.
\item (d) The sum of the residuals is always $0$ (a property of least squares), not the sum of the $y$-values.
\end{itemize}
\end{corrige}

\begin{corrige}[Exercise 2 — True/False: Correlation vs Causation]
\begin{enumerate}
\item \textbf{False.} $r = 0$ only means there is no \emph{linear} relationship. The variables may be strongly related in a non-linear way (e.g. $y = x^2$ on symmetric data gives $r = 0$).
\item \textbf{False.} The line of $x$ on $y$ minimises the sum of squares of \emph{horizontal} residuals, whereas the line of $y$ on $x$ minimises \emph{vertical} residuals. They are two different lines (they coincide only when $r = \pm 1$). Rearranging one equation does not give the other.
\item \textbf{False.} Spearman's coefficient only requires \emph{ranks}. Qualitative data that can be ordered (e.g. beauty, taste, quality of a performance) can be ranked and $R_s$ computed, even with no numerical measurements at all.
\item \textbf{True.} By definition, the least squares line is the line that \emph{minimises} the residual sum of squares. Hence its RSS is less than or equal to the RSS of any other line fitted to the same data.
\end{enumerate}
\end{corrige}

\begin{corrige}[Exercise 3 — Definitions and Direct Formulas]
\begin{enumerate}
\item The \emph{residual} of the point $(x_i, y_i)$ is the vertical signed distance between the observed value and the value predicted by the line:
\[
e_i = y_i - \hat{y}_i = y_i - (a + bx_i).
\]
\item $\displaystyle b = \frac{S_{xy}}{S_{xx}}$, where $S_{xy} = \sum x_iy_i - \frac{\sum x_i \sum y_i}{n}$ and $S_{xx} = \sum x_i^2 - \frac{(\sum x_i)^2}{n}$.
\item The slopes satisfy $\;b \times b' = r^2$, i.e. $r = \pm\sqrt{b\,b'}$, the sign of $r$ being the common sign of $b$ and $b'$. (Indeed $bb' = \dfrac{S_{xy}}{S_{xx}}\cdot\dfrac{S_{xy}}{S_{yy}} = \dfrac{S_{xy}^2}{S_{xx}S_{yy}} = r^2$.)
\item When two or more items are tied, assign to each of them the \emph{mean of the ranks they occupy}. For example, if two items tie for positions 3 and 4, each receives rank $\frac{3+4}{2} = 3.5$, and the next item receives rank 5. One then computes $d$ and $\sum d^2$ as usual in $R_s = 1 - \dfrac{6\sum d^2}{n(n^2-1)}$ (a correction factor may be applied if ties are numerous).
\end{enumerate}
\end{corrige}

\begin{corrige}[Exercise 4 — Direct Calculation: Covariance and Correlation]
Given $n = 6$, $\sum x = 42$, $\sum y = 60$, $\sum x^2 = 320$, $\sum y^2 = 650$, $\sum xy = 440$.
\begin{enumerate}
\item \textbf{Covariance:}
\[
S_{xy} = \sum xy - \frac{\sum x \sum y}{n} = 440 - \frac{42 \times 60}{6} = 440 - 420 = \boxed{20}.
\]
\item \textbf{Variances:}
\[
S_{xx} = 320 - \frac{42^2}{6} = 320 - 294 = \boxed{26}, \qquad
S_{yy} = 650 - \frac{60^2}{6} = 650 - 600 = \boxed{50}.
\]
\item \textbf{Product moment correlation coefficient:}
\[
r = \frac{S_{xy}}{\sqrt{S_{xx}S_{yy}}} = \frac{20}{\sqrt{26 \times 50}} = \frac{20}{\sqrt{1300}} = \frac{20}{36.0555\ldots} \approx \boxed{0.555}.
\]
\item \textbf{Comment:} $r \approx 0.56 > 0$, so there is a \emph{moderate positive linear relationship}: as $x$ increases, $y$ tends to increase, but the points are fairly scattered around any straight line.
\end{enumerate}
\end{corrige}

\begin{corrige}[Exercise 5 — Least Squares Regression Line: Rainfall and Crop Yield]
\textbf{Preliminary sums} ($n = 8$):
\[
\sum x = 1170,\quad \sum y = 24.9,\quad \sum x^2 = 172950,\quad \sum xy = 3714.5,\quad \sum y^2 = 80.43.
\]
\[
\bar{x} = \frac{1170}{8} = 146.25, \qquad \bar{y} = \frac{24.9}{8} = 3.1125.
\]
\begin{enumerate}
\item \textbf{Scatter diagram:} plotting yield $y$ against rainfall $x$, the points rise from the lower left to the upper right in a roughly linear band: the trend is \emph{positive and approximately linear} (more rainfall $\Rightarrow$ higher yield).
\item \textbf{Regression line of $y$ on $x$:}
\[
S_{xx} = 172950 - \frac{1170^2}{8} = 172950 - 171112.5 = 1837.5,
\]
\[
S_{xy} = 3714.5 - \frac{1170 \times 24.9}{8} = 3714.5 - 3641.625 = 72.875.
\]
\[
b = \frac{72.875}{1837.5} = 0.03966\ldots \approx 0.0397, \qquad a = 3.1125 - 0.03966\ldots \times 146.25 \approx -2.6875.
\]
\[
\boxed{y = -2.688 + 0.0397\,x}
\]
\item \textbf{Estimate for $x = 158$:}
\[
y = -2.6875 + 0.03966\ldots \times 158 \approx 3.58 \text{ t/ha}.
\]
\item \textbf{Residual at $(145, 3.0)$:} predicted value $\hat{y} = -2.6875 + 0.03966\ldots \times 145 \approx 3.063$.
\[
e = y - \hat{y} = 3.0 - 3.063 = -0.063 \text{ t/ha}.
\]
The residual is \emph{negative}: the actual yield that year was about $0.06$ t/ha \emph{below} what the model predicts — the line slightly overestimates the yield for that year.
\end{enumerate}
\end{corrige}

\begin{corrige}[Exercise 6 — Regression Line of $x$ on $y$ and Comparison]
Using the sums of Exercise 5: $S_{xy} = 72.875$, and
\[
S_{yy} = 80.43 - \frac{24.9^2}{8} = 80.43 - 77.50125 = 2.92875.
\]
\begin{enumerate}
\item \textbf{Line of $x$ on $y$:} $x = c + dy$ with
\[
d = \frac{S_{xy}}{S_{yy}} = \frac{72.875}{2.92875} \approx 24.88, \qquad c = \bar{x} - d\bar{y} = 146.25 - 24.88 \times 3.1125 \approx 68.81.
\]
\[
\boxed{x = 68.81 + 24.88\,y}
\]
\item \textbf{Rainfall for a yield of $3.6$ t/ha:}
\[
x = 68.81 + 24.88 \times 3.6 \approx 158.4 \text{ cm}.
\]
\item \textbf{Rearranging the line of $y$ on $x$:} from $y = -2.688 + 0.0397x$,
\[
x = \frac{y + 2.688}{0.0397} \approx 67.7 + 25.2\,y.
\]
This is \emph{close to but different from} $x = 68.81 + 24.88y$. The two lines are different because they answer different minimisation problems: the line of $y$ on $x$ minimises the sum of squared \emph{vertical} deviations, while the line of $x$ on $y$ minimises the sum of squared \emph{horizontal} deviations. Both pass through $(\bar{x}, \bar{y}) = (146.25, 3.1125)$ but have different slopes; they would coincide only if $r = \pm 1$ (here $r = \sqrt{0.0397 \times 24.88} \approx 0.994$, which is why they are nearly identical).
\end{enumerate}
\end{corrige}

\begin{corrige}[Exercise 7 — Mayer's Method vs Least Squares]
\textbf{Preliminary sums} ($n = 9$): $\sum x = 73$, $\sum y = 295$, $\sum x^2 = 719$, $\sum xy = 2737$, $\sum y^2 = 10607$.
\begin{enumerate}
\item \textbf{Mayer's groups.} Order the data by $x$ and split into a group of 4 and a group of 5:
\[
G_1 = \{(2,15),(4,22),(5,25),(7,30)\}, \quad G_2 = \{(8,33),(10,38),(11,40),(12,44),(14,48)\}.
\]
\[
\bar{x}_1 = \frac{18}{4} = 4.5,\quad \bar{y}_1 = \frac{92}{4} = 23; \qquad
\bar{x}_2 = \frac{55}{5} = 11,\quad \bar{y}_2 = \frac{203}{5} = 40.6.
\]
\item \textbf{Mayer's line} through $(4.5, 23)$ and $(11, 40.6)$:
\[
b_M = \frac{40.6 - 23}{11 - 4.5} = \frac{17.6}{6.5} \approx 2.708, \qquad a_M = 23 - 2.708 \times 4.5 \approx 10.82.
\]
\[
\boxed{y = 10.82 + 2.708\,x}
\]
\item \textbf{Least squares line:}
\[
S_{xx} = 719 - \frac{73^2}{9} = 126.89, \qquad S_{xy} = 2737 - \frac{73 \times 295}{9} = 344.22.
\]
\[
b = \frac{344.22}{126.89} \approx 2.713, \qquad a = \frac{295}{9} - 2.713 \times \frac{73}{9} \approx 32.78 - 22.00 = 10.77.
\]
\[
\boxed{y = 10.77 + 2.713\,x}
\]
\item \textbf{Residual sums of squares.} For the least squares line, use $\mathrm{RSS} = S_{yy} - \dfrac{S_{xy}^2}{S_{xx}}$:
\[
S_{yy} = 10607 - \frac{295^2}{9} = 937.56, \qquad \mathrm{RSS}_{LS} = 937.56 - \frac{344.22^2}{126.89} \approx 937.56 - 933.80 \approx 3.76.
\]
For Mayer's line, computing each residual $e_i = y_i - (10.82 + 2.708x_i)$ and summing the squares:
\begin{center}
\begin{tabular}{|c|c|c|c|}
\hline
$(x_i,y_i)$ & $\hat{y}_i$ (Mayer) & $e_i$ & $e_i^2$ \\
\hline
(2,15) & 16.24 & -1.24 & 1.54 \\
(4,22) & 21.65 & 0.35 & 0.12 \\
(5,25) & 24.36 & 0.64 & 0.41 \\
(7,30) & 29.78 & 0.22 & 0.05 \\
(8,33) & 32.48 & 0.52 & 0.27 \\
(10,38) & 37.90 & 0.10 & 0.01 \\
(11,40) & 40.61 & -0.61 & 0.37 \\
(12,44) & 43.32 & 0.68 & 0.46 \\
(14,48) & 48.73 & -0.73 & 0.53 \\
\hline
\end{tabular}
\end{center}
\[
\mathrm{RSS}_{M} \approx 1.54 + 0.12 + 0.41 + 0.05 + 0.27 + 0.01 + 0.37 + 0.46 + 0.53 \approx 3.76.
\]
The two fits are almost identical here, but $\mathrm{RSS}_{LS} \le \mathrm{RSS}_{M}$: the \textbf{least squares line provides the better fit}, as guaranteed by the minimisation property (Mayer's method only uses the two group means, whereas least squares uses every data point).
\end{enumerate}
\end{corrige}

\begin{corrige}[Exercise 8 — Rank Correlation: Singing Competition]
\begin{enumerate}
\item Differences $d = R_1 - R_2$:
\begin{center}
\begin{tabular}{|c|c|c|c|c|c|c|c|}
\hline
Singer & A & B & C & D & E & F & G \\
\hline
$d$ & $-1$ & $1$ & $-1$ & $1$ & $-1$ & $-1$ & $2$ \\
\hline
$d^2$ & $1$ & $1$ & $1$ & $1$ & $1$ & $1$ & $4$ \\
\hline
\end{tabular}
\end{center}
$\sum d^2 = 10$, $n = 7$:
\[
R_s = 1 - \frac{6 \times 10}{7(7^2 - 1)} = 1 - \frac{60}{336} = 1 - 0.1786 \approx \boxed{0.821}.
\]
\item \textbf{Interpretation:} $R_s \approx 0.82$ shows a \emph{strong positive agreement} between the two judges: singers ranked highly by Judge 1 tend to be ranked highly by Judge 2. The judges assess the performances in a very similar way.
\item \textbf{Adding Singer H} ranked 8th by both judges: the new difference is $d_H = 0$, so $\sum d^2$ stays equal to $10$, but $n$ increases from 7 to 8:
\[
R_s = 1 - \frac{6 \times 10}{8(64-1)} = 1 - \frac{60}{504} \approx 0.881.
\]
Since the denominator $n(n^2-1)$ grows while $\sum d^2$ is unchanged, the subtracted term shrinks and $R_s$ \emph{increases} towards 1: a new pair in perfect agreement strengthens the measured agreement.
\end{enumerate}
\end{corrige}

\begin{corrige}[Exercise 9 — GCE Style: Mixed Regression and Correlation Problem]
\textbf{Mark scheme and full solution.} First compute the summary quantities:
\[
\bar{x} = 6.5,\quad \bar{y} = 82,\quad S_{xx} = 485 - \frac{65^2}{10} = 62.5,
\]
\[
S_{yy} = 72500 - \frac{820^2}{10} = 5260, \qquad S_{xy} = 5020 - \frac{65 \times 820}{10} = -310.
\]
\begin{enumerate}
\item \textbf{[3 marks]} $\displaystyle r = \frac{-310}{\sqrt{62.5 \times 5260}} = \frac{-310}{573.4} \approx \boxed{-0.541}$.
\emph{(M1 correct formula; A1 correct $S_{xy}$, $S_{xx}$, $S_{yy}$; A1 answer to 3 s.f.)}
\item \textbf{[2 marks]} $r \approx -0.54$: a \emph{moderate negative linear correlation} — as the age of a taxi increases, its price tends to decrease, with considerable scatter. \emph{(B1 direction + strength; B1 contextual interpretation.)}
\item \textbf{[4 marks]}
\[
b = \frac{S_{xy}}{S_{xx}} = \frac{-310}{62.5} = -4.96, \qquad a = 82 - (-4.96)(6.5) = 114.24.
\]
\[
\boxed{y = 114.24 - 4.96\,x} \quad \emph{(M1 for $b$, A1; M1 for $a$ via $(\bar{x},\bar{y})$, A1.)}
\]
\item \textbf{[2 marks]} $x = 6$: $y = 114.24 - 4.96 \times 6 = 84.48$. Fair price $\approx \boxed{84\,500 \text{ FCFA}}$. \emph{(M1 substitution; A1 with units.)}
\item \textbf{[2 marks]} 12 years lies \emph{outside the range of the data} (maximum age 9 years): this is \textbf{extrapolation}. There is no evidence the linear trend continues (prices cannot fall below zero, for instance), so the estimate would be unreliable. \emph{(B1 extrapolation identified; B1 reason.)}
\item \textbf{[3 marks]}
\[
\mathrm{RSS} = S_{yy} - \frac{S_{xy}^2}{S_{xx}} = 5260 - \frac{(-310)^2}{62.5} = 5260 - 1537.6 = \boxed{3722.4}.
\]
\emph{(M1 correct formula; A1 substitution; A1 answer.)}
\item \textbf{[2 marks]} With $b = -4.96$ and the given $b' = -0.0045$: $b \times b' = 0.0223$, whereas $r^2 = 0.292$. These are \emph{not} equal, which reveals that the line given in the question is inconsistent with the data (a misprint). The correct line of $x$ on $y$ has slope $b' = \dfrac{S_{xy}}{S_{yy}} = \dfrac{-310}{5260} \approx -0.0589$, and then
\[
b \times b' = (-4.96)(-0.0589) \approx 0.292 = r^2. \checkmark
\]
\emph{Examiner expectation: M1 for computing both $bb'$ and $r^2$; A1 for the comparison and conclusion. A candidate who spots and corrects the inconsistency with justification earns full credit.}
\end{enumerate}
\end{corrige}

\begin{corrige}[Exercise 10 — GCE Style: Kendall's Rank Correlation and Interpretation]
\begin{enumerate}
\item \textbf{[5 marks]} A pair is \emph{concordant} if the two judges rank the two students in the same order, \emph{discordant} otherwise. Maths order: P $<$ Q $<$ R $<$ S $<$ T.
\begin{center}
\begin{tabular}{|l|c|c|c|c|c|c|c|c|c|c|}
\hline
Pair & PQ & PR & PS & PT & QR & QS & QT & RS & RT & ST \\
\hline
Type & D & C & C & C & C & C & C & D & D & C \\
\hline
\end{tabular}
\end{center}
(PQ: Maths P$<$Q but Physics P$>$Q — discordant; RS: Maths R$<$S but Physics R$>$S — discordant; RT: similarly discordant.) Hence $C = 7$ concordant, $D = 3$ discordant. \emph{(B1 per two correct pairs, rounded to 5 marks.)}
\item \textbf{[3 marks]}
\[
\tau = \frac{C - D}{\binom{5}{2}} = \frac{7 - 3}{10} = \boxed{0.4}.
\]
\emph{(M1 formula; A1 correct counts; A1 answer.)}
\item \textbf{[3 marks]} Differences $d$: P: $-1$, Q: $1$, R: $-2$, S: $1$, T: $1$; $\sum d^2 = 1+1+4+1+1 = 8$.
\[
R_s = 1 - \frac{6 \times 8}{5(25-1)} = 1 - \frac{48}{120} = \boxed{0.6}.
\]
\emph{(M1 differences; A1 $\sum d^2$; A1 answer.)}
\item \textbf{[3 marks]} Both coefficients are positive and indicate moderate agreement, but $R_s = 0.6 > \tau = 0.4$: the two coefficients measure agreement differently and need not be equal.
\begin{itemize}
\item \emph{Use of Kendall's $\tau$:} it has a direct probabilistic interpretation — $\tau$ measures (probability of concordance) $-$ (probability of discordance) — and it behaves better than $R_s$ for \emph{small samples} or when there are many ties.
\item \emph{Limitation:} it is more tedious to compute for large $n$ (all $\binom{n}{2}$ pairs must be examined), and its values are typically smaller in magnitude than $R_s$, so the two scales cannot be compared directly.
\end{itemize}
\emph{(B1 comparison; B1 valid use; B1 valid limitation.)}
\end{enumerate}
\end{corrige}

\begin{corrige}[Exercise 11 — GCE Style: Regression Lines and Mean Point]
\begin{enumerate}
\item \textbf{[3 marks]} Both regression lines pass through $(\bar{x}, \bar{y})$, so solve
\[
\bar{y} = 0.8\bar{x} + 3, \qquad \bar{x} = 0.45\bar{y} + 1.
\]
Substituting: $\bar{y} = 0.8(0.45\bar{y} + 1) + 3 = 0.36\bar{y} + 3.8$, so $0.64\bar{y} = 3.8$, giving $\bar{y} = 5.9375$ and $\bar{x} = 0.45 \times 5.9375 + 1 = 3.671875$.
\[
\boxed{(\bar{x}, \bar{y}) \approx (3.67,\ 5.94)} \quad \emph{(M1 using the mean-point property; A1 A1.)}
\]
\item \textbf{[3 marks]} $r^2 = b \times b' = 0.8 \times 0.45 = 0.36$. Since both slopes are positive,
\[
r = +\sqrt{0.36} = \boxed{0.6}. \quad \emph{(M1 $r^2 = bb'$; A1 value; B1 correct sign justified.)}
\]
\item \textbf{[3 marks]} $b = \dfrac{S_{xy}}{S_{xx}} = 0.8$ and $S_{xx} = 50$ give $S_{xy} = 0.8 \times 50 = 40$. Then
\[
b' = \frac{S_{xy}}{S_{yy}} = 0.45 \;\Longrightarrow\; S_{yy} = \frac{40}{0.45} \approx \boxed{88.9}.
\]
\emph{(M1 finding $S_{xy}$; M1 using $b'$; A1 answer.)}
\item \textbf{[3 marks]} The new point $(\bar{x}, \bar{y}+2)$ lies \emph{vertically above the mean point}. Its $x$-deviation from the mean is exactly zero, so it adds nothing to $S_{xy}$ and $S_{xx}$, but it adds a positive amount to $S_{yy}$. Consequently:
\begin{itemize}
\item the slope $b = S_{xy}/S_{xx}$ is \emph{unchanged};
\item $r = \dfrac{S_{xy}}{\sqrt{S_{xx}S_{yy}}}$ \emph{decreases in magnitude} (the denominator grows while the numerator stays fixed).
\end{itemize}
Intuitively, the new point does not follow the linear trend, so it weakens the measured linear correlation without tilting the line. \emph{(B1 slope unchanged with reason; B1 $r$ decreases; B1 quality of explanation.)}
\end{enumerate}
\end{corrige}

\newpage

\coursec{Summary sheet}

\begin{popfigure}
\begin{tikzpicture}[mindmap, grow cyclic, every node/.style={concept, text=white, font=\small},
  root concept/.append style={concept color=popink, font=\bfseries},
  level 1/.append style={level distance=3.6cm, sibling angle=60},
  level 1 concept/.append style={font=\small\bfseries}]
\node[root concept]{Regression \& Correlation}
  child[concept color=popblue]{ node {Scatter diagrams \& regression idea} }
  child[concept color=poporange]{ node {Mayer's method} }
  child[concept color=popgreen]{ node {Variance \& covariance} }
  child[concept color=poppurple]{ node {Least squares lines $y$ on $x$, $x$ on $y$} }
  child[concept color=poppink]{ node {RSS \& line of best fit} }
  child[concept color=popgold]{ node {PMCC $r$} }
  child[concept color=popdark]{ node {Rank correlation: Spearman, Kendall} };
\end{tikzpicture}
\legende{Mind map of the chapter: from scatter diagrams to rank correlation.}
\end{popfigure}

\begin{bilan}
\textbf{Key definitions}
\begin{itemize}
\item \cle{Bivariate data}: pairs $(x_i,y_i)$ observed on the same individuals.
\item \cle{Scatter diagram}: plot of the points $(x_i,y_i)$; it suggests the type of association (positive, negative, none, linear or not).
\item \cle{Regression line of $y$ on $x$}: line used to predict $y$ from $x$; \cle{regression line of $x$ on $y$}: line used to predict $x$ from $y$.
\item \cle{Residual}: $e_i=y_i-\hat y_i$, vertical gap between observation and fitted line.
\item \cle{Covariance}: measures how $x$ and $y$ vary together.
\item \cle{Correlation coefficient}: a number in $[-1,1]$ measuring the strength and direction of a linear (or rank) association.
\end{itemize}

\textbf{Formulas to know}
\begin{itemize}
\item Means: $\bar x=\dfrac{\sum x_i}{n}$, $\bar y=\dfrac{\sum y_i}{n}$.
\item Sample variance and covariance:
\[ s_x^2=\frac{1}{n}\sum(x_i-\bar x)^2=\frac{\sum x_i^2}{n}-\bar x^2,\qquad
s_{xy}=\frac{1}{n}\sum(x_i-\bar x)(y_i-\bar y)=\frac{\sum x_iy_i}{n}-\bar x\bar y. \]
\item Least squares line of $y$ on $x$: $y-\bar y=b(x-\bar x)$ with $b=\dfrac{s_{xy}}{s_x^2}$.
\item Least squares line of $x$ on $y$: $x-\bar x=b'(y-\bar y)$ with $b'=\dfrac{s_{xy}}{s_y^2}$.
\item Both lines pass through the \cle{mean point} $(\bar x,\bar y)$.
\item Residual sums of squares:
\[ \mathrm{RSS}_{y\text{ on }x}=\sum(y_i-\hat y_i)^2=\sum y_i^2-n\bar y^2-\frac{\bigl(\sum x_iy_i-n\bar x\bar y\bigr)^2}{\sum x_i^2-n\bar x^2}, \]
and symmetrically for $\mathrm{RSS}_{x\text{ on }y}$.
\item \cle{Line of best fit}: the line with the smaller RSS (equivalently the larger $|r|$ choice); the $y$-on-$x$ line is best for predicting $y$.
\item Product moment correlation coefficient:
\[ r=\frac{s_{xy}}{s_x s_y}=\frac{\sum x_iy_i-n\bar x\bar y}{\sqrt{\bigl(\sum x_i^2-n\bar x^2\bigr)\bigl(\sum y_i^2-n\bar y^2\bigr)}},\qquad -1\le r\le 1. \]
Also $bb'=r^2$.
\item Spearman's rank correlation (no tied ranks):
\[ r_s=1-\frac{6\sum d_i^2}{n(n^2-1)},\quad d_i=\text{rank}(x_i)-\text{rank}(y_i). \]
\item Kendall's rank correlation:
\[ \tau=\frac{2(P-Q)}{n(n-1)}=\frac{S}{\tfrac12 n(n-1)}, \]
where $P$ = number of concordant pairs, $Q$ = number of discordant pairs.
\end{itemize}

\textbf{Methods}
\begin{itemize}
\item \cle{Mayer's method}: order the data by $x$, split into two equal (or nearly equal) groups, compute the mean point of each group, and draw the line through these two points.
\item \cle{Least squares}: compute $\sum x$, $\sum y$, $\sum x^2$, $\sum y^2$, $\sum xy$ in a table; then $b$, then the line through $(\bar x,\bar y)$.
\item \cle{Ranking}: rank from highest (or lowest) consistently; give tied values the mean of the ranks they occupy.
\item \cle{Kendall}: order one variable's ranks naturally, then count, for each rank of the second variable, how many later ranks are larger ($P$) and smaller ($Q$).
\end{itemize}

\textbf{Common mistakes}
\begin{itemize}
\item Confusing $y$ on $x$ with $x$ on $y$: they are different lines (unless $|r|=1$).
\item Forgetting that both regression lines pass through $(\bar x,\bar y)$ — a useful check.
\item Using $r$ to claim causation: correlation is not causation.
\item Predicting outside the data range (extrapolation) without caution.
\item Sign errors in $\sum d_i^2$; always check $\sum d_i=0$.
\item Mishandling tied ranks in Spearman or Kendall.
\item Interpreting $r\approx 0$ as "no relation": there may be a strong non-linear relation.
\end{itemize}

\textbf{GCE tips}
\begin{itemize}
\item Show the table of sums explicitly; examiners award method marks for $\sum x^2$, $\sum xy$.
\item Quote the formula before substituting; keep at least 4 significant figures until the final answer.
\item Always state the interpretation of $r$ or $r_s$ in words ("strong positive linear correlation").
\item For prediction questions, choose the correct regression line: to estimate $y$ given $x$, use $y$ on $x$.
\item Kendall's $\tau$ is more robust with small samples and ties, but Spearman's $r_s$ is quicker to compute; know when each is appropriate.
\end{itemize}
\end{bilan}

\findecours

\end{document}
